% Production et traduction de textes sur la situation de la femme et les autres thèmes du chapitre 1 — Allemand Tle A — Cours signé OnBuch+
% Programme officiel MINESEC. Compiler avec : tectonic AL05-production-et-traduction-de-textes-sur-la-situation-de-la-fe.tex
\newcommand{\DOCMATIERE}{Allemand}
\newcommand{\DOCNIVEAU}{Tle A}
\newcommand{\DOCMODULE}{Chapitre 1 : Vie familiale et sociale}
\newcommand{\DOCLECON}{AL05}
\newcommand{\DOCTITRE}{Production et traduction de textes sur la situation de la femme et les autres thèmes du chapitre 1}
% OnBuch+ — préambule « cours signés » ENRICHI (compilé avec tectonic / XeLaTeX)
% Système visuel « Pop » d'OnBuch+ : fond crème (jamais blanc), bordures encre
% épaisses, ombres « dures » décalées, titres Archivo Black en capitales.
% Version MATHÉMATIQUES : graphes (pgfplots/tikz), boîtes pédagogiques
% (définition, propriété, méthode, attention, le savais-tu, exercice résolu,
% exercice, TP, bilan), page de garde et signature OnBuch+ ancrée en bas.
%
% Macros attendues dans main.tex AVANT \input{preamble} :
%   \DOCMATIERE  \DOCNIVEAU  \DOCMODULE  \DOCLECON (numéro)  \DOCTITRE
\documentclass[11pt]{article}

\usepackage{fontspec}
\usepackage[ngerman,french]{babel} % français = langue principale ; l'allemand via \all{..} / textebox
\usepackage[a4paper,margin=2cm,top=3.4cm,bottom=2.8cm,headheight=1.4cm,footskip=1.3cm]{geometry}
\usepackage{amsmath,amssymb,amsfonts}
\usepackage{mathtools} % \xmapsto, \coloneqq... souvent utilisés par les modèles
\usepackage{cancel} % simplifications barrées (bilans de forces)
% \abs{z} (module, valeur absolue) et \norm{u} (norme), utilisés par les modèles
\providecommand{\abs}{}\let\abs\relax\DeclarePairedDelimiter{\abs}{\lvert}{\rvert}
\providecommand{\norm}{}\let\norm\relax\DeclarePairedDelimiter{\norm}{\lVert}{\rVert}
\usepackage[table]{xcolor} % [table] : fournit aussi \rowcolors (alternance de lignes)
\usepackage{pagecolor}
\usepackage{tikz}
\usetikzlibrary{arrows.meta,positioning,calc,shapes.geometric,shapes.misc,shapes.arrows,decorations.pathmorphing,decorations.pathreplacing,decorations.markings,patterns,shadows,mindmap,trees,fit,backgrounds,matrix,intersections,angles,quotes}
\usepackage{pgfplots}
\usepackage[version=4]{mhchem} % équations nucléaires \ce{^{235}_{92}U}
\usepackage[european,siunitx]{circuitikz} % schémas électriques
\pgfplotsset{compat=1.18}
\usepgfplotslibrary{fillbetween,groupplots} % aires entre courbes (calcul intégral)
\usepackage{enumitem}
\usepackage{fancyhdr}
% [most] charge toutes les bibliothèques tcolorbox (listings, minted, raster,
% documentation...) et gonfle inutilement la mémoire TeX sur un document long
% (~300 boîtes) : on ne charge que ce qui est réellement utilisé.
\usepackage[skins,breakable]{tcolorbox}
\usepackage{graphicx}
\usepackage{colortbl}
\usepackage{booktabs}
\usepackage{tabularx}
\usepackage{diagbox} % case d'en-tête diagonale des tableaux à double entrée
% Colonnes extensibles centrée / gauche / droite (C, L, R), souvent utilisées par les modèles
\newcolumntype{C}{>{\centering\arraybackslash}X}
\newcolumntype{L}{>{\raggedright\arraybackslash}X}
\newcolumntype{R}{>{\raggedleft\arraybackslash}X}
\usepackage{array}
\usepackage{multicol}
\usepackage{multirow}
\usepackage{siunitx}
% parse-numbers=false : accepte aussi \SI{2,5 \times 10^{-3}}{...}
\sisetup{locale=FR,output-decimal-marker={,},group-minimum-digits=4,parse-numbers=false}
\DeclareSIUnit\ohm{\ensuremath{\symup{\Omega}}} % U+2126 absent de la police
\DeclareSIUnit\division{div} % réglages d'oscilloscope (ms/div, V/div)
\DeclareSIUnit\tr{tr} % tours (vitesse de rotation, ex. \SI{1500}{\tr\per\minute})
\DeclareSIUnit\molar{M} % concentration molaire (ex. \SI{3}{\molar}, \SI{1}{\micro\molar})
\DeclareSIUnit\an{an} % année (le modèle confond parfois avec \year, un compteur TeX) — ex. \SI{5}{\kilo\meter\per\an}
\DeclareSIUnit\FCFA{FCFA} % franc CFA (ex. \SI{500}{\FCFA\per\kilo\gram})
\DeclareSIUnit\degreeBrix{\textdegree Brix} % teneur en sucre d'un sirop/jus (réfractométrie)
\DeclareSIUnit\month{mois} % le modèle utilise parfois \month, un compteur TeX, comme unité de durée
\usepackage{qrcode}
% \degree : commande fréquemment utilisée par les modèles pour un angle ou une
% température en dehors de \SI{}{} (non fournie par les paquets chargés).
\providecommand{\degree}{^{\circ}}
% \qed (fin de démonstration), utilisé par les modèles sans amsthm
\providecommand{\qed}{\hfill$\square$}
% \barre : double barre des valeurs interdites dans un tableau de variations (tkz-tab)
\providecommand{\barre}{\ensuremath{\|}}
% environnement proof (amsthm non chargé) utilisé par les modèles
\ifdefined\proof\else\newenvironment{proof}[1][Démonstration]{\par\noindent\textit{#1.}\ }{\qed\par}\fi
\usepackage{lastpage}
\usepackage{hyperref}
\hypersetup{hidelinks}

% ============================================================
% Palette « Pop » OnBuch+ — mêmes hex que la classe P (plus_ui.dart)
% ============================================================
\definecolor{poporange}{HTML}{F59321}
\definecolor{popdark}{HTML}{E07A0C}
\definecolor{popcream}{HTML}{FFF6E8}
\definecolor{popink}{HTML}{1C1714}
\definecolor{poppurple}{HTML}{7A5AE0}
\definecolor{popgreen}{HTML}{2E9E6A}
\definecolor{poppink}{HTML}{E0457B}
\definecolor{popblue}{HTML}{2F7DE1}
\definecolor{popgold}{HTML}{C98A12}
\definecolor{popmuted}{HTML}{8A7F76}
\definecolor{popline}{HTML}{EADFD2}
\definecolor{popgray}{HTML}{8A7F76}
\colorlet{popgrey}{popgray}
% Alias tolérants (le modèle nomme parfois les couleurs différemment)
\colorlet{popwhite}{white}
\colorlet{popblack}{popink}
\colorlet{popred}{poppink}
\colorlet{popyellow}{popgold}
% Teintes claires (fonds de tableaux/encadrés) — le modèle nomme parfois
% les couleurs avec un suffixe L (ex. popblueL) sans les avoir définies.
\colorlet{poporangeL}{poporange!20}
\colorlet{popdarkL}{popdark!20}
\colorlet{poppurpleL}{poppurple!20}
\colorlet{popgreenL}{popgreen!20}
\colorlet{poppinkL}{poppink!20}
\colorlet{popblueL}{popblue!20}
\colorlet{popgoldL}{popgold!20}
\colorlet{popredL}{popred!20}
\colorlet{popgrayL}{popgray!20}
% Teintes claires pour fonds de tableaux / zones
\colorlet{popblueL}{popblue!10!white}
\colorlet{poporangeL}{poporange!14!white}
\colorlet{popgreenL}{popgreen!12!white}
\colorlet{poppurpleL}{poppurple!12!white}
\colorlet{poppinkL}{poppink!12!white}
\colorlet{popgoldL}{popgold!14!white}

\pagecolor{popcream}

% ============================================================
% Typographie
% ============================================================
\defaultfontfeatures{Ligatures=TeX}
\setmainfont{PlusJakartaSans-Medium.ttf}[Path=../../../fonts/,BoldFont=PlusJakartaSans-Bold.ttf]
% Maths en sans-serif (Fira Math) avec chiffres et lettres droites de Plus
% Jakarta, pour que les formules se fondent dans le texte.
\usepackage[math-style=ISO,bold-style=ISO]{unicode-math}
\setmathfont{FiraMath-Regular.otf}[Path=../../../fonts/]
\setmathfont{PlusJakartaSans-Medium.ttf}[Path=../../../fonts/,range={up/num,up/{Latin,latin},\mathup}]
\usepackage{icomma} % virgule décimale sans espace en mode math
% Caractères mathématiques Unicode absents des polices de texte (Plus Jakarta,
% Archivo Black) : rendus par leur équivalent mathématique, en texte comme en
% formule — sinon ils s'affichent en carré vide (ex. « Arithmétique dans ℤ »).
\usepackage{newunicodechar}
\newunicodechar{ℝ}{\ensuremath{\mathbb{R}}}
\newunicodechar{ℤ}{\ensuremath{\mathbb{Z}}}
\newunicodechar{ℂ}{\ensuremath{\mathbb{C}}}
\newunicodechar{ℕ}{\ensuremath{\mathbb{N}}}
\newunicodechar{ℚ}{\ensuremath{\mathbb{Q}}}
\newunicodechar{𝔻}{\ensuremath{\mathbb{D}}}
\newunicodechar{α}{\ensuremath{\alpha}}
\newunicodechar{β}{\ensuremath{\beta}}
\newunicodechar{γ}{\ensuremath{\gamma}}
\newunicodechar{δ}{\ensuremath{\delta}}
\newunicodechar{ε}{\ensuremath{\varepsilon}}
\newunicodechar{θ}{\ensuremath{\theta}}
\newunicodechar{λ}{\ensuremath{\lambda}}
\newunicodechar{μ}{\ensuremath{\mu}}
\newunicodechar{σ}{\ensuremath{\sigma}}
\newunicodechar{τ}{\ensuremath{\tau}}
\newunicodechar{φ}{\ensuremath{\varphi}}
\newunicodechar{ψ}{\ensuremath{\psi}}
\newunicodechar{ω}{\ensuremath{\omega}}
\newunicodechar{Σ}{\ensuremath{\Sigma}}
% majuscules produites par \MakeUppercase dans les titres (α-aminés -> Α)
\newunicodechar{Α}{\ensuremath{\alpha}}
\newunicodechar{Β}{\ensuremath{\beta}}
\newunicodechar{Γ}{\ensuremath{\Gamma}}
\newunicodechar{Δ}{\ensuremath{\Delta}}
\newunicodechar{Ε}{\ensuremath{\varepsilon}}
\newunicodechar{Θ}{\ensuremath{\Theta}}
\newunicodechar{Λ}{\ensuremath{\Lambda}}
\newunicodechar{Μ}{\ensuremath{\mu}}
\newunicodechar{Τ}{\ensuremath{\tau}}
\newunicodechar{Φ}{\ensuremath{\Phi}}
\newunicodechar{Ψ}{\ensuremath{\Psi}}
\newunicodechar{Ω}{\ensuremath{\Omega}}
\newunicodechar{↦}{\ensuremath{\mapsto}}
\newunicodechar{∈}{\ensuremath{\in}}
\newunicodechar{∉}{\ensuremath{\notin}}
\newunicodechar{∧}{\ensuremath{\wedge}}
\newunicodechar{⇔}{\ensuremath{\Leftrightarrow}}
\newunicodechar{⟺}{\ensuremath{\Longleftrightarrow}}
\newunicodechar{⇒}{\ensuremath{\Rightarrow}}
\newunicodechar{⟹}{\ensuremath{\Longrightarrow}}
\newunicodechar{∘}{\ensuremath{\circ}}
\newunicodechar{∪}{\ensuremath{\cup}}
\newunicodechar{∩}{\ensuremath{\cap}}
\newunicodechar{≡}{\ensuremath{\equiv}}
\newunicodechar{⊂}{\ensuremath{\subset}}
\newunicodechar{⊥}{\ensuremath{\perp}}
\newunicodechar{∃}{\ensuremath{\exists}}
\newunicodechar{∀}{\ensuremath{\forall}}
\newunicodechar{∛}{\ensuremath{\sqrt[3]{\,}}}
\newunicodechar{∜}{\ensuremath{\sqrt[4]{\,}}}
\newunicodechar{⃗}{\ensuremath{{}^{\to}}}
\newunicodechar{ⁿ}{\ifmmode^{n}\else\textsuperscript{\ensuremath{n}}\fi}
\newunicodechar{ˣ}{\ifmmode^{x}\else\textsuperscript{\ensuremath{x}}\fi}
\newunicodechar{ᵇ}{\ifmmode^{b}\else\textsuperscript{\ensuremath{b}}\fi}
\newunicodechar{ʳ}{\ifmmode^{r}\else\textsuperscript{\ensuremath{r}}\fi}
\newunicodechar{ᵘ}{\ifmmode^{u}\else\textsuperscript{\ensuremath{u}}\fi}
\newunicodechar{ʸ}{\ifmmode^{y}\else\textsuperscript{\ensuremath{y}}\fi}
\newunicodechar{ᵃ}{\ifmmode^{a}\else\textsuperscript{\ensuremath{a}}\fi}
\newunicodechar{ᵉ}{\ifmmode^{e}\else\textsuperscript{\ensuremath{e}}\fi}
\newunicodechar{ᵏ}{\ifmmode^{k}\else\textsuperscript{\ensuremath{k}}\fi}
\newunicodechar{ᵖ}{\ifmmode^{p}\else\textsuperscript{\ensuremath{p}}\fi}
\newunicodechar{ᴺ}{\ifmmode^{N}\else\textsuperscript{\ensuremath{N}}\fi}
\newunicodechar{ᶜ}{\ifmmode^{c}\else\textsuperscript{\ensuremath{c}}\fi}
\newunicodechar{ʰ}{\ifmmode^{h}\else\textsuperscript{\ensuremath{h}}\fi}
\newunicodechar{ᵠ}{\ifmmode^{q}\else\textsuperscript{\ensuremath{q}}\fi}
\newunicodechar{ₙ}{\ifmmode_{n}\else\textsubscript{\ensuremath{n}}\fi}
\newunicodechar{ₐ}{\ifmmode_{a}\else\textsubscript{\ensuremath{a}}\fi}
\newunicodechar{ᵢ}{\ifmmode_{i}\else\textsubscript{\ensuremath{i}}\fi}
\newunicodechar{ᵣ}{\ifmmode_{r}\else\textsubscript{\ensuremath{r}}\fi}
\newunicodechar{ₖ}{\ifmmode_{k}\else\textsubscript{\ensuremath{k}}\fi}
\newunicodechar{ᵦ}{\ifmmode_{\beta}\else\textsubscript{\ensuremath{\beta}}\fi}
\newunicodechar{ₓ}{\ifmmode_{x}\else\textsubscript{\ensuremath{x}}\fi}
\newfontfamily\popdisplay{ArchivoBlack-Regular.ttf}[Path=../../../fonts/]
\newfontfamily\poplogo{SpaceGrotesk-Bold.ttf}[Path=../../../fonts/]
\newfontfamily\popsemi{PlusJakartaSans-SemiBold.ttf}[Path=../../../fonts/]
\newfontfamily\popxbold{PlusJakartaSans-ExtraBold.ttf}[Path=../../../fonts/]
\newfontfamily\popxboldls{PlusJakartaSans-ExtraBold.ttf}[Path=../../../fonts/,LetterSpace=3.0]

\setlist[enumerate]{leftmargin=1.7em,itemsep=5pt,topsep=4pt}
\setlist[itemize]{leftmargin=1.7em,itemsep=4pt,topsep=4pt,label={\color{poporange}\textbullet}}
\setlength{\parindent}{0pt}
\setlength{\parskip}{5pt}
\renewcommand{\arraystretch}{1.25}

% pgfplots : style Pop par défaut
\pgfplotsset{
  every axis/.append style={axis line style={popink,line width=1pt},tick style={popink},
    grid style={popline,line width=0.5pt},label style={font=\small},tick label style={font=\footnotesize}},
  popcurve/.style={poporange,line width=1.8pt,no markers,smooth},
}
% Même style disponible dans un \draw TikZ simple (hors pgfplots)
\tikzset{popcurve/.style={poporange,line width=1.8pt}}
\tikzset{popgrid/.style={popline,very thin}}
\colorlet{popcurve}{poporange} % le nom du style est parfois utilisé comme couleur

% ============================================================
% Pill / squiggle
% ============================================================
\newcommand{\poppill}[3]{%
  \tikz[baseline=-0.55ex]\node[draw=popink,line width=1.2pt,rounded corners=2.6pt,%
    fill=#2,inner xsep=6.5pt,inner ysep=3pt]{%
    \popxboldls\fontsize{7}{7}\selectfont\color{#3}\MakeUppercase{#1}};%
}
\newcommand{\popsquiggle}[1]{%
  \tikz[baseline=-0.4ex]{
    \draw[#1,line width=2.6pt,line cap=round]
      plot [smooth] coordinates {(0,0.09) (0.34,0.01) (0.68,0.21) (1.02,0.01) (1.36,0.10)};
  }%
}
% Mot-clé surligné (marqueur orange sous le texte)
\newcommand{\cle}[1]{{\popxbold\color{popdark}#1}}

% ============================================================
% En-tête / pied
% ============================================================
\pagestyle{fancy}
\fancyhf{}
\renewcommand{\headrulewidth}{0pt}
\renewcommand{\footrulewidth}{0pt}
\lhead{\raisebox{-0.15ex}{\poplogo\fontsize{13}{13}\selectfont\color{popink}OnBuch\color{poporange}+}}
\rhead{\poppill{\DOCMATIERE\ \textbullet\ \DOCNIVEAU\ \textbullet\ Leçon \DOCLECON}{white}{popink}}
\lfoot{\popxbold\fontsize{7.6}{7.6}\selectfont\color{popmuted}\MakeUppercase{Cours signé OnBuch+}\ \textbullet\ {\popsemi\DOCTITRE}}
\rfoot{\tikz[baseline=-0.6ex]\node[rounded corners=8.5pt,fill=popink,minimum height=17pt,minimum width=17pt,inner xsep=5pt,inner sep=0]{\popxbold\fontsize{8}{8}\selectfont\color{popcream}\thepage\,/\,\pageref*{LastPage}};}

% ============================================================
% Titres
% ============================================================
\newcommand{\chaptitre}[2]{%
  \begin{tcolorbox}[enhanced,colback=poporange,colframe=popink,boxrule=2.6pt,arc=11pt,
      drop shadow=popink,left=15pt,right=15pt,top=12pt,bottom=11pt,before skip=3pt,after skip=18pt]
    \poppill{#2}{popink}{popcream}\par\vspace{9pt}
    {\raggedright\hyphenpenalty=10000\exhyphenpenalty=10000\popdisplay\fontsize{22}{24}\selectfont\color{popcream}\MakeUppercase{#1}\par}
  \end{tcolorbox}
}
% Section numérotée : pastille orange avec le numéro + titre Archivo
\newcounter{popsec}
\newcommand{\coursec}[1]{%
  \stepcounter{popsec}\setcounter{popsub}{0}%
  \par\vspace{16pt}\noindent
  \begin{minipage}{\linewidth}
  \tikz[baseline=-0.7ex]\node[draw=popink,line width=1.6pt,fill=poporange,rounded corners=4pt,minimum size=22pt,inner sep=0]{\popdisplay\fontsize{12}{12}\selectfont\color{popcream}\thepopsec};%
  \hspace{8pt}{\popdisplay\fontsize{15}{17}\selectfont\color{popink}\MakeUppercase{#1}}\par
  \vspace{4pt}\noindent{\color{poporange}\rule{36pt}{3.6pt}}
  \end{minipage}\par\nopagebreak\vspace{8pt}
}
\newcounter{popsub}
\newcommand{\courssub}[1]{%
  \stepcounter{popsub}%
  \par\vspace{9pt}\noindent{\popxbold\fontsize{11.5}{13}\selectfont\color{popink}{\color{poporange}\thepopsec.\thepopsub}\hspace{6pt}#1}\par\nopagebreak\vspace{4pt}
}

% ============================================================
% Boîtes pédagogiques — même recette : fond blanc, bordure épaisse colorée,
% ombre dure encre, badge plein. Titre optionnel [#1].
% ============================================================
\tcbset{popbase/.style={breakable,enhanced,colback=white,boxrule=2.3pt,arc=9pt,
  shadow={3pt}{-3pt}{0pt}{popink},left=14pt,right=14pt,top=11pt,bottom=11pt,before skip=11pt,after skip=15pt,
  fontupper=\popsemi\fontsize{10.6}{14.5}\selectfont\color{popink},
  fontlower=\popsemi\fontsize{10.6}{14.5}\selectfont\color{popink}}}
% Coche dessinée (le glyphe ✓ n'existe pas dans les polices du document)
\newcommand{\popcheck}[1]{\tikz[baseline=-0.5ex]\draw[#1,line width=1.6pt,line cap=round,line join=round](-0.13,0.0)--(-0.04,-0.09)--(0.13,0.1);}
\providecommand{\coche}{\popcheck{popgreen}}
\providecommand{\resultat}[1]{\par\smallskip\noindent\fcolorbox{popgreen}{popgreenL}{\parbox{\dimexpr\linewidth-2\fboxsep-2\fboxrule\relax}{\textbf{Résultat.} #1}}\par\smallskip}
\newcommand{\popboxhead}[3]{\poppill{#1}{#2}{white}\ifx\relax#3\relax\else\hspace{6pt}{\popxbold\fontsize{10.6}{12}\selectfont\color{popink}#3}\fi\par\vspace{7pt}}

\newtcolorbox{aretenir}[1][]{popbase,colframe=popgreen,before upper={\popboxhead{À retenir}{popgreen}{#1}}}
% Texte en allemand (document de lecture, dialogue, extrait) : cadre
% d'étude. Un titre optionnel (source, auteur) s'affiche dans l'en-tête.
\newtcolorbox{textebox}[1][]{popbase,colframe=popblue,before upper={\popboxhead{Text}{popblue}{#1}\begin{otherlanguage}{ngerman}},after upper={\end{otherlanguage}}}
% Marques ✓ / ✗ (forme correcte / fautive) souvent utilisées par les modèles
\providecommand{\cross}{\textcolor{poppink}{\ensuremath{\times}}}
\providecommand{\xmark}{\cross}\providecommand{\crossmark}{\cross}\providecommand{\cmark}{\ensuremath{\checkmark}}
% Ligne de réponse / justification à compléter (inventée par les modèles)
\providecommand{\justification}[1]{\par\noindent\textit{Justification~:} \dotfill\par}
% \textipa (package tipa non chargé) : la police ne couvre pas l'API -> simple passage
\providecommand{\textipa}[1]{#1}
% Soulignement (ulem non chargé) : \uline, \ul
\providecommand{\uline}[1]{\underline{#1}}\providecommand{\ul}[1]{\underline{#1}}
% variantes ASCII d'ordinaux écrites en macro
\providecommand{\iere}{\textsuperscript{re}}\providecommand{\ieme}{\textsuperscript{e}}\providecommand{\ere}{\textsuperscript{re}}\providecommand{\eme}{\textsuperscript{e}}\providecommand{\er}{\textsuperscript{er}}\providecommand{\ie}{\textsuperscript{e}}
% Ordinaux français écrits en macro par les modèles : 1\ière, 3\ième
\newcommand{\ière}{\textsuperscript{re}}\newcommand{\ième}{\textsuperscript{e}}
% \exemple (inventée par les modèles) : étiquette « Exemple : »
\providecommand{\exemple}{\textbf{Exemple~:}\ }
% \square utilisable aussi hors mode math (cases à cocher)
\AtBeginDocument{\let\squareMath\square\renewcommand{\square}{\ensuremath{\squareMath}}}
% Mot ou phrase en allemand à l'intérieur d'un paragraphe français
\newcommand{\all}[1]{\ifmmode\text{\textit{#1}}\else\begingroup\selectlanguage{ngerman}\itshape #1\endgroup\fi}
\let\esp\all % alias toléré
\newtcolorbox{exemplebox}[1][]{popbase,colframe=poporange,before upper={\popboxhead{Exemple}{poporange}{#1}}}
\newtcolorbox{definition}[1][]{popbase,colframe=popblue,before upper={\popboxhead{Définition}{popblue}{#1}}}
\newtcolorbox{propriete}[1][]{popbase,colframe=poppurple,before upper={\popboxhead{Propriété}{poppurple}{#1}}}
\newtcolorbox{methode}[1][]{popbase,colframe=popgold,before upper={\popboxhead{Méthode}{popgold}{#1}}}
\newtcolorbox{attention}[1][]{popbase,colframe=poppink,before upper={\popboxhead{Attention}{poppink}{#1}}}
\newtcolorbox{experience}[1][]{popbase,colframe=popblue,colback=popblueL,before upper={\popboxhead{Expérience}{popblue}{#1}}}
% « Le savais-tu ? » : carte violette pleine (contexte, culture, Cameroun)
\newtcolorbox{savaistu}[1][]{popbase,colback=poppurple,colframe=popink,
  fontupper=\popsemi\fontsize{10.4}{14.2}\selectfont\color{white},
  before upper={\poppill{Le savais-tu\,?}{white}{poppurple}\ifx\relax#1\relax\else\hspace{6pt}{\popxbold\color{white}#1}\fi\par\vspace{7pt}}}
% Exercice résolu : énoncé en haut, \tcblower puis la solution sur fond crème
\newcounter{popexo}\newcounter{popexr}
\newtcolorbox{exoresolu}[1][]{popbase,colframe=poporange,colbacklower=popcream,
  segmentation style={popink,line width=1.2pt,dashed},
  before upper={\stepcounter{popexr}\popboxhead{Exercice résolu \thepopexr}{poporange}{#1}},
  before lower={\poppill{Solution}{popink}{popcream}\par\vspace{6pt}}}
% Exercice (non résolu en place) — la correction est dans la partie Corrigés
\newtcolorbox{exercice}[1][]{popbase,colframe=popink,
  before upper={\stepcounter{popexo}\popboxhead{Exercice \thepopexo}{popink}{#1}}}
% Corrigé
\newtcolorbox{corrige}[1][]{popbase,colframe=popgreen,colback=popgreenL,
  before upper={\popboxhead{Corrigé}{popgreen}{#1}}}
% Objectifs / prérequis / bilan
\newtcolorbox{objectifs}{popbase,colframe=popink,colback=poporangeL,
  before upper={\popboxhead{Objectifs de la leçon}{popink}{}}}
\newtcolorbox{prerequis}{popbase,colframe=popink,colback=popblueL,
  before upper={\popboxhead{Prérequis}{popblue}{}}}
\newtcolorbox{bilan}{popbase,colframe=popink,colback=white,boxrule=2.8pt,
  before upper={\popboxhead{Fiche bilan}{poporange}{L'essentiel à connaître pour l'examen}}}
% Figure encadrée avec légende
\newtcolorbox{popfigure}[1][]{enhanced,breakable=false,colback=white,colframe=popline,boxrule=1.4pt,arc=8pt,
  left=10pt,right=10pt,top=10pt,bottom=8pt,before skip=10pt,after skip=12pt,halign=center,
  lower separated=false,halign lower=center,
  fontlower=\popsemi\fontsize{9}{11.5}\selectfont\color{popmuted},
  #1}
\newcommand{\legende}[1]{\par\vspace{6pt}{\popsemi\fontsize{9}{11.5}\selectfont\color{popmuted}#1}}

% ============================================================
% Signature OnBuch+
% ============================================================
\newcommand{\poplogoinline}[1]{{\poplogo\fontsize{#1}{#1}\selectfont\color{popink}OnBuch\color{poporange}+}}
\newcommand{\signatureonbuch}{%
  \begin{tcolorbox}[enhanced,colback=popink,colframe=popink,boxrule=2.2pt,arc=10pt,
      drop shadow=poporange,left=14pt,right=14pt,top=10pt,bottom=10pt,before skip=0pt,after skip=0pt]
    \tikz[baseline=-0.7ex]\node[circle,fill=popgreen,draw=popcream,line width=1.3pt,minimum size=22pt,inner sep=0]{\popcheck{white}};%
    \hspace{9pt}\begin{minipage}[c]{0.62\linewidth}
      {\popxbold\fontsize{12}{13}\selectfont\color{popcream}Signé\ \textbullet\ {\poplogo OnBuch\color{poporange}+}}\par\vspace{2pt}
      {\raggedright\popsemi\fontsize{8}{10.5}\selectfont\color{popline}\MakeUppercase{Équipe pédagogique OnBuch+}\\\MakeUppercase{Conforme au programme officiel MINESEC}\par}
    \end{minipage}\hfill
    {\poplogo\fontsize{22}{22}\selectfont\color{popcream}OnBuch\color{poporange}+}
  \end{tcolorbox}%
}

% ============================================================
% Page de garde
% #1 titre · #2 matière · #3 niveau · #4 module · #5 n° leçon · #6 durée ·
% #7 description / objectifs (texte ou itemize)
% ============================================================
\newcommand{\pagegarde}[7]{%
  \thispagestyle{empty}
  \newgeometry{a4paper,margin=1.7cm,top=1.6cm,bottom=1.4cm}%
  \begin{tikzpicture}[remember picture,overlay]
    \fill[poporange!18] ([xshift=-3.2cm,yshift=-3.6cm]current page.north east) circle (4.6cm);
    \fill[poppurple!14] ([xshift=2.4cm,yshift=7.6cm]current page.south west) circle (3.8cm);
  \end{tikzpicture}%
  {\poplogo\fontsize{30}{30}\selectfont\color{popink}OnBuch\color{poporange}+}\hfill
  \raisebox{0.6ex}{\poppill{Cours signé \textbullet\ Premium}{popink}{popcream}}\par
  \vspace{1pt}\popsquiggle{poppurple}\par
  \vspace{8pt}
  \begin{tcolorbox}[enhanced,colback=poporange,colframe=popink,boxrule=2.8pt,arc=13pt,
      drop shadow=popink,left=18pt,right=18pt,top=12pt,bottom=13pt,after skip=10pt]
    \poppill{#2 \textbullet\ #3}{popink}{popcream}\hspace{5pt}\poppill{#4}{white}{popink}\par\vspace{8pt}
    {\popdisplay\fontsize{14}{14}\selectfont\color{popink}LEÇON #5}\par\vspace{4pt}
    {\raggedright\hyphenpenalty=10000\exhyphenpenalty=10000\popdisplay\fontsize{25}{27}\selectfont\color{popcream}\MakeUppercase{#1}\par}
    \vspace{8pt}
    {\popxbold\fontsize{9.5}{9.5}\selectfont\color{popink}\MakeUppercase{Durée conseillée : #6}}
  \end{tcolorbox}
  \begin{tcolorbox}[enhanced,colback=white,colframe=popink,boxrule=2.2pt,arc=10pt,
      drop shadow=popink,left=16pt,right=16pt,top=10pt,bottom=10pt,after skip=8pt]
    \poppill{Dans cette leçon}{poppurple}{white}\par\vspace{5pt}
    \popsemi\fontsize{9.3}{12.6}\selectfont\color{popink}#7
  \end{tcolorbox}
  \noindent\begin{minipage}[t]{0.487\linewidth}
    \begin{tcolorbox}[enhanced,colback=popgreen,colframe=popink,boxrule=2.2pt,arc=10pt,
        drop shadow=popink,left=11pt,right=11pt,top=10pt,bottom=10pt,height=3.1cm,valign=center]
      \tcbox[colback=white,colframe=popink,boxrule=1.3pt,arc=5pt,boxsep=0pt,nobeforeafter,
        left=3pt,right=3pt,top=3pt,bottom=3pt,valign=center]{\qrcode[height=1.9cm]{https://play.google.com/store/apps/details?id=com.luvvix.onbuch&hl=fr}}%
      \hspace{8pt}\begin{minipage}[c]{0.5\linewidth}
        {\popdisplay\fontsize{11}{12.5}\selectfont\color{white}\MakeUppercase{Télécharge OnBuch}}\par\vspace{4pt}
        {\raggedright\popsemi\fontsize{8.4}{11}\selectfont\color{white}Gratuit \textbullet\ Google Play\\Tuteur IA, annales, concours, résultats\par}
      \end{minipage}
    \end{tcolorbox}
  \end{minipage}\hfill
  \begin{minipage}[t]{0.487\linewidth}
    \begin{tcolorbox}[enhanced,colback=poppurple,colframe=popink,boxrule=2.2pt,arc=10pt,
        drop shadow=popink,left=11pt,right=11pt,top=10pt,bottom=10pt,height=3.1cm,valign=center]
      \tikz[baseline=-0.7ex]\node[circle,fill=white,draw=popink,line width=1.3pt,minimum size=1.6cm,inner sep=0]{\popdisplay\fontsize{24}{24}\selectfont\color{poppurple}+};%
      \hspace{8pt}\begin{minipage}[c]{0.6\linewidth}
        {\popdisplay\fontsize{11}{12.5}\selectfont\color{white}\MakeUppercase{Passe à OnBuch+}}\par\vspace{4pt}
        {\raggedright\popsemi\fontsize{8.4}{11}\selectfont\color{white}Tous les cours signés, Tuteur IA illimité, fiches PDF — onglet OnBuch+ de l'app\par}
      \end{minipage}
    \end{tcolorbox}
  \end{minipage}
  \vspace{8pt}
  \popcontenu
  \vfill
  \signatureonbuch
  \restoregeometry
  \newpage
}

% Rangée « Ce cours contient » (page de garde)
\newcommand{\poptile}[3]{% #1 couleur #2 gros texte #3 légende
  \begin{tcolorbox}[enhanced,colback=#1,colframe=popink,boxrule=2pt,arc=9pt,shadow={2.5pt}{-2.5pt}{0pt}{popink},
      left=6pt,right=6pt,top=9pt,bottom=9pt,halign=center,valign=center,height=1.75cm,width=\linewidth,nobeforeafter]
    {\popdisplay\fontsize{12.5}{13}\selectfont\color{white}\MakeUppercase{#2}}\par\vspace{3pt}
    {\centering\popsemi\fontsize{7.8}{9.5}\selectfont\color{white}#3\par}
  \end{tcolorbox}}
\newcommand{\popcontenu}{%
  \par\noindent{\popxboldls\fontsize{7.5}{7.5}\selectfont\color{popmuted}CE COURS SIGNÉ CONTIENT}\par\vspace{7pt}
  \noindent\begin{minipage}[t]{0.235\linewidth}\poptile{popblue}{Cours}{complet, illustré, pas à pas}\end{minipage}\hfill
  \begin{minipage}[t]{0.235\linewidth}\poptile{poporange}{Méthodes}{et exercices résolus}\end{minipage}\hfill
  \begin{minipage}[t]{0.235\linewidth}\poptile{poppink}{Exercices}{type examen + corrigés}\end{minipage}\hfill
  \begin{minipage}[t]{0.235\linewidth}\poptile{popgreen}{Fiche bilan}{l'essentiel à retenir}\end{minipage}\par}

% Sommaire
\newtcolorbox{sommaire}{enhanced,colback=white,colframe=popink,boxrule=2.2pt,arc=10pt,
  drop shadow=popink,left=18pt,right=18pt,top=13pt,bottom=13pt,before skip=6pt,after skip=20pt,
  before upper={\poppill{Sommaire}{poppurple}{white}\par\vspace{9pt}\popsemi\fontsize{10.6}{14.5}\selectfont\color{popink}}}

% Signature de fin de cours — poussée tout en bas de la dernière page
\newcommand{\findecours}{%
  \par\vfill\nopagebreak
  \begin{center}\popsquiggle{poporange}\end{center}\vspace{4pt}
  \signatureonbuch
}
\AtBeginDocument{\def\llbracket{\mathopen{[\mkern-3.5mu[}}\def\rrbracket{\mathclose{]\mkern-3.5mu]}}} % intervalles d'entiers (glyphe ⟦ absent de la police)
% Glyphes absents de Fira Math : redéfinis à partir de glyphes disponibles
\AtBeginDocument{%
  \def\setminus{\mathbin{\backslash}}\let\smallsetminus\setminus
  \def\vdots{\mathord{\vbox{\baselineskip3.2pt\lineskiplimit0pt\kern4pt\hbox{.}\hbox{.}\hbox{.}}}}%
  \def\ddots{\mathinner{\mkern1mu\raise6.5pt\hbox{.}\mkern2mu\raise3.6pt\hbox{.}\mkern2mu\raise0.7pt\hbox{.}\mkern1mu}}%
  \def\triangle{\mathord{\scalebox{0.9}{$\symup{\Delta}$}}}%
  \def\nabla{\mathord{\raisebox{\depth}{\scalebox{1}[-1]{$\symup{\Delta}$}}}}%
  \def\checkmark{\popcheck{popdark}}%
  \def\star{\mathbin{\ast}}\def\bigstar{\mathord{\ast}}%
  \def\diamond{\mathbin{\scalebox{0.6}{$\Diamond$}}}%
  \def\bigcup{\mathop{\mathchoice{\raisebox{-0.3em}{\scalebox{1.7}{$\cup$}}}{\scalebox{1.25}{$\cup$}}{\cup}{\cup}}\displaylimits}%
  \def\bigcap{\mathop{\mathchoice{\raisebox{-0.3em}{\scalebox{1.7}{$\cap$}}}{\scalebox{1.25}{$\cap$}}{\cap}{\cap}}\displaylimits}%
  \def\lesseqqgtr{\mathrel{\lessgtr}}\def\bigoplus{\mathop{\oplus}\displaylimits}%
}
\providecommand{\ssi}{si et seulement si} % abréviation fréquente
\AtBeginDocument{\providecommand{\wideparen}{\overparen}} % arc de cercle

%% FIGFIX-BEGIN v5 — correctifs globaux des figures (docs/onbuch-generation/tools/figfix)
\usepackage{adjustbox}\usepackage{etoolbox}\tcbuselibrary{raster}
% toute figure TikZ/pgfplots plus large que la ligne est réduite à la largeur disponible
\BeforeBeginEnvironment{tikzpicture}{\begin{adjustbox}{max width=\linewidth}}
\AfterEndEnvironment{tikzpicture}{\end{adjustbox}}
% légendes pgfplots lisibles par-dessus les courbes
\pgfplotsset{every axis/.append style={legend style={fill=white,fill opacity=0.92,text opacity=1,draw=black!15,inner sep=3pt}}}
% étiquettes d'arêtes (syntaxe edge["..."]) avec fond blanc
\tikzset{every edge quotes/.append style={fill=white,inner sep=1.2pt}}
%% FIGFIX-END
\begin{document}

\pagegarde{Production et traduction de textes sur la situation de la femme et les autres thèmes du chapitre 1}{Allemand}{Tle A}{Chapitre 1 : Vie familiale et sociale}{AL05}{2 h}{Leçon mixte de production et de traduction de textes sur la situation de la femme et les thèmes de la vie familiale et sociale. Elle combine lecture d'un texte support, lexique thématique, révisions grammaticales (ordre des mots, cas, temps du passé) et entraînement au thème et à la version en vue du Bac.
\vspace{4pt}
\begin{multicols}{2}\small
\begin{itemize}
\item À la fin de cette leçon, je sais comprendre et commenter un texte allemand sur la situation de la femme.
\item Je sais utiliser le lexique de la famille, du mariage, du travail et de l'égalité hommes-femmes.
\item Je sais respecter l'ordre des mots allemand (verbe en 2e position, verbe à la fin en subordonnée) dans mes productions.
\item Je sais employer correctement les quatre cas (Nominativ, Akkusativ, Dativ, Genitiv) après verbes et prépositions.
\item Je sais raconter au passé en choisissant entre Perfekt, Präteritum et Plusquamperfekt.
\item Je sais rédiger une lettre personnelle, un court essai argumenté et un dialogue en allemand.
\end{itemize}\end{multicols}}

\begin{sommaire}
\begin{multicols}{2}
\begin{enumerate}
\item Texte support : Die Frau in der heutigen Gesellschaft
\item Lexique thématique : famille, mariage, travail, égalité
\item Grammaire 1 : ordre des mots et cas
\item Grammaire 2 : les temps du passé
\item Structures de production : lettre, essai, dialogue
\item Traduction : thème et version
\item Activité d'intégration
\item Méthodes et exercices résolus
\item Exercices
\item Corrigés des exercices
\item Fiche bilan
\end{enumerate}
\end{multicols}
\end{sommaire}

\begin{objectifs}
\begin{itemize}
\item À la fin de cette leçon, je sais comprendre et commenter un texte allemand sur la situation de la femme.
\item Je sais utiliser le lexique de la famille, du mariage, du travail et de l'égalité hommes-femmes.
\item Je sais respecter l'ordre des mots allemand (verbe en 2e position, verbe à la fin en subordonnée) dans mes productions.
\item Je sais employer correctement les quatre cas (Nominativ, Akkusativ, Dativ, Genitiv) après verbes et prépositions.
\item Je sais raconter au passé en choisissant entre Perfekt, Präteritum et Plusquamperfekt.
\item Je sais rédiger une lettre personnelle, un court essai argumenté et un dialogue en allemand.
\item Je sais traduire de courtes phrases et un texte du français vers l'allemand (thème) et de l'allemand vers le français (version).
\item Je sais relire et corriger ma production : orthographe, majuscules, ordre des mots, cas.
\end{itemize}
\end{objectifs}

\begin{prerequis}
\begin{itemize}
\item Conjugaison des verbes forts et faibles au présent et au Perfekt (Première)
\item Déclinaison des articles et des adjectifs (Seconde/Première)
\item Vocabulaire de base de la famille et de la vie quotidienne (Première)
\item Notions d'argumentation : pour/contre, connecteurs logiques simples (Première)
\item Structure de la phrase déclarative et interrogative (Seconde)
\end{itemize}
\end{prerequis}

\newpage

\coursec{Texte support : Die Frau in der heutigen Gesellschaft}

Aminatou, élève de Terminale à Yaoundé, veut écrire une lettre à une amie allemande de Stuttgart pour lui raconter le mariage de sa grande sœur et débattre de la place des femmes dans la société camerounaise. En Allemagne, en Autriche et en Suisse, l'égalité entre hommes et femmes est un sujet de débat public permanent : congé parental, salaires, partage des tâches ménagères. Comment exprimer tout cela en allemand correct, avec le bon ordre des mots et les bons cas ? Cette leçon vous donne les outils pour produire et traduire des textes sur la vie familiale et sociale.

\textbf{Problématique :} comment lire, comprendre, puis produire et traduire un texte sur la situation de la femme dans les pays germanophones et au Cameroun ?

\courssub{Lecture globale : de quoi parle le texte ?}

Avant toute lecture détaillée, un bon lecteur pose trois questions au texte : \emph{de quel type est-il ?}, \emph{d'où vient-il ?}, \emph{à qui s'adresse-t-il ?} Lisez d'abord le texte en entier, sans dictionnaire.

\begin{textebox}[Die Frau in der heutigen Gesellschaft]
Früher war die Rolle der Frau klar: Sie blieb zu Hause, kümmerte sich um die Kinder und kochte für die Familie. Man sagte damals: „Kinder, Küche, Kirche" — das war die Welt der Frau. Heute hat sich viel geändert. Die meisten Frauen in Deutschland, Österreich und der Schweiz gehen zur Schule, machen eine Ausbildung und arbeiten in einem Beruf. Viele studieren an der Universität und werden Ärztinnen, Lehrerinnen oder Ingenieurinnen.

Trotzdem verdienen Frauen oft weniger als Männer und machen mehr Hausarbeit. Einerseits gibt es Gesetze für die Gleichberechtigung, andererseits ist der Alltag oft noch ungerecht. Deshalb diskutieren Politiker, Journalisten und Familien über dieses Thema. Viele junge Paare teilen sich heute die Aufgaben: Der Vater kann auch zu Hause bleiben und sich um die Kinder kümmern. Die Gleichberechtigung ist ein Ziel, aber der Weg ist noch lang.
\end{textebox}

\begin{definition}[Glossaire du texte]
\begin{itemize}
\item \cle{sich kümmern um} (+ Akk.) = s'occuper de
\item \all{die Küche, -n} = la cuisine ; \all{die Kirche, -n} = l'église
\item \all{die Ausbildung, -en} = la formation (professionnelle)
\item \all{der Beruf, -e} = le métier, la profession
\item \all{verdienen} = gagner (un salaire)
\item \all{die Hausarbeit} (Sg.) = les tâches ménagères
\item \all{das Gesetz, -e} = la loi
\item \all{die Gleichberechtigung} (Sg.) = l'égalité des droits
\item \all{ungerecht} = injuste
\item \all{sich teilen} = se partager
\item \all{die Aufgabe, -n} = la tâche
\item \all{das Ziel, -e} = le but
\end{itemize}
\end{definition}

\begin{methode}[Lecture globale d'un article de journal]
\begin{enumerate}
\item \textbf{Le titre} annonce le thème : \all{Die Frau in der heutigen Gesellschaft} = « la femme dans la société d'aujourd'hui ». Thème : la situation de la femme.
\item \textbf{Le type de texte} : c'est un article de journal (un reportage ou un texte d'information), écrit à la troisième personne, avec des faits et des constats, pas une lettre ni un dialogue.
\item \textbf{La source} : un journal ou un magazine germanophone (Allemagne, Autriche, Suisse) — le texte cite explicitement ces trois pays.
\item \textbf{Le destinataire} : le grand public, des lecteurs adultes intéressés par les questions de société.
\item \textbf{L'intention} : informer et faire réfléchir sur les progrès et les limites de l'égalité.
\end{enumerate}
\end{methode}

\courssub{Lecture analytique : la femme d'hier et d'aujourd'hui}

Le texte est construit sur une \textbf{opposition temporelle} : \all{früher} (autrefois) contre \all{heute} (aujourd'hui). Observons les deux portraits.

\textbf{Hier — le modèle traditionnel.} \all{Früher war die Rolle der Frau klar: Sie blieb zu Hause, kümmerte sich um die Kinder und kochte für die Familie.} « Autrefois, le rôle de la femme était clair : elle restait à la maison, s'occupait des enfants et cuisinait pour la famille. » La formule \all{„Kinder, Küche, Kirche"} (« les enfants, la cuisine, l'église ») résume ce modèle des trois K : la femme était définie par la maison.

\textbf{Aujourd'hui — le changement.} \all{Heute hat sich viel geändert.} « Aujourd'hui, beaucoup de choses ont changé. » Notez le verbe réfléchi \all{sich ändern} (changer) au \cle{Perfekt} : \all{hat sich geändert}. Les femmes vont à l'école, font une formation, exercent un métier : \all{Die meisten Frauen gehen zur Schule, machen eine Ausbildung und arbeiten in einem Beruf.}

\textbf{Mais des problèmes persistent.} \all{Trotzdem verdienen Frauen oft weniger als Männer und machen mehr Hausarbeit.} « Pourtant, les femmes gagnent souvent moins que les hommes et font plus de tâches ménagères. » La comparaison s'exprime avec \all{weniger als} (moins que) et \all{mehr} (plus).

\begin{center}
\begin{tabularx}{\linewidth}{|X|X|X|}
\hline
\rowcolor{popblueL} Aspect & Früher (hier) & Heute (aujourd'hui) \\
\hline
Rôle & Hausfrau zu Hause & Beruf und Familie \\
\hline
Éducation & keine Ausbildung & Schule, Ausbildung, Universität \\
\hline
Tâches & Kinder, Küche, Kirche & Paare teilen sich die Aufgaben \\
\hline
Droits & keine Gleichberechtigung & Gesetze für Gleichberechtigung \\
\hline
Problème & Abhängigkeit vom Mann & weniger Lohn, mehr Hausarbeit \\
\hline
\end{tabularx}
\end{center}

\begin{attention}[Piège : sich kümmern um]
Ne traduisez jamais mot à mot \all{sich kümmern um} : c'est un \textbf{verbe réfléchi avec préposition fixe} (\all{um} + Akkusativ). On dit : \all{Sie kümmert sich um die Kinder.} « Elle s'occupe des enfants. » Jamais \all{*Sie kümmert die Kinder}. Autres verbes de ce type : \all{sich interessieren für} (s'intéresser à), \all{sich freuen über} (se réjouir de), \all{warten auf} (attendre). Apprenez toujours le verbe \emph{avec} sa préposition.
\end{attention}

\courssub{Comparaison avec la situation au Cameroun}

Pour réagir à un texte (compétence d'expression orale et écrite), il faut savoir comparer avec son propre pays. Voici le vocabulaire utile pour parler du Cameroun en allemand :

\begin{itemize}
\item \all{In Kamerun gehen heute viele Mädchen zur Schule.} « Au Cameroun, beaucoup de filles vont aujourd'hui à l'école. »
\item \all{Früher heirateten viele Mädchen sehr jung.} « Autrefois, beaucoup de filles se mariaient très jeunes. » (\all{heiraten} = se marier, épouser)
\item \all{Viele Frauen arbeiten auf dem Markt oder auf dem Feld.} « Beaucoup de femmes travaillent au marché ou aux champs. »
\item \all{Auch in Kamerun ist der Weg zur Gleichberechtigung noch lang.} « Au Cameroun aussi, le chemin vers l'égalité des droits est encore long. »
\end{itemize}

Structure de comparaison à réutiliser : \all{In Deutschland ist es ..., aber in Kamerun ...} / \all{Wie in Deutschland, auch in Kamerun ...} Notez que \all{Kamerun} est neutre : \all{in Kamerun} (sans article), comme \all{in Deutschland} ; mais on dit \all{in der Schweiz} (féminin, avec article !).

\courssub{Questions de compréhension}

Répondez d'abord aux questions globales (les « W-Fragen »), puis aux questions de détail.

\begin{exercice}[Compréhension du texte]
\begin{enumerate}
\item \textbf{Wer?} — Von wem spricht der Text? (De qui parle le texte ?)
\item \textbf{Was?} — Was ist das Thema des Textes?
\item \textbf{Wo?} — In welchen Ländern spielt der Text?
\item \textbf{Wann?} — Welche zwei Zeiten vergleicht der Text?
\item \textbf{Warum?} — Warum diskutieren viele Leute über dieses Thema?
\item Wie war die Rolle der Frau früher?
\item Was machen die meisten Frauen heute?
\item Welches Problem gibt es trotz der Gesetze?
\end{enumerate}
\end{exercice}

\begin{corrige}[Exercice : Compréhension du texte]
\begin{enumerate}
\item Der Text spricht von den Frauen. (Le texte parle des femmes.)
\item Das Thema ist die Situation der Frau in der Gesellschaft — früher und heute.
\item Der Text spielt in Deutschland, Österreich und der Schweiz.
\item Er vergleicht \all{früher} (autrefois) und \all{heute} (aujourd'hui).
\item Weil der Alltag oft noch ungerecht ist: Frauen verdienen weniger und machen mehr Hausarbeit.
\item Früher blieb die Frau zu Hause, kümmerte sich um die Kinder und kochte.
\item Heute gehen die meisten Frauen zur Schule, machen eine Ausbildung und arbeiten in einem Beruf.
\item Trotz der Gesetze verdienen Frauen oft weniger als Männer.
\end{enumerate}
\end{corrige}

\courssub{Relevé des connecteurs du texte}

Les connecteurs sont les « charnières » du texte : ils organisent l'argumentation. Relevez-les systématiquement, vous les réutiliserez dans vos propres productions.

\begin{center}
\begin{tabularx}{\linewidth}{|X|X|X|}
\hline
\rowcolor{popblueL} Connecteur & Français & Fonction \\
\hline
\all{früher – heute} & autrefois – aujourd'hui & opposition temporelle \\
\hline
\all{einerseits – andererseits} & d'une part – d'autre part & balance d'arguments \\
\hline
\all{deshalb} & c'est pourquoi, donc & conséquence \\
\hline
\all{trotzdem} & pourtant, néanmoins & concession \\
\hline
\all{aber} & mais & opposition simple \\
\hline
\end{tabularx}
\end{center}

\begin{propriete}[Connecteurs et place du verbe]
\all{Deshalb} et \all{trotzdem} occupent la \textbf{première place} de la phrase : le verbe conjugué vient ensuite, puis le sujet (inversion !). \all{Deshalb diskutieren viele Leute über dieses Thema.} « C'est pourquoi beaucoup de gens discutent de ce sujet. » En français on dit « donc beaucoup de gens discutent » ; en allemand : \all{deshalb} + verbe + sujet. En revanche, \all{aber} n'occupe pas une position : \all{Der Weg ist noch lang, aber das Ziel ist klar.}
\end{propriete}

\begin{exoresolu}[Connecteurs]
Énoncé : Complétez avec \all{deshalb}, \all{trotzdem} ou \all{einerseits ... andererseits}. — \all{Frauen haben heute gute Berufe. \dots\ verdienen sie oft weniger als Männer.}
\tcblower
Solution : il y a une opposition entre « avoir de bons métiers » et « gagner moins » : c'est une concession. La bonne réponse est \all{trotzdem} : \all{Frauen haben heute gute Berufe. Trotzdem verdienen sie oft weniger als Männer.} Vérifiez l'inversion : \all{verdienen} (verbe) vient avant \all{sie} (sujet). \all{Deshalb} exprimerait une conséquence, ce qui serait illogique ici.
\end{exoresolu}

\courssub{Carte mentale : le vocabulaire du thème}

Pour mémoriser le lexique, organisez-le en carte mentale autour du thème central.

\begin{popfigure}
\begin{center}
\begin{tikzpicture}[mindmap, grow cyclic, every node/.style={concept, text=white, font=\small},
root concept/.append style={concept color=popblue, font=\small\bfseries},
level 1/.append style={level distance=3.2cm, sibling angle=90},
level 2/.append style={level distance=2.2cm, sibling angle=45, concept color=popmuted}]
\node[root concept]{Die Frau heute}
  child[concept color=poporange]{ node {Familie}
    child { node {heiraten} }
    child { node {die Kinder} }
  }
  child[concept color=popgreen]{ node {Beruf}
    child { node {verdienen} }
    child { node {arbeiten} }
  }
  child[concept color=poppurple]{ node {Bildung}
    child { node {die Schule} }
    child { node {die Ausbildung} }
  }
  child[concept color=poppink]{ node {Rechte}
    child { node {das Gesetz} }
    child { node {die Gleich-berechtigung} }
  };
\end{tikzpicture}
\end{center}
\legende{Carte mentale du vocabulaire : la femme aujourd'hui (Familie, Beruf, Bildung, Rechte)}
\end{popfigure}

\begin{aretenir}[L'essentiel de la section]
\begin{itemize}
\item Un article de journal s'identifie par son titre, sa source, son destinataire et son intention.
\item Le texte oppose \all{früher} (modèle „Kinder, Küche, Kirche") et \all{heute} (école, formation, métier), tout en montrant les inégalités qui persistent.
\item Les connecteurs \all{früher – heute}, \all{einerseits – andererseits}, \all{deshalb}, \all{trotzdem} structurent l'argumentation ; après \all{deshalb} et \all{trotzdem}, inversion verbe-sujet.
\item \all{sich kümmern um} + Akkusativ = s'occuper de ; à apprendre avec sa préposition.
\end{itemize}
\end{aretenir}

\begin{savaistu}[Le congé parental en Allemagne]
En Allemagne, le congé parental (\all{die Elternzeit}) peut être pris par le père \emph{ou} par la mère, et même partagé entre les deux. Aujourd'hui, environ 25 \% des pères allemands prennent un congé parental pour s'occuper de leur bébé. On parle même de \all{Vätermonate} (« mois des pères »). C'est un signe concret du changement des rôles dans la famille allemande — un excellent argument pour vos essais et débats !
\end{savaistu}

\coursec{Lexique thématique : famille, mariage, travail, égalité}

Dans la section précédente, nous avons lu et commenté le texte \all{Die Frau in der heutigen Gesellschaft}. Ce texte était riche en vocabulaire : il parlait de la \all{Familie}, de l'\all{Ehe}, du \all{Beruf}, de la \all{Gleichberechtigung} et de la \all{Hausarbeit}. Pour pouvoir, à votre tour, produire un essai, une lettre ou un dialogue sur la situation de la femme — et pour traduire correctement dans les deux sens —, il faut d'abord posséder solidement ce lexique. C'est l'objet de cette section : nous allons organiser le vocabulaire en cinq champs lexicaux, l'observer dans des phrases authentiques, l'apprendre avec ses articles et ses pluriels, puis nous entraîner à l'employer dans des expressions de débat.

\courssub{Observer d'abord : le vocabulaire en contexte}

Avant toute liste, observons comment ces mots vivent dans des phrases. Lisez ce court portrait :

\begin{textebox}[Ein Porträt : Mama Esther aus Yaoundé]
Meine Mutter Esther ist 45 Jahre alt und arbeitet als Verkäuferin auf dem Markt in Yaoundé. Sie ist mit meinem Vater verheiratet; sie haben vor zwanzig Jahren geheiratet. Meine Eltern haben vier Kinder, zwei Söhne und zwei Töchter. Meine Großeltern leben im Dorf, und mein jüngster Bruder besucht sie oft.

Meine Mutter verdient ihr eigenes Geld, aber sie macht auch fast die ganze Hausarbeit: Sie kocht, putzt und kauft jeden Morgen ein. Mein Vater kümmert sich um die Schulsachen der Kinder. Meiner Meinung nach ist das nicht ganz gerecht: Einerseits arbeitet meine Mutter viel, andererseits hat sie weniger Freizeit als mein Vater. Es ist wichtig, dass Männer und Frauen die Aufgaben im Haushalt teilen. Die Gleichberechtigung beginnt in der Familie!
\end{textebox}

\textbf{Glossaire :} \all{der Markt, die Märkte} (le marché) ; \all{verheiratet sein mit} (être marié avec) ; \all{das Dorf, die Dörfer} (le village) ; \all{das Geld} (l'argent) ; \all{die Freizeit} (le temps libre) ; \all{gerecht} (juste) ; \all{teilen} (partager).

Relevez dans ce texte tous les noms : ils portent tous une majuscule et un article. C'est la première grande loi du vocabulaire allemand.

\begin{propriete}[La règle d'or : article + pluriel]
Tout nom allemand s'apprend avec \cle{son article} (son genre) et \cle{son pluriel}. On note : \all{die Ehe (-n)}, \all{das Recht (-e)}, \all{die Aufgabe (-n)}, \all{der Beruf (-e)}. Un nom appris sans article est un nom à moitié appris : sans le genre, impossible de décliner correctement (\all{die Ehe} → \all{der Ehe}, \all{das Recht} → \all{des Rechts}).
\end{propriete}

\courssub{Champ 1 : la famille — die Familie}

\begin{definition}[Les mots de la famille]
\begin{itemize}
\item \all{die Familie (-n)} : la famille
\item \all{die Eltern} (pluriel seulement) : les parents
\item \all{der Vater (¨-)} / \all{die Mutter (¨-)} : le père / la mère
\item \all{die Großeltern} (pl.) : les grands-parents ; \all{der Großvater}, \all{die Großmutter}
\item \all{der Ehemann (¨-er)} : le mari, l'époux ; \all{die Ehefrau (-en)} : la femme, l'épouse
\item \all{das Kind (-er)} : l'enfant ; \all{der Sohn (¨-e)} : le fils ; \all{die Tochter (¨-)} : la fille
\item \all{der Enkel (-)} / \all{die Enkelin (-nen)} : le petit-fils / la petite-fille
\item \all{der Haushalt (-e)} : le ménage, le foyer
\end{itemize}
\end{definition}

\begin{exemplebox}[La famille en phrases]
\all{Meine Großeltern leben in Douala.} « Mes grands-parents habitent à Douala. »

\all{Ihr Ehemann arbeitet bei einer Bank.} « Son mari travaille dans une banque. »

\all{Die Enkelin besucht ihre Großmutter jeden Sonntag.} « La petite-fille rend visite à sa grand-mère chaque dimanche. »

\all{In diesem Haushalt leben sechs Personen.} « Six personnes vivent dans ce foyer. »
\end{exemplebox}

\begin{attention}[Eltern, Großeltern : des pluriels sans singulier courant]
\all{Die Eltern} n'a pas de singulier usuel (on dit \all{der Vater} ou \all{die Mutter}). Attention à l'accord du verbe : \all{Meine Eltern \textbf{wohnen} in Yaoundé.} et jamais \emph{meine Eltern wohnt}. De même, \all{der Vater} et \all{die Mutter} font leur pluriel en Umlaut : \all{die Väter}, \all{die Mütter}.
\end{attention}

\courssub{Champ 2 : mariage et vie de couple — die Ehe}

\begin{definition}[Les mots du mariage]
\begin{itemize}
\item \all{heiraten (+ Akk.)} : épouser, se marier avec
\item \all{die Hochzeit (-en)} : le mariage (la fête, la cérémonie)
\item \all{die Ehe (-n)} : le mariage (l'institution, l'union)
\item \all{verlobt sein (mit)} : être fiancé (avec) ; \all{die Verlobung (-en)} : les fiançailles
\item \all{verheiratet sein (mit)} : être marié (avec)
\item \all{sich scheiden lassen} : divorcer ; \all{die Scheidung (-en)} : le divorce
\item \all{der Partner (-)} / \all{die Partnerin (-nen)} : le partenaire / la partenaire
\end{itemize}
\end{definition}

\begin{propriete}[heiraten : un verbe à construire sans préposition]
\all{heiraten} se construit avec l'\cle{accusatif direct}, sans préposition : \all{Er hat sie geheiratet.} « Il l'a épousée. » Conjugaison au Präsens : \all{ich heirate, du heiratest, er heiratet, wir heiraten}. Au Perfekt : \all{hat geheiratet} (verbe faible, régulier). Le participe \all{verheiratet} s'emploie avec \all{sein} + \all{mit} : \all{Sie ist mit einem Lehrer verheiratet.} « Elle est mariée à un enseignant. »
\end{propriete}

\begin{attention}[Le piège du francophone : « se marier avec »]
En français on dit « se marier \emph{avec} quelqu'un », et l'on est tenté de traduire : \emph{er hat mit ihr geheiratet}. C'est une erreur ! Retenez :

\begin{center}
\begin{tabularx}{\linewidth}{|X|X|X|}
\hline
\rowcolor{popblueL} Français & Deutsch & Remarque \\
\hline
Il l'a épousée. & Er hat sie geheiratet. & \all{heiraten} + Akkusativ, sans préposition \\
\hline
Il s'est marié avec elle. & Er hat sie geheiratet. / Er ist mit ihr verheiratet. & deux formulations possibles \\
\hline
Elle est mariée avec lui. & Sie ist mit ihm verheiratet. & ici \all{mit} est obligatoire \\
\hline
Ils ont divorcé. & Sie haben sich scheiden lassen. & \all{sich lassen} + infinitif \\
\hline
\end{tabularx}
\end{center}
\end{attention}

\begin{exemplebox}[Mariage et couple en phrases]
\all{Sie haben letztes Jahr geheiratet.} « Ils se sont mariés l'année dernière. »

\all{Die Hochzeit war sehr schön; zweihundert Gäste waren da.} « Le mariage était très beau ; deux cents invités étaient là. »

\all{Nach zehn Jahren Ehe haben sie sich scheiden lassen.} « Après dix ans de mariage, ils ont divorcé. »

\all{Meine Schwester ist mit ihrem Kollegen verlobt.} « Ma sœur est fiancée à son collègue. »
\end{exemplebox}

\courssub{Champ 3 : travail et formation — der Beruf}

\begin{definition}[Les mots du travail]
\begin{itemize}
\item \all{der Beruf (-e)} : le métier, la profession
\item \all{arbeiten} : travailler ; \all{die Arbeit (-en)} : le travail
\item \all{die Ausbildung (-en)} : la formation (professionnelle)
\item \all{die Stelle (-n)} : le poste, l'emploi
\item \all{verdienen} : gagner (de l'argent) ; \all{das Gehalt (¨-er)} : le salaire
\item \all{der Chef (-s)} / \all{die Chefin (-nen)} : le chef / la cheffe
\item \all{arbeitslos} : au chômage ; \all{die Arbeitslosigkeit} : le chômage
\end{itemize}
\end{definition}

Remarquez le féminin des noms de métiers : on ajoute \all{-in} au masculin, souvent avec Umlaut : \all{der Lehrer → die Lehrerin}, \all{der Arzt → die Ärztin}, \all{der Chef → die Chefin}. C'est un point essentiel pour parler de la situation de la femme !

\begin{exemplebox}[Le travail en phrases]
\all{Meine Mutter arbeitet als Lehrerin.} « Ma mère travaille comme enseignante. »

\all{Nach der Ausbildung hat sie eine Stelle in einem Krankenhaus gefunden.} « Après sa formation, elle a trouvé un poste dans un hôpital. »

\all{Die Frauen verdienen oft weniger als die Männer.} « Les femmes gagnent souvent moins que les hommes. »

\all{Ihr Gehalt ist nicht hoch, aber sie liebt ihren Beruf.} « Son salaire n'est pas élevé, mais elle aime son métier. »

\all{Viele junge Leute sind arbeitslos.} « Beaucoup de jeunes sont au chômage. »
\end{exemplebox}

\begin{attention}[arbeiten als, pas « arbeiten wie »]
Pour dire « travailler comme », l'allemand utilise \all{als} sans article : \all{Sie arbeitet als Ärztin.} « Elle travaille comme médecin. » Ne dites pas \emph{sie arbeitet wie eine Ärztin} (cela signifierait « elle travaille à la manière d'un médecin »). Et notez l'absence d'article après \all{als} pour les professions.
\end{attention}

\courssub{Champ 4 : égalité et droits — die Gleichberechtigung}

\begin{definition}[Les mots de l'égalité]
\begin{itemize}
\item \all{die Gleichberechtigung} (sans pluriel) : l'égalité des droits
\item \all{das Recht (-e)} : le droit ; \all{die Rechte} (pl.) : les droits
\item \all{die Pflicht (-en)} : le devoir
\item \all{gerecht} / \all{ungerecht} : juste / injuste
\item \all{die Rolle (-n)} : le rôle ; \all{das Rollenbild (-er)} : le modèle de rôle
\item \all{die Tradition (-en)} : la tradition ; \all{traditionell} : traditionnel
\item \all{die Emanzipation} : l'émancipation
\end{itemize}
\end{definition}

\begin{savaistu}[Un mot composé typiquement allemand]
\all{Die Gleichberechtigung} est un mot composé : \all{gleich} (égal) + \all{das Recht} (le droit) + le suffixe \all{-ung}. L'allemand adore ces longs mots : \all{die Frauenrechte} (les droits des femmes), \all{das Rollenbild} (l'image du rôle). Le genre d'un mot composé est toujours celui du \textbf{dernier} élément : \all{das Recht} → \all{die Gleichberechtigung} ? Non ! Attention : le suffixe \all{-ung} est toujours féminin, donc \all{die Gleichberechtigung}. En revanche \all{das Frauenrecht} reste neutre, comme \all{das Recht}.
\end{savaistu}

\begin{exemplebox}[Égalité et droits en phrases]
\all{Die Gleichberechtigung von Mann und Frau steht in der Verfassung.} « L'égalité des droits entre l'homme et la femme figure dans la Constitution. »

\all{Jeder Mensch hat Rechte und Pflichten.} « Tout être humain a des droits et des devoirs. »

\all{Es ist ungerecht, dass sie für dieselbe Arbeit weniger verdient.} « Il est injuste qu'elle gagne moins pour le même travail. »

\all{Die traditionelle Rolle der Frau verändert sich.} « Le rôle traditionnel de la femme change. »
\end{exemplebox}

\courssub{Champ 5 : tâches et vie quotidienne — die Hausarbeit}

\begin{definition}[Les mots du quotidien]
\begin{itemize}
\item \all{die Hausarbeit} (sans pluriel) : les tâches ménagères
\item \all{kochen} : cuisiner, faire la cuisine
\item \all{putzen} : nettoyer, faire le ménage
\item \all{einkaufen} : faire les courses (verbe à particule séparable : \all{ich kaufe ein})
\item \all{die Kinder erziehen} : élever les enfants (\all{erziehen} : verbe fort, \all{erzog, hat erzogen})
\item \all{sich kümmern um (+ Akk.)} : s'occuper de
\end{itemize}
\end{definition}

\begin{exemplebox}[La vie quotidienne en phrases]
\all{In vielen Familien macht die Frau die ganze Hausarbeit allein.} « Dans beaucoup de familles, la femme fait seule toutes les tâches ménagères. »

\all{Mein Vater kauft jeden Samstag auf dem Markt ein.} « Mon père fait les courses au marché chaque samedi. »

\all{Beide Eltern kümmern sich um die Kinder.} « Les deux parents s'occupent des enfants. »

\all{Wer kocht heute? — Heute putze ich, und du kochst!} « Qui cuisine aujourd'hui ? — Aujourd'hui je nettoie, et toi tu cuisines ! »
\end{exemplebox}

\courssub{Carte mentale et tableau récapitulatif}

\begin{popfigure}
\begin{tikzpicture}[mindmap, grow cyclic, every node/.style={concept, font=\small}, concept color=popblueL, root concept/.append style={concept color=popblue, text=white, font=\bfseries}, level 1/.append style={level distance=3.4cm, sibling angle=72}]
\node[root concept]{Familie und Gesellschaft}
  child[concept color=popgreenL] { node[align=center]{die Familie\\ Eltern, Kinder,\\ Großeltern} }
  child[concept color=poporangeL] { node[align=center]{die Ehe\\ heiraten,\\ Hochzeit, Scheidung} }
  child[concept color=poppinkL] { node[align=center]{der Beruf\\ arbeiten, Stelle,\\ Gehalt, Chefin} }
  child[concept color=popgoldL] { node[align=center]{die Gleichberechtigung\\ Rechte, Pflichten,\\ gerecht} }
  child[concept color=poppurpleL] { node[align=center]{die Hausarbeit\\ kochen, putzen,\\ einkaufen} };
\end{tikzpicture}
\legende{Carte mentale des cinq champs lexicaux du chapitre}
\end{popfigure}

\begin{center}
\begin{tabularx}{\linewidth}{|X|X|X|X|X|}
\hline
\rowcolor{popblueL} Famille & Mariage & Travail & Égalité & Quotidien \\
\hline
die Familie (-n) la famille & die Ehe (-n) le mariage & der Beruf (-e) le métier & das Recht (-e) le droit & die Hausarbeit le ménage \\
\hline
die Eltern (pl.) les parents & die Hochzeit (-en) la noce & die Stelle (-n) le poste & die Pflicht (-en) le devoir & kochen cuisiner \\
\hline
die Großeltern (pl.) les grands-parents & heiraten épouser & die Ausbildung (-en) la formation & gerecht / ungerecht juste / injuste & putzen nettoyer \\
\hline
der Ehemann (¨-er) le mari & sich scheiden lassen divorcer & verdienen gagner & die Rolle (-n) le rôle & einkaufen faire les courses \\
\hline
die Ehefrau (-en) l'épouse & die Scheidung (-en) le divorce & das Gehalt (¨-er) le salaire & die Tradition (-en) la tradition & die Kinder erziehen élever les enfants \\
\hline
das Kind (-er) l'enfant & der Partner (-) le partenaire & die Chefin (-nen) la cheffe & die Emanzipation l'émancipation & sich kümmern um s'occuper de \\
\hline
der Haushalt (-e) le foyer & verlobt sein être fiancé & arbeitslos au chômage & die Gleichberechtigung l'égalité des droits & der Alltag le quotidien \\
\hline
\end{tabularx}
\end{center}

\courssub{Expressions utiles pour le débat}

Pour rédiger un essai ou participer à un débat sur la situation de la femme, il ne suffit pas de connaître des mots : il faut des \cle{phrases toutes faites} pour exprimer son opinion.

\begin{definition}[Redemittel : donner son avis]
\begin{itemize}
\item \all{Meiner Meinung nach ...} : à mon avis... (suivi du verbe en 2\textsuperscript{e} position : \all{Meiner Meinung nach ist das ungerecht.})
\item \all{Ich bin der Meinung, dass ...} : je suis d'avis que... (subordonnée : verbe à la fin)
\item \all{Einerseits ..., andererseits ...} : d'une part..., d'autre part...
\item \all{Es ist wichtig, dass ...} : il est important que... (subordonnée, verbe à la fin)
\item \all{Ich finde, dass ...} : je trouve que...
\item \all{Man muss ...} : il faut...
\end{itemize}
\end{definition}

\begin{exemplebox}[Le débat en phrases]
\all{Meiner Meinung nach müssen Männer mehr im Haushalt helfen.} « À mon avis, les hommes doivent davantage aider au foyer. »

\all{Ich bin der Meinung, dass Mädchen und Jungen dieselbe Ausbildung brauchen.} « Je suis d'avis que les filles et les garçons ont besoin de la même formation. »

\all{Einerseits arbeiten viele Frauen, andererseits machen sie noch die ganze Hausarbeit.} « D'une part beaucoup de femmes travaillent, d'autre part elles font encore tout le ménage. »

\all{Es ist wichtig, dass jeder Mensch seine Rechte kennt.} « Il est important que chacun connaisse ses droits. »
\end{exemplebox}

\begin{methode}[Mémoriser durablement : la phrase personnelle]
Pour retenir chaque mot nouveau, suivez ces étapes :
\begin{enumerate}
\item Écrivez le mot avec son article et son pluriel : \all{die Scheidung (-en)}.
\item Construisez une phrase \textbf{personnelle}, vraie pour vous ou votre entourage : \all{In meiner Familie gibt es keine Scheidung.} « Dans ma famille, il n'y a pas de divorce. »
\item Relisez votre phrase à voix haute et vérifiez : majuscule au nom, verbe en 2\textsuperscript{e} position, cas corrects.
\item Le lendemain, cachez la colonne allemande de votre liste et retrouvez les mots à partir du français.
\end{enumerate}
Une phrase personnelle vaut dix répétitions mécaniques : le cerveau retient ce qui a du sens pour lui.
\end{methode}

\begin{exoresolu}[Complétez avec le mot qui convient]
Énoncé. Complétez les phrases avec : \all{Gehalt}, \all{Scheidung}, \all{Gleichberechtigung}, \all{Haushalt}, \all{Ausbildung}.

a) Nach der \dots\ hat er eine Stelle als Mechaniker gefunden.

b) Sie verdient 200\,000 Francs im Monat; ihr \dots\ ist nicht schlecht.

c) Für mich ist die \dots\ von Mann und Frau sehr wichtig.

d) Im \dots\ hilft jeder: Die Kinder putzen, der Vater kocht.

e) Nach der \dots\ lebt sie allein mit den zwei Kindern.

\tcblower
Solution détaillée.

a) \all{Nach der \textbf{Ausbildung} hat er eine Stelle als Mechaniker gefunden.} « Après sa formation, il a trouvé un poste comme mécanicien. » — C'est la formation qui précède logiquement le poste ; datif féminin : \all{der Ausbildung}.

b) \all{Sie verdient 200\,000 Francs im Monat; ihr \textbf{Gehalt} ist nicht schlecht.} « Elle gagne 200 000 francs par mois ; son salaire n'est pas mauvais. » — Ce qu'on gagne chaque mois, c'est le salaire, \all{das Gehalt}.

c) \all{Für mich ist die \textbf{Gleichberechtigung} von Mann und Frau sehr wichtig.} « Pour moi, l'égalité des droits entre l'homme et la femme est très importante. » — Mot-clé du chapitre, féminin à cause du suffixe \all{-ung}.

d) \all{Im \textbf{Haushalt} hilft jeder: Die Kinder putzen, der Vater kocht.} « Au foyer, tout le monde aide : les enfants nettoient, le père cuisine. » — \all{im} = \all{in dem}, datif masculin : \all{dem Haushalt}.

e) \all{Nach der \textbf{Scheidung} lebt sie allein mit den zwei Kindern.} « Après le divorce, elle vit seule avec les deux enfants. » — \all{die Scheidung}, datif féminin : \all{der Scheidung}.
\end{exoresolu}

\begin{exercice}[À vous de jouer]
1. Traduisez en allemand : a) Ma grand-mère s'occupe des enfants. b) Ils se sont mariés l'année dernière. c) À mon avis, les tâches ménagères sont l'affaire de tous. d) Il est important que les femmes gagnent le même salaire.

2. Rédigez trois phrases personnelles avec : \all{die Ausbildung}, \all{verdienen}, \all{sich kümmern um}.
\end{exercice}

\begin{corrige}[Exercice « À vous de jouer »]
1. a) \all{Meine Großmutter kümmert sich um die Kinder.} — \all{sich kümmern um} + Akkusativ ; verbe en 2\textsuperscript{e} position.

b) \all{Sie haben letztes Jahr geheiratet.} — Perfekt avec \all{haben}, sans préposition.

c) \all{Meiner Meinung nach ist die Hausarbeit die Sache aller.} — Après \all{meiner Meinung nach}, le verbe reste en 2\textsuperscript{e} position.

d) \all{Es ist wichtig, dass die Frauen dasselbe Gehalt verdienen.} — Subordonnée avec \all{dass} : verbe conjugué à la fin.

2. Exemples possibles : \all{Meine Schwester macht eine Ausbildung als Krankenschwester.} / \all{Mein Onkel verdient sein Geld als Taxifahrer in Douala.} / \all{Ich kümmere mich um meinen kleinen Bruder.}
\end{corrige}

\begin{attention}[Bilan des pièges lexicaux de cette section]
\begin{itemize}
\item « Se marier avec » ≠ \all{mit} + \all{heiraten} : dites \all{Er hat sie geheiratet.}
\item « Travailler comme » = \all{arbeiten als} (sans article) : \all{Sie arbeitet als Lehrerin.}
\item \all{die Eltern} est toujours pluriel : \all{meine Eltern wohnen}...
\item \all{einkaufen} est séparable : \all{Ich kaufe heute ein.}
\item Ne confondez pas \all{die Hochzeit} (la fête) et \all{die Ehe} (l'union) : \all{Die Hochzeit dauert einen Tag, die Ehe ein Leben.} « La noce dure un jour, le mariage une vie. »
\end{itemize}
\end{attention}

Vous disposez maintenant d'un lexique complet et structuré sur la famille, le mariage, le travail, l'égalité et la vie quotidienne. Ces mots seront vos outils dans les sections suivantes : nous allons d'abord réviser l'ordre des mots et les cas, puis les temps du passé, avant de les réinvestir dans la production de lettres, d'essais et de dialogues, et enfin dans la traduction en thème et en version.

\coursec{Grammaire 1 : ordre des mots et cas}

Dans la section précédente, nous avons enrichi notre vocabulaire sur la famille, le mariage, le travail et l'égalité. Mais un vocabulaire riche ne suffit pas : pour produire un texte correct sur la situation de la femme, il faut savoir \emph{assembler} les mots. C'est exactement l'objet de cette section : l'\cle{ordre des mots} (la place du verbe) et les \cle{cas} (la fonction des noms). Ce sont les deux piliers de la phrase allemande, et les deux sources principales d'erreurs des élèves francophones.

\courssub{Observer avant de comprendre : quatre phrases, quatre structures}

Lisons d'abord ces quatre phrases, toutes construites avec le même vocabulaire :

\begin{exemplebox}[Quatre structures fondamentales]
\begin{enumerate}
\item \all{Die Frau arbeitet heute.} — « La femme travaille aujourd'hui. »
\item \all{Heute arbeitet die Frau.} — « Aujourd'hui, la femme travaille. »
\item \all{Arbeitet die Frau heute?} — « La femme travaille-t-elle aujourd'hui ? »
\item \all{Ich weiß, dass die Frau heute arbeitet.} — « Je sais que la femme travaille aujourd'hui. »
\end{enumerate}
\end{exemplebox}

Observons : dans les phrases 1 et 2, le verbe conjugué \all{arbeitet} occupe la \textbf{deuxième position}, quoi qu'il y ait en tête de phrase. Dans la phrase 3 (question sans mot interrogatif), il passe en \textbf{première position}. Dans la phrase 4, introduite par \all{dass}, il va se placer \textbf{tout à la fin}. Le verbe conjugué est donc le pivot de la phrase allemande : sa place change selon le type de phrase, mais elle est \emph{toujours} prévisible.

\begin{propriete}[La position du verbe conjugué : les trois lois]
\begin{enumerate}
\item \textbf{Phrase déclarative} : le verbe conjugué est en \cle{2e position}. La première position peut être occupée par le sujet (\all{Die Frau arbeitet heute.}) ou par un complément (\all{Heute arbeitet die Frau.}) — dans ce cas, le sujet passe après le verbe : c'est l'\cle{inversion}.
\item \textbf{Question sans mot interrogatif} (réponse oui/non) : le verbe conjugué est en \cle{1re position} : \all{Arbeitet die Frau heute?} Avec un mot interrogatif, celui-ci compte comme première position : \all{Wann arbeitet die Frau?}
\item \textbf{Subordonnée} (après \all{weil, dass, obwohl, wenn, als}...) : le verbe conjugué est à la \cle{dernière position} : \all{..., weil sie weniger verdienen.} (« ... parce qu'elles gagnent moins. »)
\end{enumerate}
\end{propriete}

\begin{attention}[L'erreur n° 1 du francophone : l'inversion oubliée]
En français, on dit : « Aujourd'hui, \textbf{la femme travaille} » (complément + sujet + verbe). En allemand, c'est impossible : après un complément en tête (\all{heute, gestern, dann, deshalb, in Kamerun, leider}...), le verbe doit \emph{immédiatement} suivre, puis le sujet.

\all{Heute arbeiten viele Frauen.} — « Aujourd'hui, beaucoup de femmes travaillent. »

Ne dites jamais : \emph{Heute viele Frauen arbeiten} (faute grave). Retenez le schéma : \textbf{Complément + VERBE + sujet}. Autres exemples : \all{Deshalb kämpfen viele Frauen für ihre Rechte.} (« C'est pourquoi beaucoup de femmes luttent pour leurs droits. ») ; \all{In Yaoundé arbeitet meine Schwester als Lehrerin.} (« À Yaoundé, ma sœur travaille comme enseignante. »)
\end{attention}

\begin{popfigure}
\begin{tikzpicture}[node distance=0.35cm, every node/.style={font=\small}]
\tikzset{
  case/.style={draw=popblue, fill=popblueL, rounded corners=2pt, minimum height=1.1cm, align=center, minimum width=3.2cm},
  caseV/.style={case, fill=poporangeL, draw=poporange, minimum width=2.6cm},
  caseM/.style={case, minimum width=3.4cm},
  caseE/.style={case, fill=popgreenL, draw=popgreen, minimum width=3.2cm}
}
\node[case, align=center] (c1){\textbf{Position 1}\\ Sujet \emph{ou} complément\\ \all{Die Frau} / \all{Heute}};
\node[caseV, right=of c1, align=center] (c2){\textbf{Position 2}\\ VERBE conjugué\\ \all{arbeitet} / \all{hat}};
\node[caseM, right=of c2, align=center] (c3){\textbf{Positions centrales}\\ sujet, compléments\\ \all{die Frau heute}};
\node[caseE, right=of c3, align=center] (c4){\textbf{Dernière position}\\ infinitif / participe II\\ verbe de subordonnée\\ \all{lernen} / \all{gelernt} / \all{verdienen}};
\draw[-{Stealth}, thick, popdark] (c1) -- (c2);
\draw[-{Stealth}, thick, popdark] (c2) -- (c3);
\draw[-{Stealth}, thick, popdark] (c3) -- (c4);
\end{tikzpicture}
\legende{Le « cadre » de la phrase allemande : le verbe conjugué en position 2, le second élément verbal (infinitif, participe II) ou le verbe de la subordonnée tout à la fin.}
\end{popfigure}

\courssub{Phrases avec deux verbes : le cadre verbal}

Quand la phrase contient deux verbes (auxiliaire + participe, modal + infinitif), l'allemand construit un \cle{cadre} (\all{Satzklammer}) : le verbe conjugué ouvre la phrase en position 2, et l'autre élément verbal la ferme à la toute fin. Tout le reste se loge \emph{à l'intérieur} du cadre.

\begin{exemplebox}[Le cadre verbal]
\begin{itemize}
\item Futur / modalité : \all{Sie will einen Beruf lernen.} — « Elle veut apprendre un métier. » (infinitif \all{lernen} à la fin)
\item Passé composé : \all{Sie hat einen Beruf gelernt.} — « Elle a appris un métier. » (participe II \all{gelernt} à la fin)
\item Avec compléments à l'intérieur du cadre : \all{Viele Frauen in Douala haben früher keine Schule besucht.} — « Beaucoup de femmes à Douala n'ont pas fréquenté l'école autrefois. »
\item Subordonnée avec deux verbes : le verbe conjugué reprend la dernière place, juste après le participe/infinitif : \all{..., weil sie einen Beruf lernen will.} — « ... parce qu'elle veut apprendre un métier. » ; \all{..., obwohl sie hart gearbeitet hat.} — « ... bien qu'elle ait travaillé dur. »
\end{itemize}
\end{exemplebox}

\begin{attention}[Ne jamais coller les deux verbes]
Le francophone a tendance à placer l'infinitif juste après le modal, comme en français (« elle veut apprendre »). En allemand, \all{Sie will lernen einen Beruf} est fautif : l'infinitif va \emph{à la fin}. De même au passé composé : \all{Sie hat gelernt einen Beruf} est fautif. Le cadre est sacré : \all{Sie will einen Beruf lernen.} / \all{Sie hat einen Beruf gelernt.}
\end{attention}

\courssub{Les quatre cas : la grammaire des fonctions}

Le français indique la fonction d'un nom par sa \emph{place} (« La mère aide l'enfant » / « L'enfant aide la mère » : l'ordre change le sens). L'allemand, lui, indique la fonction par la \emph{forme} de l'article et du nom : ce sont les \cle{cas}. Grâce aux cas, \all{Der Vater liebt das Kind} et \all{Das Kind liebt der Vater} signifient tous deux « Le père aime l'enfant » : c'est \all{der} (Nominativ) qui désigne le sujet, pas la position.

\begin{definition}[Les quatre cas et leurs fonctions]
\begin{itemize}
\item \cle{Nominativ} : le \textbf{sujet} (qui fait l'action ?). \all{Die Mutter kocht.} — « La mère cuisine. »
\item \cle{Akkusativ} : le \textbf{COD} (qui/quoi ?) et le \textbf{mouvement} vers un lieu. \all{Sie besucht ihre Mutter.} — « Elle rend visite à sa mère. »
\item \cle{Dativ} : le \textbf{COI} (à qui ?) et le \textbf{lieu} (où ?). \all{Er hilft dem Kind.} — « Il aide l'enfant. »
\item \cle{Genitiv} : la \textbf{possession} (de qui ?). \all{Das ist das Haus meiner Eltern.} — « C'est la maison de mes parents. »
\end{itemize}
\end{definition}

\begin{center}
\begin{tabular}{|l|l|l|l|l|}
\hline
\rowcolor{popblueL} \textbf{Cas} & \textbf{Masculin} & \textbf{Féminin} & \textbf{Neutre} & \textbf{Pluriel} \\
\hline
Nominativ & der / ein & die / eine & das / ein & die / -- \\
\hline
Akkusativ & den / einen & die / eine & das / ein & die / -- \\
\hline
Dativ & dem / einem & der / einer & dem / einem & den (\textit{+ n au nom}) / -- \\
\hline
Genitiv & des (\textit{+ s/es au nom}) / eines & der / einer & des (\textit{+ s/es au nom}) / eines & der / -- \\
\hline
\end{tabular}
\end{center}

Points clés du tableau : seul le masculin change à l'Akkusativ (\all{der} $\rightarrow$ \all{den}) ; le Dativ pluriel ajoute un \all{-n} au nom quand c'est possible (\all{den Kindern}) ; le Genitiv masculin et neutre ajoute \all{-(e)s} au nom (\all{des Vaters}, \all{des Kindes}).

\begin{exemplebox}[Les cas en contexte]
\begin{itemize}
\item \all{Die Mutter hilft dem Kind.} — « La mère aide l'enfant. » (\all{dem Kind} : Dativ, car \all{helfen} exige le Dativ)
\item \all{Sie kocht für die Familie.} — « Elle cuisine pour la famille. » (\all{für} + Akkusativ)
\item \all{Das ist das Haus meiner Eltern.} — « C'est la maison de mes parents. » (Genitiv de possession)
\item \all{Der Mann dankt seiner Frau.} — « L'homme remercie sa femme. » (\all{danken} + Dativ)
\item \all{Die Tochter besucht den Vater im Krankenhaus.} — « La fille rend visite au père à l'hôpital. » (\all{den Vater} : Akkusativ, COD)
\end{itemize}
\end{exemplebox}

\begin{attention}[Verbes à Dativ : helfen, danken, gefallen, gehören, antworten, fehlen, passen]
Certains verbes allemands se construisent avec le Dativ alors que leur équivalent français prend un COD. C'est un piège classique : on dit \all{Ich helfe meiner Mutter.} (« J'aide ma mère ») et jamais \emph{meine Mutter} après \all{helfen}. De même : \all{Ich danke meinem Vater.} (« Je remercie mon père ») ; \all{Das Kleid gefällt meiner Schwester.} (« La robe plaît à ma sœur ») ; \all{Das Haus gehört meiner Tante.} (« La maison appartient à ma tante ») ; \all{Sie antwortet dem Lehrer.} (« Elle répond au professeur ») ; \all{Mir fehlt ein Wort.} (« Il me manque un mot / Je ne trouve pas un mot ») ; \all{Die Jacke passt ihm gut.} (« La veste lui va bien »).
\end{attention}

\courssub{Prépositions et cas : fixes ou à deux cas}

\begin{propriete}[Prépositions à cas fixe]
\begin{itemize}
\item \textbf{+ Akkusativ} : \all{für} (pour), \all{ohne} (sans), \all{gegen} (contre), \all{durch} (à travers), \all{um} (autour de). Ex. : \all{Sie kämpft für ihre Rechte.} — « Elle lutte pour ses droits. »
\item \textbf{+ Dativ} : \all{mit} (avec), \all{nach} (après, vers une ville/pays sans article), \all{bei} (chez), \all{zu} (vers une personne), \all{von} (de, provenance), \all{aus} (de, origine), \all{seit} (depuis). Ex. : \all{Sie arbeitet mit ihrem Mann.} — « Elle travaille avec son mari. » ; \all{Nach der Arbeit kocht sie.} — « Après le travail, elle cuisine. » ; \all{Sie kommt aus Kamerun.} — « Elle vient du Cameroun. »
\item \textbf{+ Genitiv} : \all{trotz} (malgré), \all{während} (pendant), \all{wegen} (à cause de), \all{statt} (au lieu de). Ex. : \all{Trotz der Schwierigkeiten studiert sie.} — « Malgré les difficultés, elle fait des études. »
\end{itemize}
\end{propriete}

\begin{propriete}[Prépositions à deux cas (\all{Wechselpräpositionen}) : in, an, auf, über, unter, vor, hinter, neben, zwischen]
La question décide du cas :
\begin{itemize}
\item \textbf{Wo ?} (où ? = lieu, position statique) $\rightarrow$ \textbf{Dativ} : \all{Die Frau arbeitet in der Schule.} — « La femme travaille dans l'école (elle y est). »
\item \textbf{Wohin ?} (vers où ? = mouvement, direction) $\rightarrow$ \textbf{Akkusativ} : \all{Die Frau geht in die Schule.} — « La femme va à l'école (elle s'y rend). »
\end{itemize}
Comparez : \all{Das Kind sitzt auf dem Stuhl.} (« L'enfant est assis sur la chaise ») / \all{Das Kind setzt sich auf den Stuhl.} (« L'enfant s'assoit sur la chaise »).
\end{propriete}

\courssub{Comparaison français / allemand : deux logiques différentes}

\begin{center}
\begin{tabularx}{\linewidth}{|X|X|X|}
\hline
\rowcolor{popblueL} \textbf{Français} & \textbf{Deutsch} & \textbf{Remarque} \\
\hline
Aujourd'hui, beaucoup de femmes travaillent. & Heute arbeiten viele Frauen. & Inversion obligatoire après le complément en tête. \\
\hline
Je sais que la femme travaille aujourd'hui. & Ich weiß, dass die Frau heute arbeitet. & Verbe à la fin dans la subordonnée. \\
\hline
Elle veut apprendre un métier. & Sie will einen Beruf lernen. & Infinitif à la fin (cadre verbal). \\
\hline
La mère aide l'enfant. & Die Mutter hilft dem Kind. & Le sens vient de l'ordre en français, du cas (Dativ) en allemand. \\
\hline
Elle va à l'école. / Elle est à l'école. & Sie geht in die Schule. / Sie ist in der Schule. & Mouvement = Akkusativ ; lieu = Dativ. \\
\hline
\end{tabularx}
\end{center}

\begin{aretenir}[Les cinq réflexes du bon germaniste]
\begin{enumerate}
\item Phrase déclarative : verbe conjugué en 2e position, toujours.
\item Complément en tête $\rightarrow$ inversion : \all{Heute arbeitet sie.}
\item Subordonnée (\all{weil, dass, obwohl, wenn}) $\rightarrow$ verbe à la fin.
\item Deux verbes $\rightarrow$ cadre : \all{Sie hat einen Beruf gelernt.}
\item Fonction du nom $\rightarrow$ cas : sujet = Nominativ, COD = Akkusativ, COI et \all{helfen/danken} = Dativ, possession = Genitiv.
\end{enumerate}
\end{aretenir}

\begin{exoresolu}[Remettre la phrase en ordre et justifier]
Énoncé : Remettez les mots suivants dans l'ordre correct, puis traduisez : \all{verdienen / weil / weniger / viele Frauen / als Männer / , / Meine Tante sagt}
\tcblower
Solution détaillée :

\textbf{1. Analyser les éléments.}
Nous avons deux verbes conjugués : \all{sagt} (dire) et \all{verdienen} (gagner). Il y aura donc deux propositions.
La conjonction \all{weil} (parce que) introduit une \textbf{subordonnée} : le verbe \all{verdienen} doit aller \textbf{à la fin} de cette subordonnée.
Les éléments de la subordonnée : \all{weil} + sujet \all{viele Frauen} + comparatif \all{weniger als Männer} + verbe \all{verdienen}.
$\rightarrow$ \all{weil viele Frauen weniger als Männer verdienen}

Les éléments restants forment la \textbf{principale} : \all{Meine Tante sagt}.

\textbf{2. Assembler.}
En allemand, une subordonnée peut occuper la \textbf{position 1} de la principale (Vorfeld). C'est la structure la plus naturelle ici pour lier la cause (\all{weil...}) et la conséquence/réaction (\all{Meine Tante sagt}).
Structure : [Subordonnée] , [Verbe principal] [Sujet principal] .
$\rightarrow$ \all{Weil viele Frauen weniger als Männer verdienen, sagt meine Tante.}

\textbf{3. Vérification des règles.}
\begin{itemize}
\item Subordonnée \all{weil} : verbe \all{verdienen} en dernière position. \checkmark
\item Subordonnée en position 1 $\rightarrow$ verbe principal \all{sagt} en position 2, sujet \all{meine Tante} après (inversion). \checkmark
\item Comparatif : \all{weniger als} (moins que). \checkmark
\end{itemize}

\textbf{Traduction :} « Parce que beaucoup de femmes gagnent moins que les hommes, ma tante le dit / ma tante l'affirme. »
\end{exoresolu}

\begin{exercice}[À vous de jouer]
\begin{enumerate}
\item Traduisez en allemand : « Aujourd'hui, beaucoup de femmes travaillent comme médecins. »
\item Traduisez : « Je sais que ma sœur aide souvent ma mère. » (attention à \all{helfen} !)
\item Remettez en ordre : \all{hat / einen Beruf / meine Cousine / in Yaoundé / gelernt}
\item Complétez avec le bon cas (article, adjectif, nom) : \\
\all{Sie kämpft für (ihre Rechte).} \\
\all{Er wohnt bei (seine Eltern).} \\
\all{Trotz (die Arbeit) hilft sie (ihre Familie).}
\end{enumerate}
\end{exercice}

\begin{corrige}[Exercice « À vous de jouer »]
\begin{enumerate}
\item \all{Heute arbeiten viele Frauen als Ärztinnen.} (Inversion après \all{Heute} ; « comme » = \all{als} + Nominativ ; \all{Ärztin} pluriel \all{Ärztinnen}).
\item \all{Ich weiß, dass meine Schwester meiner Mutter oft hilft.} (Subordonnée \all{dass} : verbe \all{hilft} à la fin ; \all{helfen} + Dativ : \all{meiner Mutter} ; place de \all{oft} avant le verbe à la fin).
\item \all{Meine Cousine hat in Yaoundé einen Beruf gelernt.} (Cadre : \all{hat} en 2e position, \all{gelernt} à la fin ; le complément de lieu \all{in Yaoundé} et le COD \all{einen Beruf} à l'intérieur du cadre).
\item \all{Sie kämpft für ihre Rechte.} (\all{für} + Akkusativ pluriel) ; \all{Er wohnt bei seinen Eltern.} (\all{bei} + Dativ pluriel : article \all{seinen} + \all{-n} à \all{Eltern}) ; \all{Trotz der Arbeit hilft sie ihrer Familie.} (\all{trotz} + Genitiv féminin \all{der Arbeit} ; inversion \all{hilft sie} car complément en tête ; \all{helfen} + Dativ féminin \all{ihrer Familie}).
\end{enumerate}
\end{corrige}

Ces mécanismes — place du verbe et choix des cas — seront vos outils principaux dans les sections suivantes : les temps du passé d'abord, puis la rédaction de lettres, d'essais et de dialogues, et enfin la traduction, où chaque phrase devra passer le contrôle des cinq réflexes ci-dessus.

\coursec{Grammaire 2 : les temps du passé}

Après avoir révisé l'ordre des mots et les cas (Grammaire 1), nous abordons maintenant les temps du passé, indispensables pour raconter, décrire l'évolution de la situation de la femme ou traduire des textes narratifs. En allemand, le choix du temps dépend du type de texte (oral ou écrit) et de la fonction du verbe (auxiliaire, modal, verbe d'action).

\courssub{Observation : un texte mêlant passé oral et passé écrit}

\begin{textebox}[Évolution de la situation de la femme en Allemagne]
In der Vergangenheit blieb die Frau meist zu Hause. Sie kümmerte sich um die Kinder und den Haushalt. Der Mann ging arbeiten und verdiente das Geld. Das war die traditionelle Rollenverteilung.

Heute hat sich vieles geändert. Viele Frauen haben studiert und arbeiten jetzt in Berufen, die früher nur Männern vorbehalten waren. Sie sind Ärztinnen, Ingenieurinnen oder Politikerinnen geworden. Dennoch gibt es noch Herausforderungen: Frauen verdienen oft weniger als Männer für die gleiche Arbeit. Nachdem die Bundesregierung das Entgelttransparenzgesetz verabschiedet hatte, wurde die Lage etwas besser, aber der Weg zur echten Gleichberechtigung ist noch lang.
\end{textebox}

\noindent\textbf{Glossaire :} \all{die Vergangenheit} (le passé) – \all{bleiben, blieb, ist geblieben} (rester) – \all{sich kümmern um + Akk.} (s'occuper de) – \all{der Haushalt, -e} (le ménage) – \all{die Rollenverteilung, -en} (la répartition des rôles) – \all{sich ändern} (changer) – \all{studieren} (faire des études supérieures) – \all{vorbehalten sein + Dat.} (être réservé à) – \all{die Ärztin, -nen} (la femme médecin) – \all{die Ingenieurin, -nen} (la femme ingénieur) – \all{die Politikerin, -nen} (la femme politique) – \all{die Herausforderung, -en} (le défi) – \all{verdienen} (gagner [de l'argent]) – \all{das Entgelttransparenzgesetz, -e} (la loi sur la transparence des rémunérations) – \all{verabschieden} (adopter [une loi]) – \all{die Gleichberechtigung} (l'égalité des droits).

\begin{exoresolu}[Repérage des temps du passé]
\textbf{Énoncé :} Relevez dans le texte ci-dessus tous les verbes conjugués au passé. Pour chacun, indiquez : l'infinitif, le temps (Perfekt, Präteritum, Plusquamperfekt), l'auxiliaire (haben/sein) et la valeur (récit écrit, conversation, antériorité).

\tcblower
\textbf{Solution détaillée :}

\begin{center}
\begin{tabularx}{\linewidth}{|X|X|l|l|X|}
\hline
\rowcolor{popblueL}
Forme trouvée & Infinitif & Temps & Auxiliaire & Valeur / Contexte \\ \hline
blieb & bleiben & Präteritum & — & Récit écrit (narration) \\ \hline
kümmerte & sich kümmern & Präteritum & — & Récit écrit \\ \hline
ging & gehen & Präteritum & — & Récit écrit \\ \hline
verdiente & verdienen & Präteritum & — & Récit écrit \\ \hline
war & sein & Präteritum & — & Récit écrit (sein = Präteritum obligatoire) \\ \hline
hat ... geändert & sich ändern & Perfekt & haben & Constat actuel / bilan \\ \hline
haben ... studiert & studieren & Perfekt & haben & Expérience vécue, bilan \\ \hline
sind ... geworden & werden & Perfekt & sein & Changement d'état (devenir) \\ \hline
verabschiedet hatte & verabschieden & Plusquamperfekt & haben & Antériorité (après que le gouvernement eut adopté) \\ \hline
wurde & werden & Präteritum & — & Récit écrit (passif : \all{wurde ... besser}) \\ \hline
\end{tabularx}
\end{center}

\emph{Observation :} Le texte alterne Präteritum (récit du passé : \all{blieb, kümmerte, ging, verdiente, war}) et Perfekt (bilan actuel, expérience : \all{hat geändert, haben studiert, sind geworden}). Le Plusquamperfekt (\all{verabschiedet hatte}) marque l'antériorité par rapport au Präteritum (\all{wurde}). Les formes \all{gibt, verdienen, ist} sont au présent : elles décrivent la situation actuelle.
\end{exoresolu}

\courssub{Le Perfekt : le passé de la conversation et de la lettre}

\begin{propriete}[Formation du Perfekt]
Le Perfekt se forme avec l'auxiliaire \textbf{haben} ou \textbf{sein} au présent + le \textbf{Partizip II} (participe passé) placé \textbf{à la fin de la phrase}.

\begin{center}
\begin{tabular}{|l|l|l|l|}
\hline
\rowcolor{popblueL}
Sujet & Auxiliaire (Präsens) & ... & Partizip II \\ \hline
ich & habe / bin & ... & gemacht / gegangen \\ \hline
du & hast / bist & ... & gemacht / gegangen \\ \hline
er/sie/es & hat / ist & ... & gemacht / gegangen \\ \hline
wir & haben / sind & ... & gemacht / gegangen \\ \hline
ihr & habt / seid & ... & gemacht / gegangen \\ \hline
sie/Sie & haben / sind & ... & gemacht / gegangen \\ \hline
\end{tabular}
\end{center}

\textbf{Choix de l'auxiliaire :}
\begin{itemize}
\item \textbf{haben} : la majorité des verbes (transitifs, réflexifs, verbes d'état, verbes de durée, verbes modaux).
\item \textbf{sein} : verbes de \textbf{mouvement} (gehen, kommen, laufen, fahren, fliegen, reisen...) et verbes de \textbf{changement d'état} (aufstehen, einschlafen, aufwachen, sterben, werden), ainsi que \textbf{sein} et \textbf{bleiben} (exception à mémoriser).
\end{itemize}
\end{propriete}

\begin{exemplebox}[Perfekt avec haben / sein]
\all{Gestern habe ich meiner Mutter im Haushalt geholfen.} « Hier, j'ai aidé ma mère dans le ménage. » (haben – verbe transitif ; \all{helfen} + datif) \\
\all{Meine Schwester ist nach Berlin gefahren.} « Ma sœur est allée à Berlin. » (sein – mouvement) \\
\all{Die Frau ist Ärztin geworden.} « La femme est devenue médecin. » (sein – changement d'état) \\
\all{Wir haben lange über die Gleichberechtigung diskutiert.} « Nous avons longtemps discuté de l'égalité des droits. » (haben – verbe de communication) \\
\all{Früher sind die Frauen zu Hause geblieben.} « Autrefois, les femmes restaient à la maison. » (sein – \all{bleiben} prend toujours \all{sein})
\end{exemplebox}

\begin{attention}[Sein, haben, modaux : pas de Perfekt dans le récit !]
En narration écrite (article, récit, essai), \textbf{on n'utilise jamais le Perfekt avec sein, haben et les verbes modaux}. On leur préfère le \textbf{Präteritum}, même à l'oral dans un contexte formel.

\begin{center}
\begin{tabularx}{\linewidth}{|X|X|}
\hline
\rowcolor{poporangeL}
\bfseries Incorrect (Perfekt artificiel) & \bfseries Correct (Präteritum) \\ \hline
Sie \textbf{ist müde gewesen}. & Sie \textbf{war} müde. \\ \hline
Er \textbf{hat keine Zeit gehabt}. & Er \textbf{hatte} keine Zeit. \\ \hline
Sie \textbf{hat arbeiten gewollt}. & Sie \textbf{wollte} arbeiten. \\ \hline
Wir \textbf{haben nicht kommen gekonnt}. & Wir \textbf{konnten} nicht kommen. \\ \hline
\end{tabularx}
\end{center}

\emph{Règle d'or :} dans un texte narratif, écrivez toujours \all{war, hatte, konnte, musste, wollte, sollte, durfte, mochte} et jamais \all{ist gewesen, hat gehabt, hat gekonnt...}.
\end{attention}

\courssub{Le Präteritum : le passé du récit écrit}

\begin{propriete}[Formation du Präteritum]
Le Präteritum (ou Imperfekt) est le temps du \textbf{récit écrit} (journal, roman, compte rendu, essai historique). Il se forme sans auxiliaire.

\begin{itemize}
\item \textbf{Verbes faibles (réguliers) :} radical + \textbf{-te} + terminaisons personnelles. \\ \all{machen → ich machte, du machtest, er machte, wir machten, ihr machtet, sie machten}
\item \textbf{Verbes forts (irréguliers) :} changement de voyelle radicale + terminaisons (sans -e aux 1\textsuperscript{re} et 3\textsuperscript{e} pers. sg). \\ \all{gehen → ich ging, du gingst, er ging, wir gingen, ihr gingt, sie gingen} \\
\all{sehen → ich sah, du sahst, er sah, wir sahen, ihr saht, sie sahen} \\
\all{geben → ich gab, du gabst, er gab, wir gaben, ihr gabt, sie gaben}
\item \textbf{Verbes mixtes :} changement de voyelle + -te. \\ \all{bringen → ich brachte, du brachtest, er brachte...} ; \all{denken → ich dachte...} ; \all{kennen → ich kannte...}
\end{itemize}

\textbf{Verbes indispensables (irréguliers) :}
\begin{center}
\begin{tabular}{|l|l|l|l|}
\hline
\rowcolor{popblueL}
Infinitif & Präteritum (ich/er/sie/es) & Partizip II & Français \\ \hline
sein & war & gewesen & être \\ \hline
haben & hatte & gehabt & avoir \\ \hline
werden & wurde & geworden & devenir / aux. passif \\ \hline
können & konnte & gekonnt & pouvoir \\ \hline
müssen & musste & gemusst & devoir \\ \hline
wollen & wollte & gewollt & vouloir \\ \hline
sollen & sollte & gesollt & devoir (conseil) \\ \hline
dürfen & durfte & gedurft & avoir le droit \\ \hline
mögen & mochte & gemocht & aimer \\ \hline
gehen & ging & gegangen & aller \\ \hline
kommen & kam & gekommen & venir \\ \hline
sehen & sah & gesehen & voir \\ \hline
geben & gab & gegeben & donner \\ \hline
nehmen & nahm & genommen & prendre \\ \hline
bleiben & blieb & geblieben & rester \\ \hline
\end{tabular}
\end{center}
\end{propriete}

\begin{exemplebox}[Präteritum dans un récit sur la condition féminine]
\all{Früher blieb die Frau zu Hause.} « Autrefois, la femme restait à la maison. » \\
\all{Sie kümmerte sich um die Kinder.} « Elle s'occupait des enfants. » \\
\all{Der Mann ging arbeiten.} « L'homme allait travailler. » \\
\all{Er verdiente das Geld.} « Il gagnait l'argent. » \\
\all{Das war die traditionelle Rollenverteilung.} « C'était la répartition traditionnelle des rôles. » \\
\all{Die Frauen hatten keine Wahl.} « Les femmes n'avaient pas le choix. » \\
\all{Sie durften nicht wählen.} « Elles n'avaient pas le droit de voter. » \\
\all{Nach 1918 erhielten sie das Wahlrecht.} « Après 1918, elles obtinrent le droit de vote. »
\end{exemplebox}

\courssub{Le Plusquamperfekt : l'antériorité dans le passé}

\begin{propriete}[Formation et emploi du Plusquamperfekt]
Le Plusquamperfekt exprime une action \textbf{antérieure} à une autre action déjà passée (repère au Präteritum ou au Perfekt). Il se forme comme le Perfekt, mais l'auxiliaire est au \textbf{Präteritum} : \textbf{hatte / war + Partizip II}.

\begin{center}
\begin{tabular}{|l|l|l|}
\hline
\rowcolor{popblueL}
Sujet & Auxiliaire (Präteritum) & Partizip II \\ \hline
ich & hatte / war & gemacht / gegangen \\ \hline
du & hattest / warst & gemacht / gegangen \\ \hline
er/sie/es & hatte / war & gemacht / gegangen \\ \hline
wir & hatten / waren & gemacht / gegangen \\ \hline
ihr & hattet / wart & gemacht / gegangen \\ \hline
sie/Sie & hatten / waren & gemacht / gegangen \\ \hline
\end{tabular}
\end{center}

\textbf{Emploi typique :} subordonnées temporelles avec \all{nachdem} (après que), \all{bevor} (avant que), \all{als} (lorsque), \all{bis} (jusqu'à ce que). Attention : dans la subordonnée, le verbe conjugué (\all{hatte/war}) va \textbf{à la fin} et le Partizip II juste avant lui.
\end{propriete}

\begin{exemplebox}[Plusquamperfekt – antériorité]
\all{Nachdem sie die Schule \textbf{beendet hatte}, machte sie eine Ausbildung.} « Après qu'elle \textbf{eut terminé} l'école, elle fit une formation. » \\
\all{Bevor sie heirateten, \textbf{hatten sie lange zusammen gewohnt}.} « Avant de se marier, ils \textbf{avaient longtemps habité} ensemble. » \\
\all{Als ich ankam, \textbf{war sie schon gegangen}.} « Quand je suis arrivé, elle \textbf{était déjà partie}. » \\
\all{Die Regierung \textbf{hatte das Gesetz verabschiedet}, bevor die Proteste begannen.} « Le gouvernement \textbf{avait adopté} la loi avant que les manifestations ne commencent. »
\end{exemplebox}

\courssub{Tableau comparatif : correspondance français / allemand}

\begin{propriete}[Correspondance des temps du passé FR ↔ DE]
\begin{center}
\begin{tabularx}{\linewidth}{|X|X|X|X|}
\hline
\rowcolor{popblueL}
\bfseries Français & \bfseries Allemand & \bfseries Exemple & \bfseries Remarque \\ \hline
Passé composé & \textbf{Perfekt} & \all{Ich habe gearbeitet.} « J'ai travaillé. » & Oral, lettres, bilan, expérience vécue. \\ \hline
Imparfait & \textbf{Präteritum} & \all{Ich arbeitete. / Ich war müde.} « Je travaillais / J'étais fatigué. » & Récit écrit, description passée ; \textbf{obligatoire pour sein/haben/modaux}. \\ \hline
Passé simple & \textbf{Präteritum} & \all{Sie ging nach Hause.} « Elle rentra chez elle. » & Narration écrite (journal, roman, histoire). \\ \hline
Plus-que-parfait & \textbf{Plusquamperfekt} & \all{Sie hatte gearbeitet.} « Elle avait travaillé. » & Antériorité par rapport à un passé. \\ \hline
Passé récent (venir de) & \textbf{Perfekt} + \all{gerade} & \all{Ich habe gerade gegessen.} « Je viens de manger. » & \all{gerade} + Perfekt = « venir de ». \\ \hline
\end{tabularx}
\end{center}

\emph{Point clé :} le français utilise l'imparfait pour la description \textbf{et} l'habitude passée ; l'allemand utilise le Präteritum pour les deux, mais le Perfekt peut remplacer le Präteritum à l'oral pour les verbes d'action (sauf sein/haben/modaux).
\end{propriete}

\courssub{Frise chronologique : Plusquamperfekt → Präteritum/Perfekt → Präsens}

\begin{popfigure}
\begin{tikzpicture}[
  timeline/.style={very thick, popblue},
  event/.style={circle, fill=popblue, inner sep=2pt},
  labelnode/.style={align=center, font=\small, text width=3.8cm},
  arrow/.style={-Stealth, thick, popdark}
]
% Ligne de temps
\draw[timeline] (0,0) -- (14,0);

% Points temporels
\node[event] (pp) at (1,0) {};
\node[event] (p) at (7,0) {};
\node[event] (pr) at (13,0) {};

% Étiquettes temps
\node[labelnode, above=0.5cm of pp, align=center]{\textbf{Plusquamperfekt}\\(antériorité)\\\all{hatte gemacht}\\« avait fait »};
\node[labelnode, above=0.5cm of p, align=center]{\textbf{Präteritum / Perfekt}\\(passé)\\\all{machte / hat gemacht}\\« fit / a fait »};
\node[labelnode, above=0.5cm of pr, align=center]{\textbf{Präsens}\\(présent)\\\all{macht}\\« fait »};

% Flèches
\draw[arrow] (pp) -- (p) node[midway, below=0.3cm, font=\small] {antérieur à};
\draw[arrow] (p) -- (pr) node[midway, below=0.3cm, font=\small] {antérieur à};

% Exemples contextuels
\node[labelnode, below=1.2cm of pp, text width=4.2cm, text=popdark] {
\textbf{Exemple :}\\
\all{Nachdem sie \textbf{studiert hatte},}\\
\all{ging sie arbeiten.}\\
« Après qu'elle \textbf{eut étudié}, elle alla travailler. »};

\node[labelnode, below=1.2cm of p, text width=4.2cm, text=popdark] {
\textbf{Exemple :}\\
\all{Früher \textbf{blieb} sie zu Hause.}\\
\all{Heute \textbf{hat sie} einen Job.}\\
« Autrefois elle \textbf{restait} à la maison. Aujourd'hui elle \textbf{a} un emploi. »};

\node[labelnode, below=1.2cm of pr, text width=4.2cm, text=popdark] {
\textbf{Exemple :}\\
\all{Heute \textbf{arbeitet} sie}\\
\all{als Ingenieurin.}\\
« Aujourd'hui elle \textbf{travaille} comme ingénieure. »};

% Marqueurs temporels
\node[font=\small, text=popmuted] at (1,-3.2) {nachdem / bevor / als};
\node[font=\small, text=popmuted] at (7,-3.2) {früher / gestern / damals};
\node[font=\small, text=popmuted] at (13,-3.2) {heute / jetzt / derzeit};
\end{tikzpicture}
\legende{Frise chronologique des temps du passé : l'antériorité (Plusquamperfekt) précède le passé narratif (Präteritum) ou conversationnel (Perfekt), qui précède le présent.}
\end{popfigure}

\courssub{Les 20 verbes forts les plus fréquents : tableau de référence}

\begin{aretenir}[Verbes forts à mémoriser impérativement]
\begin{center}
\begin{tabular}{|l|l|l|l|}
\hline
\rowcolor{popblueL}
Infinitiv & Präteritum & Partizip II & Français \\ \hline
sein & war & ist gewesen & être \\ \hline
haben & hatte & hat gehabt & avoir \\ \hline
werden & wurde & ist geworden & devenir \\ \hline
gehen & ging & ist gegangen & aller \\ \hline
kommen & kam & ist gekommen & venir \\ \hline
sehen & sah & hat gesehen & voir \\ \hline
geben & gab & hat gegeben & donner \\ \hline
nehmen & nahm & hat genommen & prendre \\ \hline
bleiben & blieb & ist geblieben & rester \\ \hline
finden & fand & hat gefunden & trouver \\ \hline
halten & hielt & hat gehalten & tenir \\ \hline
lassen & ließ & hat gelassen & laisser \\ \hline
sprechen & sprach & hat gesprochen & parler \\ \hline
treffen & traf & hat getroffen & rencontrer \\ \hline
trinken & trank & hat getrunken & boire \\ \hline
schreiben & schrieb & hat geschrieben & écrire \\ \hline
lesen & las & hat gelesen & lire \\ \hline
essen & aß & hat gegessen & manger \\ \hline
fahren & fuhr & ist gefahren & aller (en véhicule) \\ \hline
stehen & stand & hat gestanden & être debout \\ \hline
\end{tabular}
\end{center}
Liste non exhaustive, à mémoriser impérativement pour le Präteritum et le Perfekt. Notez bien l'auxiliaire (\all{ist} / \all{hat}) indiqué devant chaque Partizip II.
\end{aretenir}

\courssub{Exercice résolu : transformation et choix du temps}

\begin{exoresolu}[Du français vers l'allemand : choisir le bon temps]
\textbf{Énoncé :} Traduisez les phrases suivantes en allemand en justifiant le choix du temps (Perfekt, Präteritum, Plusquamperfekt).

\begin{enumerate}
\item Autrefois, les femmes restaient à la maison. (contexte : récit historique)
\item Hier, j'ai aidé ma mère dans le ménage. (contexte : conversation)
\item Après qu'elle eut fini ses études, elle trouva du travail. (antériorité + récit)
\item Elle est devenue médecin l'année dernière. (changement d'état, conversation)
\item Nous n'avions pas le droit de voter. (récit, modal)
\item Avant de se marier, ils avaient longtemps vécu ensemble. (antériorité)
\end{enumerate}

\tcblower
\textbf{Solution détaillée et justifications :}

\begin{enumerate}
\item \all{Früher blieben die Frauen zu Hause.} \\
\textbf{Präteritum} (\all{blieben}) – récit historique / description passée. Verbe fort \all{bleiben} → \all{blieb} (sg) / \all{blieben} (pl).
\item \all{Gestern habe ich meiner Mutter im Haushalt geholfen.} \\
\textbf{Perfekt} (\all{habe ... geholfen}) – conversation / événement ponctuel passé. \all{helfen} + datif (\all{meiner Mutter}) → Partizip II \all{geholfen}, auxiliaire \all{haben}.
\item \all{Nachdem sie ihr Studium \textbf{beendet hatte}, fand sie eine Arbeit.} \\
\textbf{Plusquamperfekt} (\all{hatte ... beendet}) dans la subordonnée \all{nachdem} (antériorité) + \textbf{Präteritum} (\all{fand}) dans la principale (récit). \all{beenden} (faible) → \all{beendet} ; \all{finden} (fort) → \all{fand}.
\item \all{Sie ist letztes Jahr Ärztin geworden.} \\
\textbf{Perfekt} avec \textbf{sein} (\all{ist ... geworden}) – changement d'état (\all{werden}), contexte conversationnel. Variante narrative : \all{Sie wurde letztes Jahr Ärztin.} (Präteritum).
\item \all{Wir durften nicht wählen.} \\
\textbf{Präteritum} (\all{durften}) – \textbf{obligatoire pour les modaux} en narration. \all{dürfen} → Präteritum \all{durfte} (sg) / \all{durften} (pl).
\item \all{Bevor sie heirateten, \textbf{hatten sie lange zusammen gewohnt}.} \\
\textbf{Plusquamperfekt} (\all{hatten ... gewohnt}) après \all{bevor} (antériorité). \all{wohnen} (faible) → Partizip II \all{gewohnt}, auxiliaire \all{haben}. Attention à l'ordre des mots : le Partizip II se place à la fin, après \all{zusammen}.
\end{enumerate}
\end{exoresolu}

\begin{attention}[Piège majeur : \textbf{bleiben} et \textbf{sein} prennent \textbf{sein} au Perfekt !]
Beaucoup de francophones écrivent *\all{hat geblieben} ou *\all{hat gewesen} par analogie avec « j'ai resté / j'ai été ». \textbf{Erreur !}

\begin{center}
\begin{tabularx}{\linewidth}{|X|X|}
\hline
\rowcolor{poporangeL}
\bfseries Incorrect & \bfseries Correct \\ \hline
Sie \textbf{hat zu Hause geblieben}. & Sie \textbf{ist zu Hause geblieben}. \\ \hline
Er \textbf{hat krank gewesen}. & Er \textbf{ist krank gewesen}. (Perfekt) / Er \textbf{war krank}. (Präteritum, préférable) \\ \hline
Wir \textbf{haben nach Hause gegangen}. & Wir \textbf{sind nach Hause gegangen}. (\all{gehen} = mouvement → sein) \\ \hline
\end{tabularx}
\end{center}

\emph{Mnémotechnique :} \textbf{sein} et \textbf{bleiben} expriment un \textbf{état} → ils prennent \textbf{sein} (comme les verbes de mouvement). Pensez : \all{Ich bin zu Hause} (état) → \all{Ich bin zu Hause geblieben}.
\end{attention}

\courssub{Synthèse : quel temps choisir pour votre production ?}

\begin{methode}[Choisir le bon temps du passé pour écrire / traduire]
\begin{enumerate}
\item \textbf{Identifiez le type de texte :}
\begin{itemize}
\item Lettre personnelle, dialogue, récit oral, e-mail → \textbf{Perfekt} (sauf sein/haben/modaux).
\item Article, essai, compte rendu, nouvelle, texte historique → \textbf{Präteritum}.
\end{itemize}
\item \textbf{Identifiez le verbe :}
\begin{itemize}
\item \all{sein, haben, können, müssen, wollen, sollen, dürfen, mögen} → \textbf{toujours Präteritum} en narration.
\item Verbe de mouvement (gehen, kommen...) ou de changement d'état (werden, aufwachen...), ainsi que \all{sein} et \all{bleiben} → \textbf{Perfekt avec sein}.
\item Autre verbe d'action → \textbf{Perfekt avec haben} (oral) / \textbf{Präteritum} (écrit).
\end{itemize}
\item \textbf{Y a-t-il antériorité ?} (mots-clés : \all{nachdem, bevor, als, bis}) → \textbf{Plusquamperfekt} dans la subordonnée, verbe conjugué à la fin.
\item \textbf{Relisez-vous :} vérifiez les majuscules (noms), l'ordre des mots (verbe en 2\textsuperscript{e} position, Partizip II à la fin), les cas (accusatif/datif après prépositions), les Umlaute (ä/ö/ü) et le ß.
\end{enumerate}
\end{methode}

\begin{exemplebox}[Application mixte : lettre formelle sur l'égalité]
\textbf{Contexte :} Vous écrivez une lettre au maire pour demander une crèche municipale. Style formel, mais personnel.

\all{Sehr geehrter Herr Bürgermeister,}\\
\all{ich schreibe Ihnen, weil sich die Situation für berufstätige Mütter in unserer Stadt nicht verbessert hat.} (Perfekt – bilan actuel)\\
\all{Als ich vor fünf Jahren hierher zog, gab es nur eine Kita.} (Präteritum – récit passé : \all{zog, gab})\\
\all{Nachdem die alte Kita geschlossen worden war, mussten viele Frauen ihre Jobs aufgeben.} (Plusquamperfekt passif – antériorité : \all{war ... geschlossen worden} ; Präteritum modal : \all{mussten})\\
\all{Heute arbeiten viele Frauen, aber sie finden keine Betreuung für ihre Kinder.} (Präsens – situation actuelle)\\
\all{Ich bitte Sie dringend, eine neue Kita bauen zu lassen.}\\
\all{Mit freundlichen Grüßen}\\
\all{Marie Nkomo}
\end{exemplebox}

\begin{savaistu}[Le genre grammatical et les noms de métiers au féminin]
En allemand, le féminin des noms de métier se forme généralement par \textbf{-in} (au pluriel \textbf{-innen}) : \all{der Arzt → die Ärztin}, \all{der Ingenieur → die Ingenieurin}, \all{der Lehrer → die Lehrerin}. Depuis les années 1980, l'usage recommande souvent la \textbf{forme paire} (\all{Ärztinnen und Ärzte}) ou des formes inclusives comme le \textbf{Genre-Sternchen} (\all{Ärzt*innen}) ou les \textbf{deux-points} (\all{Ärzt:innen}). Dans les textes administratifs récents (Bundestag, universités), vous rencontrerez aussi \all{die Studierenden} (participe présent substantivé) au lieu de \all{Studenten/Studentinnen}. Pour le baccalauréat, maîtrisez la forme classique en \textbf{-in/-innen} et sachez reconnaître les formes inclusives.
\end{savaistu}

\coursec{Structures de production : lettre, essai, dialogue}

Dans la section précédente, nous avons révisé les temps du passé (Perfekt et Präteritum), indispensables pour raconter des événements. Mais raconter ne suffit pas : à l'examen comme dans la vie réelle, on vous demandera de \cle{produire} des textes organisés — une lettre à un ami, un court essai sur la place de la femme dans la société, un dialogue entre deux personnages. Chacun de ces genres obéit à des règles de construction précises. C'est ce que nous allons apprendre maintenant, en réinvestissant l'ordre des mots, les cas et les temps déjà étudiés.

\courssub{La lettre personnelle (der persönliche Brief)}

\textbf{Observation.} Lisez d'abord ce début de lettre écrit par une élève de Yaoundé à son amie allemande :

\begin{textebox}[Extrait de lettre personnelle]
Yaoundé, den 15. März 2025

Liebe Anna,

vielen Dank für deinen Brief. Ich habe mich sehr gefreut! Bei uns in Kamerun gibt es viele Diskussionen über die Rolle der Frau. Meine Mutter arbeitet als Lehrerin, aber sie macht auch fast alle Hausarbeit. Ich finde das ungerecht.

Wie ist das bei euch in Deutschland? Schreib mir bald!

Viele Grüße

Amina
\end{textebox}

Glossaire : \all{sich freuen} : se réjouir ; \all{die Rolle, -n} : le rôle ; \all{die Hausarbeit} : les tâches ménagères ; \all{ungerecht} : injuste ; \all{die Diskussion, -en} : la discussion ; \all{arbeiten als} : travailler comme (+\textbf{Nominatif}).

Observez la disposition : la date en haut à droite, la salutation suivie d'une \cle{virgule}, puis la ligne suivante qui commence par une \cle{minuscule} (\all{vielen Dank...}), et enfin une formule de politesse avant la signature.

\begin{propriete}[Structure de la lettre personnelle allemande]
\begin{enumerate}
\item \textbf{Lieu et date} (en haut à droite) : \all{Yaoundé, den 15. März 2025}. On utilise l'accusatif avec \all{den} + date ordinale : \all{den 3. Mai}, \all{den 22. Dezember 2024}.
\item \textbf{Salutation} : \all{Liebe Anna,} (à une femme) / \all{Lieber Paul,} (à un homme) / \all{Liebe Familie Mbarga,} — toujours suivie d'une virgule.
\item \textbf{Introduction} : remerciement ou nouvelles (\all{vielen Dank für deinen Brief}, \all{wie geht es dir?}).
\item \textbf{Corps} : le sujet développé en un ou deux paragraphes.
\item \textbf{Conclusion} : question ou souhait (\all{Schreib mir bald!}).
\item \textbf{Formule finale} : \all{Viele Grüße} / \all{Herzliche Grüße} / \all{Liebe Grüße}, puis la signature, sans virgule.
\end{enumerate}
\end{propriete}

\begin{attention}[La minuscule après la salutation]
Contrairement au français (« Chère Anna, \textbf{J}e te remercie... »), la première ligne du corps de la lettre allemande commence par une \cle{minuscule} : \all{Liebe Anna,} puis à la ligne \all{vielen Dank für deinen Brief.} C'est une erreur classique et sanctionnée à l'examen. Seule exception : si le premier mot est un nom (majuscule obligatoire) ou le pronom formel \all{Sie}.
\end{attention}

\begin{popfigure}
\begin{tikzpicture}
\node[draw=popblue, fill=popblueL, rounded corners, text width=8.5cm, align=left, font=\small] (brief) at (0,0) {Yaoundé, den 15. März 2025\\[2pt]Liebe Anna,\\[2pt]vielen Dank für deinen Brief. Ich habe mich sehr gefreut...\\[2pt]Viele Grüße\\Amina};
\node[draw=poporange, fill=poporangeL, rounded corners, font=\small, align=center] (d) at (6.5,1.6) {Ort und Datum\\(lieu et date)};
\node[draw=poporange, fill=poporangeL, rounded corners, font=\small, align=center] (a) at (6.5,0.6) {Anrede + Komma\\(salutation)};
\node[draw=poporange, fill=poporangeL, rounded corners, font=\small, align=center] (t) at (6.5,-0.6) {Text : minuscule !\\(corps)};
\node[draw=poporange, fill=poporangeL, rounded corners, font=\small, align=center] (g) at (6.5,-1.8) {Grußformel\\(formule finale)};
\draw[-{Stealth}, poporange, thick] (d.west) -- (4.35,1.2);
\draw[-{Stealth}, poporange, thick] (a.west) -- (4.35,0.35);
\draw[-{Stealth}, poporange, thick] (t.west) -- (4.35,-0.5);
\draw[-{Stealth}, poporange, thick] (g.west) -- (4.35,-1.5);
\end{tikzpicture}
\legende{Modèle de lettre personnelle avec ses quatre parties annotées.}
\end{popfigure}

\courssub{Le court essai argumenté (der kurze Aufsatz)}

L'essai argumenté est le texte roi de la classe de Terminale : on vous donne un sujet de société — par exemple \all{Die Frau in der heutigen Gesellschaft} — et vous devez présenter le sujet, peser le pour et le contre, puis donner votre opinion.

\begin{propriete}[Plan en trois parties]
\begin{enumerate}
\item \textbf{Einleitung} (introduction) : présenter le sujet et la question. Exemple : \all{Heute sprechen alle über die Gleichberechtigung von Mann und Frau. Aber was ist die Realität?}
\item \textbf{Hauptteil} (développement) : arguments pour et contre, avec des exemples concrets. On organise avec \all{einerseits ... andererseits} (d'un côté... de l'autre).
\item \textbf{Schluss} (conclusion) : votre opinion personnelle, justifiée : \all{Meiner Meinung nach ...} / \all{Ich bin der Ansicht, dass ...}
\end{enumerate}
\end{propriete}

\begin{popfigure}
\begin{tikzpicture}[node distance=0.6cm]
\node[draw=popblue, fill=popblueL, rounded corners, text width=3.2cm, align=center, font=\small] (ein) {\textbf{Einleitung}\\Sujet + question};
\node[draw=poppurple, fill=poppurpleL, rounded corners, text width=2.4cm, align=center, font=\small, below left=0.7cm and 0.4cm of ein] (pro) {\textbf{Hauptteil}\\Pro\\(arguments pour)};
\node[draw=poporange, fill=poporangeL, rounded corners, text width=2.4cm, align=center, font=\small, below right=0.7cm and 0.4cm of ein] (contra) {\textbf{Hauptteil}\\Contra\\(arguments contre)};
\node[draw=popgreen, fill=popgreenL, rounded corners, text width=3.2cm, align=center, font=\small, below=2.6cm of ein] (schluss) {\textbf{Schluss}\\Meine Meinung};
\draw[-{Stealth}, thick, popblue] (ein.south) -- (pro.north);
\draw[-{Stealth}, thick, popblue] (ein.south) -- (contra.north);
\draw[-{Stealth}, thick, popgreen] (pro.south) -- (schluss.north west);
\draw[-{Stealth}, thick, popgreen] (contra.south) -- (schluss.north east);
\end{tikzpicture}
\legende{Structure de l'essai argumenté : Einleitung, Hauptteil (Pro/Contra), Schluss.}
\end{popfigure}

\textbf{Les connecteurs d'argumentation.} Ce sont les articulations de votre texte ; sans eux, l'essai est une simple liste de phrases. Attention à l'ordre des mots qu'ils provoquent !

\begin{center}
\begin{tabularx}{\linewidth}{|l|X|X|}
\hline
\rowcolor{popblueL} Connecteur & Sens & Effet sur l'ordre des mots \\
\hline
\all{zuerst, dann, danach} & d'abord, ensuite & verbe en 2e position : \all{Zuerst zeigt er...} \\
\hline
\all{außerdem} & de plus & verbe en 2e position : \all{Außerdem arbeiten viele Frauen.} \\
\hline
\all{jedoch, aber} & cependant, mais & \all{aber} : position 0 (conjonction) ; \all{jedoch} : adverbe $\rightarrow$ inversion (V2) \\
\hline
\all{trotzdem} & malgré cela & verbe en 2e position : \all{Trotzdem macht sie die Hausarbeit.} \\
\hline
\all{deshalb, darum} & c'est pourquoi & verbe en 2e position : \all{Deshalb kämpfen sie.} \\
\hline
\all{meiner Meinung nach} & à mon avis & verbe en 2e position \\
\hline
\all{ich bin der Ansicht, dass...} & je suis d'avis que... & subordonnée : verbe à la fin \\
\hline
\end{tabularx}
\end{center}

\begin{exemplebox}[Phrases d'argumentation traduites]
\all{Meiner Meinung nach ist die Gleichberechtigung sehr wichtig.} « À mon avis, l'égalité des droits est très importante. »

\all{Einerseits arbeiten viele Frauen, andererseits machen sie mehr Hausarbeit.} « D'un côté beaucoup de femmes travaillent, de l'autre elles font plus de tâches ménagères. »

\all{Viele Mädchen gehen zur Schule, jedoch heiraten sie oft sehr früh.} « Beaucoup de filles vont à l'école, cependant elles se marient souvent très tôt. »

\all{Die Situation ist schwierig, trotzdem gibt es Hoffnung.} « La situation est difficile, il y a pourtant de l'espoir. »
\end{exemplebox}

\begin{attention}[« Je pense que » : verbe à la fin !]
« Je pense que » se traduit par \all{ich denke, dass...} ou \all{ich meine, dass...} avec le verbe conjugué \cle{à la fin} de la subordonnée : \all{Ich meine, dass die Frauen mehr Rechte haben müssen.} (« Je pense que les femmes doivent avoir plus de droits. »). Ne dites jamais \all{Ich meine, dass die Frauen müssen mehr Rechte haben} — calque du français, faute grave. Astuce : sans \all{dass}, on peut aussi dire \all{Ich denke, die Frauen müssen mehr Rechte haben} (verbe en 2e position).
\end{attention}

\courssub{Le dialogue (das Gespräch)}

Le dialogue écrit obéit à des règles typographiques allemandes précises :

\begin{propriete}[Règles du dialogue]
\begin{itemize}
\item Chaque réplique commence par un \cle{tiret} (ou est mise entre guillemets allemands „...“) et change de ligne à chaque changement de locuteur.
\item Les répliques sont courtes et orales : \all{Stimmt!}, \all{Wirklich?}, \all{Das sehe ich anders.}
\item Les verbes de parole (incises) : \all{sagen} (dire), \all{fragen} (demander), \all{antworten} (répondre), \all{meinen} (penser, dire). Dans l'incise, inversion sujet-verbe : \all{„Das ist ungerecht“, sagt Amina.}
\end{itemize}
\end{propriete}

\begin{exemplebox}[Mini-dialogue]
— \all{Was denkst du über die Hausarbeit?} (« Que penses-tu des tâches ménagères ? »)

— \all{Ich meine, dass Männer mehr helfen müssen}, antwortet Paul. (« Je pense que les hommes doivent aider davantage », répond Paul.)

— \all{Stimmt! Mein Vater kocht jetzt jeden Sonntag.} (« C'est vrai ! Mon père cuisine maintenant tous les dimanches. »)
\end{exemplebox}

\courssub{Méthode : le commentaire de texte}

\begin{methode}[Commenter un texte en 4 étapes]
\begin{enumerate}
\item \textbf{Einleitende Bemerkungen} : présenter le type de texte (article de journal, lettre, dialogue), le thème et, si possible, l'auteur ou la source. Formule : \all{Der Text ist ein Zeitungsartikel über...}
\item \textbf{Inhalt} : résumer le contenu en 2 ou 3 phrases, avec vos propres mots, au présent. Formule : \all{Der Autor zeigt zuerst..., dann...}
\item \textbf{Analyse} : étudier la structure et les moyens linguistiques : connecteurs, oppositions (\all{früher – heute}), champ lexical, arguments. Formule : \all{Er benutzt Gegensätze wie...}
\item \textbf{Eigene Meinung} : donner votre avis personnel et le justifier, en reliant le texte à votre réalité (le Cameroun !). Formule : \all{Ich finde das Thema wichtig, weil...}
\end{enumerate}
\end{methode}

\begin{textebox}[Commentaire modèle guidé (extrait)]
Der Text ist ein Zeitungsartikel über die Rolle der Frau in der heutigen Gesellschaft. Der Autor zeigt zuerst die traditionelle Rolle der Frau, dann die Veränderungen von heute. Er benutzt Gegensätze wie „früher – heute“ und viele Beispiele aus dem Alltag. Ich finde das Thema wichtig, weil die Gleichberechtigung auch in Kamerun ein großes Problem ist.
\end{textebox}

Glossaire : \all{die Veränderung, -en} : le changement ; \all{der Gegensatz, -e} : l'opposition, le contraste ; \all{der Alltag} : la vie quotidienne ; \all{die Gleichberechtigung} : l'égalité des droits ; \all{traditionell} : traditionnel.

Remarquez la progression en quatre paragraphes : présentation (type + thème), résumé (\all{zuerst... dann...}), analyse (\all{Gegensätze, Beispiele}), avis personnel justifié (\all{weil...}).

\courssub{Relecture et correction : la checklist}

Avant de rendre toute production, relisez-vous avec cette liste de contrôle :

\begin{aretenir}[Checklist de relecture]
\begin{enumerate}
\item \textbf{Majuscules} : tous les noms communs ont-ils une majuscule ? (\all{die Frau, das Recht, die Arbeit})
\item \textbf{Verbe en 2e position} dans chaque phrase principale, même après un connecteur (\all{Deshalb kämpft sie.}).
\item \textbf{Verbe à la fin} dans les subordonnées avec \all{dass, weil, obwohl}.
\item \textbf{Cas après prépositions} : \all{mit dem Vater} (datif), \all{für die Frauen} (akkusativ), \all{über die Rolle} (akkusativ).
\item \textbf{Accords} : article-adjectif-nom (\all{die traditionelle Rolle}), sujet-verbe.
\item \textbf{Orthographe} : Umlaute (\all{Mädchen, öffnen}), \all{ß/ss} (\all{groß, müssen}), lettre : minuscule après la salutation.
\end{enumerate}
\end{aretenir}

\begin{exoresolu}[Corriger une production d'élève]
Énoncé. Un élève a écrit dans son essai : \all{Ich denke, dass die frauen müssen mehr Rechte haben, trotzdem sie haben wenig Zeit.} Relevez et corrigez toutes les fautes (4 fautes).
\tcblower
Solution détaillée.
\begin{enumerate}
\item \all{die frauen} : majuscule obligatoire au nom → \all{die Frauen}.
\item \all{dass die Frauen müssen mehr Rechte haben} : dans la subordonnée en \all{dass}, le verbe conjugué va à la fin → \all{dass die Frauen mehr Rechte haben müssen}.
\item \all{trotzdem sie haben} : après \all{trotzdem} (adverbe connecteur), le verbe doit être en 2e position → \all{trotzdem haben sie wenig Zeit}.
\item La ponctuation est correcte, mais on préfère deux phrases ou un point-virgule pour la clarté.
\end{enumerate}
Phrase corrigée : \all{Ich denke, dass die Frauen mehr Rechte haben müssen; trotzdem haben sie wenig Zeit.} (« Je pense que les femmes doivent avoir plus de droits ; pourtant, elles ont peu de temps. »)
\end{exoresolu}

\begin{exercice}[À vous de produire]
\begin{enumerate}
\item Rédigez le début d'une lettre à votre ami Paul (lieu : Douala, date : 10 avril 2025) dans laquelle vous lui parlez du mariage de votre sœur. Respectez la structure et la minuscule après la salutation.
\item Écrivez l'introduction et un paragraphe \all{einerseits / andererseits} d'un essai sur le sujet : \all{Sollen Männer mehr im Haushalt helfen?} (« Les hommes doivent-ils aider davantage à la maison ? »)
\item Transformez en dialogue direct : Amina demande à Paul ce qu'il pense de l'égalité ; Paul répond qu'elle est très importante.
\end{enumerate}
\end{exercice}

\begin{corrige}[Exercice « À vous de produire »]
1. \all{Douala, den 10. April 2025 — Lieber Paul, — vielen Dank für deinen letzten Brief. Ich muss dir etwas Schönes erzählen: meine Schwester hat am Samstag geheiratet! ... Viele Grüße — [Nom].} Points de contrôle : \all{den 10. April} (accusatif), virgule après \all{Lieber Paul,}, minuscule à \all{vielen}, Perfekt \all{hat geheiratet}.

2. Exemple : \all{Heute arbeiten viele Frauen außer Haus. Aber wer macht die Hausarbeit? Einerseits haben die Männer oft keine Zeit, weil sie viel arbeiten. Andererseits sind die Frauen auch müde nach der Arbeit. Meiner Meinung nach müssen die Männer mehr helfen.}

3. \all{— Was denkst du über die Gleichberechtigung?, fragt Amina. — Ich meine, dass sie sehr wichtig ist, antwortet Paul.} (verbe à la fin dans \all{dass sie sehr wichtig ist} ; pas de virgule après le point d'interrogation dans le dialogue).
\end{corrige}

\begin{aretenir}[L'essentiel]
Lettre : date avec \all{den}, salutation + virgule, minuscule à la ligne, \all{Viele Grüße}. Essai : Einleitung — Hauptteil (\all{einerseits/andererseits}) — Schluss (\all{Meiner Meinung nach}). Dialogue : tirets, répliques courtes, verbes de parole avec inversion. Dans tous les cas : verbe en 2e position, verbe à la fin après \all{dass/weil}, majuscules aux noms, et relecture systématique avec la checklist.
\end{aretenir}

\coursec{Traduction : thème et version}

Dans la section précédente, nous avons appris à structurer une lettre, un court essai et un dialogue. Mais au baccalauréat comme dans la vie professionnelle, il ne suffit pas de produire : il faut aussi savoir passer d'une langue à l'autre sans trahir ni le sens ni la grammaire. C'est l'objet de cette dernière section : la \cle{version} (de l'allemand vers le français) et le \cle{thème} (du français vers l'allemand), appliqués aux thèmes du chapitre 1, et en particulier à la situation de la femme.

\courssub{Qu'est-ce que traduire ? Observer avant d'agir}

Traduire n'est pas remplacer chaque mot par son équivalent dans un dictionnaire. Traduire, c'est \cle{transférer un sens} d'une langue à une autre en respectant les structures propres à chacune. Observons d'abord deux phrases :

\begin{exemplebox}[Observer et comparer]
\begin{itemize}
\item \all{Die Frauen verdienen oft weniger als die Männer.} → « Les femmes gagnent souvent moins que les hommes. »
\item « Ma sœur s'est mariée l'année dernière. » → \all{Meine Schwester hat letztes Jahr geheiratet.}
\end{itemize}
Dans les deux cas, l'ordre des mots change, les temps ne se recouvrent pas exactement, et pourtant le sens est parfaitement conservé. C'est cela, une bonne traduction.
\end{exemplebox}

\begin{definition}[Version et thème]
La \cle{version} consiste à traduire un texte allemand vers le français. Le \cle{thème} consiste à traduire un texte français vers l'allemand. Dans les deux cas, on traduit par \textbf{phrases entières}, jamais mot à mot.
\end{definition}

\courssub{La méthode générale de la traduction}

Quel que soit le sens de la traduction, la démarche comporte quatre grandes étapes, résumées dans l'organigramme suivant :

\begin{popfigure}
\begin{tikzpicture}[
  node distance=0.6cm,
  etape/.style={rectangle, rounded corners, draw=popblue, fill=popblueL, thick, align=center, minimum width=4.2cm, minimum height=0.9cm, font=\small},
  detail/.style={rectangle, rounded corners, draw=poporange, fill=poporangeL, align=center, minimum width=5.6cm, minimum height=0.9cm, font=\small},
  fleche/.style={-{Stealth[length=3mm]}, thick, popdark}
]
\node[etape] (lire) {\textbf{1. LIRE} tout le texte};
\node[detail, below=of lire, align=center] (analyser){\textbf{2. ANALYSER} : Sujet / Verbe / Compléments\\ repérer les temps et les cas};
\node[etape, below=of analyser] (traduire) {\textbf{3. TRADUIRE} phrase par phrase};
\node[detail, below=of traduire, align=center] (relire){\textbf{4. RELIRE} : ordre des mots, cas,\\ orthographe, majuscules};
\draw[fleche] (lire) -- (analyser);
\draw[fleche] (analyser) -- (traduire);
\draw[fleche] (traduire) -- (relire);
\end{tikzpicture}
\legende{Organigramme de la méthode de traduction : lire, analyser, traduire, relire.}
\end{popfigure}

\begin{methode}[La version en 4 étapes (allemand → français)]
\begin{enumerate}
\item \textbf{Lire} le texte entier une première fois, sans rien écrire, pour comprendre le sens global (qui ? quoi ? quand ?).
\item \textbf{Analyser} chaque phrase : souligner le verbe conjugué, identifier le sujet (Nominativ) et les compléments (Akkusativ, Dativ, prépositions). Repérer le temps du verbe.
\item \textbf{Traduire} par phrases entières, en français naturel. Un Präteritum allemand devient souvent un passé composé ou un imparfait français selon le sens.
\item \textbf{Relire} : le français est-il correct et idiomatique ? A-t-on oublié un élément (négation, particule, complément) ?
\end{enumerate}
\end{methode}

\begin{methode}[Le thème en 5 étapes (français → allemand)]
\begin{enumerate}
\item \textbf{Temps du verbe français} : identifier le temps de chaque verbe (présent, passé composé, imparfait, futur).
\item \textbf{Équivalent allemand} : choisir le temps allemand correspondant (Präsens, Perfekt, Präteritum, Futur) selon le contexte (voir le tableau ci-dessous).
\item \textbf{Ordre des mots} : placer le verbe conjugué en \textbf{2\textsuperscript{e} position} dans la phrase principale, en \textbf{position finale} dans la subordonnée.
\item \textbf{Cas et prépositions} : vérifier le cas exigé par chaque verbe (\all{jmdm. helfen}) et chaque préposition (\all{mit + Dativ}, \all{für + Akkusativ}).
\item \textbf{Relecture} : majuscules à tous les noms, orthographe (ß, Umlaute), particules séparables renvoyées en fin de phrase.
\end{enumerate}
\end{methode}

\courssub{Le problème des temps : correspondances français–allemand}

C'est la difficulté majeure du thème. Le français et l'allemand n'ont pas les mêmes systèmes de temps. Le tableau suivant donne les équivalences à connaître :

\begin{center}
\begin{tabularx}{\linewidth}{|X|X|X|}
\hline
\rowcolor{popblueL} \textbf{Français} & \textbf{Allemand} & \textbf{Remarque} \\
\hline
Présent : elle travaille & Präsens : \all{sie arbeitet} & correspondance directe \\
\hline
Passé composé : elle a travaillé & Perfekt : \all{sie hat gearbeitet} & contexte oral, lettre, dialogue \\
\hline
Passé composé / passé simple (récit) & Präteritum : \all{sie arbeitete} & récit écrit, article, essai \\
\hline
Imparfait : elle travaillait & Präteritum : \all{sie arbeitete} & l'allemand n'a pas d'imparfait ! \\
\hline
Futur : elle travaillera & Futur I : \all{sie wird arbeiten} & ou Präsens + complément de temps \\
\hline
Conditionnel : elle travaillerait & Konjunktiv II : \all{sie würde arbeiten} & \\
\hline
\end{tabularx}
\end{center}

\begin{attention}[Passé composé français : Perfekt ou Präteritum ?]
Le passé composé français a deux traductions possibles en allemand. Dans une \textbf{lettre, un dialogue, un texte oral} : le \cle{Perfekt} (\all{Meine Schwester hat geheiratet.}). Dans un \textbf{récit, un article, un essai} : le \cle{Präteritum} (\all{Die Frauen kämpften für ihre Rechte.}). Pour \all{sein}, \all{haben} et les modaux, le Präteritum est presque toujours préféré : \all{Autrefois, les femmes restaient à la maison.} → \all{Früher blieben die Frauen zu Hause.} (imparfait français = Präteritum allemand).
\end{attention}

\courssub{Exemples guidés et commentés}

\begin{exemplebox}[Thème : analyse pas à pas]
\textbf{Phrase 1} : « Autrefois, les femmes restaient à la maison. »
\begin{itemize}
\item Temps : imparfait → Präteritum de \all{bleiben} : \all{blieben}.
\item Ordre : \all{Autrefois} en tête → le verbe reste en 2\textsuperscript{e} position, le sujet après le verbe.
\item Traduction : \all{Früher blieben die Frauen zu Hause.}
\end{itemize}
\textbf{Phrase 2} : « Ma sœur s'est mariée l'année dernière. »
\begin{itemize}
\item Temps : passé composé, contexte courant → Perfekt : \all{hat geheiratet} (\all{heiraten} se conjugue avec \all{haben}).
\item Traduction : \all{Meine Schwester hat letztes Jahr geheiratet.}
\end{itemize}
\textbf{Phrase 3} : « Je pense que les hommes doivent aussi faire la cuisine. »
\begin{itemize}
\item \all{que} → \all{dass} : subordonnée, verbe conjugué (\all{müssen}) en position finale.
\item Traduction : \all{Ich meine, dass die Männer auch kochen müssen.}
\end{itemize}
\end{exemplebox}

\begin{exemplebox}[Version : analyse pas à pas]
\textbf{Phrase 1} : \all{Viele junge Paare teilen sich die Hausarbeit.}
\begin{itemize}
\item Sujet : \all{Viele junge Paare} ; verbe : \all{teilen sich} (verbe réfléchi, à ne pas oublier !).
\item Traduction : « Beaucoup de jeunes couples se partagent les tâches ménagères. »
\end{itemize}
\textbf{Phrase 2} : \all{Trotzdem verdienen die Frauen oft weniger als die Männer.}
\begin{itemize}
\item \all{trotzdem} = pourtant ; \all{verdienen} = gagner (de l'argent) ; \all{weniger als} = moins que.
\item Traduction : « Pourtant, les femmes gagnent souvent moins que les hommes. »
\end{itemize}
\end{exemplebox}

\courssub{Les pièges de traduction : faux amis, cas et expressions}

\begin{attention}[Faux amis à connaître par cœur]
\begin{itemize}
\item \all{bekommen} = \textbf{recevoir} (et non « devenir » !). « Devenir » se dit \all{werden}. \all{Sie hat ein Geschenk bekommen.} = « Elle a reçu un cadeau. »
\item \all{aktuell} = \textbf{actuel, en ce moment} (et non « réel, véritable »). \all{das aktuelle Problem} = « le problème actuel ».
\item \all{die Mappe, -n} = la chemise, la pochette (et non « la mappe »).
\item \all{die Fabrik} = l'usine (et non « la fabrique » au sens abstrait).
\end{itemize}
\end{attention}

\begin{propriete}[« Il y a » = \all{es gibt} + Akkusativ]
L'expression française « il y a » se traduit par \cle{es gibt}, et le nom qui suit est à l'\textbf{Akkusativ} : \all{Es gibt viele Probleme.} = « Il y a beaucoup de problèmes. » Au passé : \all{Es gab viele Probleme.} Ne dites \textbf{jamais} \all{es hat} pour « il y a » : c'est un germanisme inexistant, erreur typique des francophones.
\end{propriete}

\begin{propriete}[« Beaucoup de » : \all{viele} ou \all{viel} ?]
Devant un nom \textbf{comptable} (pluriel), on utilise \all{viele} : \all{viele Frauen, viele Probleme, viele junge Paare}. Devant un nom \textbf{non comptable} (singulier), on utilise \all{viel} : \all{viel Arbeit} (beaucoup de travail), \all{viel Zeit} (beaucoup de temps), \all{viel Geld} (beaucoup d'argent).
\end{propriete}

\begin{attention}[Les cas après les verbes : le piège du Dativ]
Certains verbes allemands régissent le \textbf{Dativ} là où le français utilise un complément direct :
\begin{itemize}
\item « aider quelqu'un » = \all{jmdm. (Dat.) helfen} : \all{Ich helfe meiner Mutter.} (et non \all{meine Mutter} !)
\item « remercier quelqu'un » = \all{jmdm. danken} ; « téléphoner à quelqu'un » = \all{jmdm. telefonieren}.
\item « s'occuper de » = \all{sich kümmern um + Akkusativ} : \all{Sie kümmert sich um die Kinder.} = « Elle s'occupe des enfants. »
\end{itemize}
Avant de traduire, demandez-vous toujours : quel cas ce verbe allemand exige-t-il ?
\end{attention}

\begin{exoresolu}[Thème guidé : cinq phrases sur la situation de la femme]
Traduisez en allemand : 1) « Il y a encore beaucoup d'inégalités entre les hommes et les femmes. » 2) « Ma grand-mère n'a pas pu aller à l'université. » 3) « De plus en plus de femmes travaillent aujourd'hui. » 4) « Mon père aide ma mère à la maison. » 5) « Je crois que l'égalité est importante. »
\tcblower
\textbf{Solutions commentées :}
\begin{enumerate}
\item « Il y a » → \all{es gibt} + Akk. ; « inégalités » = \all{die Ungleichheiten} : \all{Es gibt noch viele Ungleichheiten zwischen den Männern und den Frauen.} (\all{zwischen} + Dativ !)
\item Passé composé → Perfekt ; « pouvoir » au participe : \all{gekonnt} ou construction avec infinitif : \all{Meine Großmutter hat nicht zur Universität gehen können.} (double infinitif avec modal au Perfekt).
\item \all{Immer mehr Frauen arbeiten heute.} (« de plus en plus de » = \all{immer mehr} + nom sans article).
\item \all{helfen} + Dativ : \all{Mein Vater hilft meiner Mutter zu Hause.} (« à la maison » = \all{zu Hause}).
\item Subordonnée avec \all{dass}, verbe en finale : \all{Ich glaube, dass die Gleichberechtigung wichtig ist.}
\end{enumerate}
\end{exoresolu}

\begin{exoresolu}[Version guidée : cinq phrases]
Traduisez en français : 1) \all{Früher durften die Mädchen in vielen Familien nicht zur Schule gehen.} 2) \all{Heute studieren mehr Frauen als früher.} 3) \all{Mein Bruder kümmert sich um die Kinder, wenn seine Frau arbeitet.} 4) \all{Die Frauen haben lange für das Wahlrecht gekämpft.} 5) \all{Es gibt Länder, wo die Mädchen noch keine Schule besuchen dürfen.}
\tcblower
\textbf{Solutions commentées :}
\begin{enumerate}
\item \all{durften} = Präteritum de \all{dürfen} : « Autrefois, dans beaucoup de familles, les filles n'avaient pas le droit d'aller à l'école. »
\item « Aujourd'hui, il y a plus de femmes qui font des études qu'autrefois. »
\item \all{sich kümmern um} = s'occuper de : « Mon frère s'occupe des enfants quand sa femme travaille. »
\item \all{das Wahlrecht} = le droit de vote ; Perfekt → passé composé : « Les femmes ont longtemps lutté pour le droit de vote. »
\item \all{Es gibt} = il y a : « Il y a des pays où les filles ne peuvent toujours pas aller à l'école. »
\end{enumerate}
\end{exoresolu}

\courssub{Textes d'entraînement : version et thème}

\begin{textebox}[Version : \all{Die Frauen in Kamerun}]
\all{In Kamerun spielen die Frauen eine wichtige Rolle in der Wirtschaft. Auf den Märkten von Yaoundé und Douala verkaufen sie Obst, Gemüse und Stoffe. Viele von ihnen haben keine Schule besucht, aber sie ernähren ihre Familien und bezahlen die Schulgebühren ihrer Kinder. Früher blieben die meisten Frauen zu Hause und kümmerten sich nur um die Kinder und das Feld. Heute arbeiten immer mehr Frauen als Lehrerinnen, Ärztinnen oder Ingenieurinnen. Trotzdem gibt es noch Probleme: Die Frauen verdienen oft weniger als die Männer, und in manchen Dörfern dürfen die Mädchen nicht lange zur Schule gehen. Viele Organisationen kämpfen deshalb für die Gleichberechtigung der Frau.}

\textbf{Glossaire :} die Wirtschaft (l'économie) ; der Stoff, -e (le tissu) ; ernähren (nourrir) ; die Schulgebühr, -en (les frais de scolarité) ; das Feld, -er (le champ) ; die Ärztin, -nen (la femme médecin) ; das Dorf, -er (le village) ; die Gleichberechtigung (l'égalité des droits).
\end{textebox}

\begin{exercice}[Version : traduisez le texte « Die Frauen in Kamerun »]
Traduisez le texte ci-dessus en français. Appliquez la méthode : lecture globale, analyse des verbes (\all{spielen, haben besucht, ernähren, blieben, kümmerten, arbeiten, verdienen, dürfen, kämpfen}), traduction phrase par phrase, relecture.
\end{exercice}

\begin{corrige}[Version : proposition de traduction]
« Au Cameroun, les femmes jouent un rôle important dans l'économie. Sur les marchés de Yaoundé et de Douala, elles vendent des fruits, des légumes et des tissus. Beaucoup d'entre elles ne sont pas allées à l'école, mais elles nourrissent leurs familles et paient les frais de scolarité de leurs enfants. Autrefois, la plupart des femmes restaient à la maison et ne s'occupaient que des enfants et du champ. Aujourd'hui, de plus en plus de femmes travaillent comme enseignantes, médecins ou ingénieures. Pourtant, il y a encore des problèmes : les femmes gagnent souvent moins que les hommes, et dans certains villages, les filles n'ont pas le droit d'aller longtemps à l'école. C'est pourquoi de nombreuses organisations luttent pour l'égalité des droits de la femme. »
\end{corrige}

\begin{exercice}[Thème : traduisez ce texte en allemand]
« Autrefois, ma grand-mère ne pouvait pas choisir son mari. Elle s'est mariée très jeune et elle a eu huit enfants. Elle n'est jamais allée à l'école, mais elle a toujours voulu que ses filles fassent des études. Aujourd'hui, ma tante est professeur à l'université de Yaoundé. Je pense que la situation des femmes a beaucoup changé, mais il y a encore beaucoup de travail. Les hommes et les femmes doivent se partager les tâches ménagères et s'occuper ensemble des enfants. »
\end{exercice}

\begin{corrige}[Thème : proposition de traduction]
\all{Früher konnte meine Großmutter ihren Mann nicht wählen. Sie hat sehr jung geheiratet und sie hat acht Kinder bekommen. Sie ist nie zur Schule gegangen, aber sie hat immer gewollt, dass ihre Töchter studieren. Heute ist meine Tante Professorin an der Universität von Yaoundé. Ich meine, dass sich die Situation der Frauen sehr verändert hat, aber es gibt noch viel Arbeit. Die Männer und die Frauen müssen sich die Hausarbeit teilen und sich zusammen um die Kinder kümmern.}

\textbf{Points de vigilance :} \all{konnte} (Präteritum du modal, récit) ; \all{hat acht Kinder bekommen} (\all{bekommen} = avoir/recevoir des enfants) ; \all{ist nie zur Schule gegangen} (Perfekt avec \all{sein}) ; \all{sich verändern} (verbe réfléchi) ; \all{sich die Hausarbeit teilen} ; \all{sich um die Kinder kümmern} (+ Akkusativ).
\end{corrige}

\begin{aretenir}[Les dix commandements du traducteur]
\begin{enumerate}
\item Je lis tout le texte avant de traduire.
\item Je traduis par phrases entières, jamais mot à mot.
\item Je repère le sujet, le verbe et les compléments.
\item Passé composé français : Perfekt (oral/lettre) ou Präteritum (récit).
\item Imparfait français = Präteritum allemand.
\item Verbe conjugué en 2\textsuperscript{e} position, en finale après \all{dass}, \all{weil}, \all{wenn}.
\item « Il y a » = \all{es gibt} + Akkusativ.
\item \all{bekommen} = recevoir ; « devenir » = \all{werden}.
\item \all{helfen}, \all{danken} + Dativ ; \all{sich kümmern um} + Akkusativ.
\item Je relis : majuscules, Umlaute, particules séparables, cas.
\end{enumerate}
\end{aretenir}

\begin{savaistu}[Version et thème au baccalauréat]
Au baccalauréat camerounais, la version et le thème portent très souvent sur de courts textes de société : la famille, l'école, la situation de la femme, le travail des jeunes. La grille de correction valorise autant le respect de l'\textbf{ordre des mots} et des \textbf{cas} que la richesse du vocabulaire. Une phrase simple mais grammaticalement parfaite rapporte plus de points qu'une phrase ambitieuse pleine de fautes de cas. Entraînez-vous donc à traduire régulièrement de courtes phrases : c'est la gymnastique quotidienne du bon germaniste !
\end{savaistu}

\newpage

\coursec{Activité d'intégration}

\courssub{Objectifs de l'activité}

À l'issue de cette activité d'intégration, tu seras capable de :

\begin{itemize}
\item rédiger une \cle{lettre personnelle} complète en allemand (12 à 15 lignes) respectant la structure apprise : en-tête avec lieu et date, formule d'appel, introduction, corps, conclusion, formule de politesse ;
\item raconter un \cle{événement familial} au passé en choisissant correctement entre \cle{Perfekt} et \cle{Präteritum} et en accordant le participe passé avec le bon auxiliaire (\all{haben} ou \all{sein}) ;
\item comparer la \cle{situation des femmes} au Cameroun et en Allemagne en mobilisant le lexique thématique du chapitre (\all{die Gleichberechtigung}, \all{der Beruf}, \all{die Hausarbeit}, \all{die Tradition}) ;
\item utiliser au moins deux \cle{connecteurs d'argumentation} (\all{einerseits \dots\ andererseits}, \all{trotzdem}, \all{deshalb}) en plaçant correctement le verbe ;
\item relire, corriger et \cle{traduire} la lettre d'un camarade (version allemand $\rightarrow$ français) à l'aide d'une checklist.
\end{itemize}

\courssub{Situation}

Tu participes depuis deux ans à un échange scolaire entre ton lycée de Yaoundé et le Gymnasium de Bayreuth (Bavière). Ta correspondante allemande, Lena, t'a écrit la semaine dernière. Dans sa lettre, elle te pose des questions sur ta famille et sur la place des femmes dans la société camerounaise, car elle prépare un exposé (\all{ein Referat}) sur le thème \all{„Frauen in Afrika und in Europa“} pour son cours de géographie. Voici un extrait de sa lettre :

\begin{textebox}[Extrait de la lettre de Lena]
Liebe Amina,

vielen Dank für deinen letzten Brief! Bei uns ist alles gut. Nächste Woche muss ich in der Schule ein Referat über das Thema „Frauen in Afrika und in Europa“ halten, und dabei kannst du mir bestimmt helfen!

Erzähl mir doch bitte zuerst von deiner Familie: Was habt ihr in den letzten Wochen erlebt? Du hast doch sicher ein Fest oder eine Feier gehabt?

Und dann interessiert mich: Wie ist die Situation der Frauen in Kamerun? Arbeiten die meisten Frauen in einem Beruf, oder bleiben sie zu Hause? Haben sie die gleichen Rechte wie die Männer? Hier in Deutschland diskutieren wir oft über Gleichberechtigung, zum Beispiel über Frauen in Chefsessel-Positionen.

Schreib mir bald zurück!

Viele liebe Grüße aus Bayreuth

Deine Lena
\end{textebox}

\textbf{Glossaire :} \all{das Referat, -e} (l'exposé) ; \all{erleben} (vivre, éprouver) ; \all{die Feier, -n} (la fête) ; \all{die gleichen Rechte} (les mêmes droits) ; \all{die Gleichberechtigung} (l'égalité des droits) ; \all{der Chefsessel} (le fauteuil de chef, ici : poste de direction).

\textbf{Ta mission :} tu réponds à Lena par une lettre de 12 à 15 lignes dans laquelle tu (1) racontes un événement familial récent (mariage, naissance, fête de famille) et (2) compares la situation des femmes au Cameroun et en Allemagne. Puis tu échanges ta lettre avec un autre binôme qui la traduit en français et la corrige.

\courssub{Consignes}

\textbf{Étape 1 — Compréhension et repérage (10 min).} Lis la lettre de Lena. Souligne les trois questions qu'elle te pose. Note dans ton cahier les mots du champ lexical de la famille et de l'égalité que tu reconnais.

\textbf{Étape 2 — Préparation du plan (10 min).} Par binômes, complétez ensemble le plan de votre réponse :
\begin{enumerate}
\item En-tête : lieu et date (\all{Yaoundé, den 15. Mai 2025}).
\item Formule d'appel : \all{Liebe Lena,} — attention, la phrase suivante commence par une \textbf{minuscule}.
\item Remerciement et transition : remercie Lena pour sa lettre.
\item Paragraphe 1 : raconte un événement familial récent au \textbf{passé} (au moins 4 verbes au Perfekt ou au Präteritum).
\item Paragraphe 2 : compare la situation des femmes au Cameroun et en Allemagne avec \textbf{au moins deux connecteurs} : \all{einerseits \dots\ andererseits}, \all{trotzdem}, \all{deshalb}.
\item Conclusion : pose une question à Lena et annonce ta prochaine lettre.
\item Formule de politesse : \all{Viele liebe Grüße} + ta signature.
\end{enumerate}

\textbf{Étape 3 — Rédaction (25 min).} Rédigez la lettre (12 à 15 lignes). Chaque membre du binôme écrit au moins un paragraphe, puis vous harmonisez le texte. Vérifiez au fur et à mesure : verbe conjugué en \textbf{2e position}, majuscule à tous les \textbf{noms}, participe passé à la \textbf{fin} de la phrase au Perfekt.

\textbf{Étape 4 — Échange, version et correction (20 min).} Échangez votre lettre avec un autre binôme. Ce binôme :
\begin{enumerate}
\item traduit la lettre reçue en français (version) sur une feuille séparée ;
\item corrige la lettre à l'aide de la checklist ci-dessous et note ses remarques en marge.
\end{enumerate}

\begin{methode}[Checklist de correction]
\begin{enumerate}
\item \textbf{Majuscules} : tous les noms commencent-ils par une majuscule ?
\item \textbf{Place du verbe} : le verbe conjugué est-il en 2e position dans chaque phrase principale ? Après \all{trotzdem} et \all{deshalb}, le verbe suit-il immédiatement (position 2) ?
\item \textbf{Passé} : auxiliaire correct (\all{haben} ou \all{sein} pour les verbes de mouvement et de changement) ? Participe passé bien formé (\all{ge-...-t} / \all{ge-...-en}) ?
\item \textbf{Cas} : article et adjectif corrects après les prépositions (\all{mit der Familie}, \all{für deinen Brief}) ?
\item \textbf{Connecteurs} : au moins deux connecteurs d'argumentation sont-ils présents et correctement employés ?
\item \textbf{Structure de la lettre} : en-tête, appel, minuscule après l'appel, formule finale, signature ?
\end{enumerate}
\end{methode}

\textbf{Étape 5 — Restitution orale (10 min).} Deux binômes volontaires lisent leur lettre à voix haute. La classe relève les points forts (lexique riche, connecteurs bien placés, passé correct) et propose une amélioration par lettre.

\courssub{Ressources à mobiliser}

\begin{aretenir}[Lexique utile pour la lettre]
\begin{itemize}
\item Famille et fêtes : \all{die Hochzeit, -en} (le mariage) ; \all{die Geburt, -en} (la naissance) ; \all{die Feier, -n} ; \all{heiraten} (se marier) ; \all{ein Kind bekommen} (avoir un enfant).
\item Femmes et société : \all{die Frau, -en} ; \all{die Gleichberechtigung} ; \all{der Beruf, -e} (le métier) ; \all{berufstätig sein} (travailler, exercer un métier) ; \all{die Hausarbeit} (les tâches ménagères) ; \all{die Tradition, -en} ; \all{die Rechte} (les droits) ; \all{verdienen} (gagner).
\item Connecteurs : \all{einerseits \dots\ andererseits} (d'une part \dots\ d'autre part) ; \all{trotzdem} (pourtant, néanmoins) ; \all{deshalb} (c'est pourquoi) ; \all{außerdem} (de plus) ; \all{meiner Meinung nach} (à mon avis).
\end{itemize}
\end{aretenir}

\begin{propriete}[Rappels grammaticaux essentiels]
\begin{itemize}
\item \textbf{Perfekt} (langage courant, lettres) : auxiliaire \all{haben}/\all{sein} en 2e position + participe passé en fin de phrase. \all{Meine Schwester hat letzten Monat geheiratet.} « Ma sœur s'est mariée le mois dernier. » — \all{Wir sind zu der Feier gefahren.} « Nous sommes allés à la fête. »
\item \textbf{Präteritum} : pour \all{sein}, \all{haben}, les modaux et les verbes fréquents : \all{Es war ein schönes Fest.} « C'était une belle fête. » — \all{Es gab viel zu essen.} « Il y avait beaucoup à manger. »
\item \textbf{Connecteurs} : \all{trotzdem} et \all{deshalb} occupent la 1re position, le verbe reste en 2e position (inversion du sujet) : \all{Viele Frauen arbeiten, trotzdem machen sie die meiste Hausarbeit.} « Beaucoup de femmes travaillent, pourtant elles font la plupart des tâches ménagères. »
\item \textbf{Après l'appel} d'une lettre (\all{Liebe Lena,}), la phrase commence par une \textbf{minuscule}.
\end{itemize}
\end{propriete}

\begin{attention}[Pièges fréquents des francophones]
\begin{itemize}
\item Ne dis pas \all{*Ich habe gefahren} : les verbes de déplacement exigent \all{sein} : \all{ich bin gefahren}.
\item Après \all{trotzdem}, pas d'ordre sujet-verbe français : \all{trotzdem machen sie\dots} et non \all{*trotzdem sie machen\dots}.
\item Faux ami : \all{bekommen} signifie « recevoir, obtenir », pas « devenir » (\all{werden}).
\item \all{die Frau} = « la femme » (épouse ou femme en général) ; « une femme » au sens de « épouse » se traduit aussi par \all{Ehefrau} dans les contextes précis.
\end{itemize}
\end{attention}

\courssub{Proposition de production et corrigé commenté}

\begin{exoresolu}[Lettre modèle : réponse à Lena]
\textbf{Production modèle (à lire après ta propre rédaction) :}

\all{Yaoundé, den 15. Mai 2025}

\all{Liebe Lena,}

\all{vielen Dank für deinen Brief! Dein Referat ist eine tolle Idee, und ich helfe dir gern.}

\all{Zuerst erzähle ich dir von meiner Familie: Letzten Monat hat meine ältere Schwester Sandrine geheiratet. Die Hochzeit war ein großes Fest. Wir sind am Morgen in die Kirche gegangen, und am Nachmittag haben wir mit mehr als zweihundert Gästen gefeiert. Es gab Poulet DG, Ndolé und viel Musik. Meine Mutter hat den ganzen Tag getanzt!}

\all{Jetzt zu deinen Fragen: In Kamerun arbeiten heute viele Frauen in einem Beruf, zum Beispiel als Lehrerinnen, Ärztinnen oder Händlerinnen auf dem Markt. Einerseits haben sie die gleichen Rechte wie die Männer, andererseits verdienen sie oft weniger Geld. Viele Frauen haben einen Beruf, trotzdem machen sie zu Hause die meiste Hausarbeit. Deshalb sagen viele junge Frauen bei uns: Die Gleichberechtigung muss auch in der Familie beginnen. Meiner Meinung nach ist das in Deutschland ähnlich, oder?}

\all{Und wie läuft dein Referat? Schreib mir bald!}

\all{Viele liebe Grüße aus Yaoundé}

\all{Deine Amina}

\tcblower

\textbf{Commentaire des choix de langue :}
\begin{itemize}
\item \textbf{Structure} : en-tête avec \all{den} + date (accusatif), appel suivi d'une minuscule (\all{vielen Dank\dots}), formule finale sans virgule avant la signature.
\item \textbf{Temps du passé} : Perfekt pour les événements racontés (\all{hat geheiratet}, \all{haben gefeiert}, \all{hat getanzt}) ; Präteritum pour \all{sein} et \all{es gab} (\all{war}, \all{gab}), plus naturel à l'écrit. Auxiliaire \all{sein} avec le verbe de mouvement : \all{wir sind \dots\ gegangen}.
\item \textbf{Cas} : \all{für deinen Brief} (accusatif après \all{für}) ; \all{mit mehr als zweihundert Gästen} (datif pluriel après \all{mit}) ; \all{in die Kirche} (accusatif, mouvement).
\item \textbf{Connecteurs} : \all{einerseits \dots\ andererseits} en tête de proposition avec verbe en 2e position ; \all{trotzdem} et \all{deshalb} suivis immédiatement du verbe (\all{trotzdem machen sie}, \all{deshalb sagen viele junge Frauen}).
\item \textbf{Lexique du chapitre} : \all{die Hochzeit}, \all{die gleichen Rechte}, \all{die Hausarbeit}, \all{die Gleichberechtigung}, \all{berufstätig} (ici : \all{in einem Beruf arbeiten}).
\end{itemize}

\textbf{Version française (traduction attendue à l'étape 4) :}

« Yaoundé, le 15 mai 2025. Chère Lena, merci beaucoup pour ta lettre ! Ton exposé est une excellente idée, et je t'aide volontiers. D'abord, je te raconte ce qui s'est passé dans ma famille : le mois dernier, ma grande sœur Sandrine s'est mariée. Le mariage était une grande fête. Le matin, nous sommes allés à l'église, et l'après-midi, nous avons fait la fête avec plus de deux cents invités. Il y avait du poulet DG, du ndolé et beaucoup de musique. Ma mère a dansé toute la journée ! Maintenant, à tes questions : au Cameroun, beaucoup de femmes exercent aujourd'hui un métier, par exemple comme institutrices, médecins ou commerçantes au marché. D'une part, elles ont les mêmes droits que les hommes, d'autre part, elles gagnent souvent moins d'argent. Beaucoup de femmes ont un métier, pourtant elles font la plupart des tâches ménagères à la maison. C'est pourquoi beaucoup de jeunes femmes disent chez nous : l'égalité des droits doit aussi commencer dans la famille. À mon avis, c'est pareil en Allemagne, non ? Et comment avance ton exposé ? Écris-moi vite ! Amitiés de Yaoundé. Ta Amina. »
\end{exoresolu}

\courssub{Bilan de l'activité}

Cette activité t'a permis de réinvestir l'ensemble des acquis du chapitre dans une tâche authentique : écrire à une vraie destinataire, raconter au passé et argumenter. Vérifie que tu maîtrises désormais :

\begin{itemize}
\item la \textbf{structure de la lettre} allemande (en-tête, appel + minuscule, corps, formule finale) ;
\item le choix entre \textbf{Perfekt} et \textbf{Präteritum} et le bon auxiliaire (\all{haben}/\all{sein}) ;
\item l'\textbf{ordre des mots} après les connecteurs d'argumentation (\all{trotzdem}, \all{deshalb} : verbe en 2e position) ;
\item le \textbf{lexique} de la famille et de la situation des femmes ;
\item la \textbf{traduction} fidèle d'une lettre (version) et la correction méthodique d'une production à l'aide d'une checklist.
\end{itemize}

\begin{experience}[Pour aller plus loin : débat en classe]
En groupes de quatre, préparez un mini-débat sur la question : \all{„Soll die Hausarbeit zwischen Mann und Frau geteilt werden?“} (« Les tâches ménagères doivent-elles être partagées entre l'homme et la femme ? »). Chaque groupe présente deux arguments \all{dafür} (pour) et deux arguments \all{dagegen} (contre), en utilisant obligatoirement \all{einerseits \dots\ andererseits}, \all{trotzdem} et \all{deshalb}. La classe vote ensuite pour le groupe le plus convaincant.
\end{experience}

\newpage

\coursec{Méthodes et exercices résolus}

\coursec{Leçon AL05 : Production et traduction — La situation de la femme et vie familiale/sociale}

\courssub{Texte support : Die Frau in der heutigen Gesellschaft}

\begin{textebox}[Texte original — Contexte camerounais et germanophone]
In Kamerun wie in Deutschland hat sich die Rolle der Frau in den letzten Jahrzehnten stark gewandelt. Früher war die Frau fast nur für den Haushalt und die Erziehung der Kinder zuständig. Heute sind Frauen in fast allen Berufen vertreten: als Ärztinnen, Ingenieurinnen, Lehrerinnen, Unternehmerinnen oder Politikerinnen. In Yaoundé und Douala leiten viele Frauen ihre eigenen Geschäfte auf dem Marché Central oder im quartier d'affaires. Auch in der Politik gibt es Fortschritte: im kamerunischen Parlament sitzen heute mehr Abgeordneten als vor zwanzig Jahren.

In Deutschland ist die Gleichberechtigung im Grundgesetz verankert (Artikel 3). Dennoch gibt es noch Herausforderungen: der Gender Pay Gap (Lohnlücke) beträgt immer noch etwa 18 \%, und Frauen übernehmen den größten Teil der unbezahlten Care-Arbeit (Haushalt, Pflege). In der Schweiz und in Österreich sieht die Lage ähnlich aus. Der Weg zur echten Gleichstellung ist also noch nicht zu Ende, aber die Entwicklung geht in die richtige Richtung. Bildung bleibt der Schlüssel.
\end{textebox}

\courssub{Lexique thématique : Famille, mariage, travail, égalité}

\begin{aretenir}[Vocabulaire de base — Noms (article + pluriel), verbes, adjectifs, connecteurs]
\begin{center}
\begin{tabularx}{\linewidth}{|X|X|X|}
\hline
\rowcolor{popblueL} \textbf{Français} & \textbf{Deutsch} & \textbf{Remarque} \\
\hline
la femme / les femmes & \all{die Frau, die Frauen} & Umlaut au pluriel \\
l'homme / les hommes & \all{der Mann, die Männer} & Umlaut au pluriel \\
la fille / les filles & \all{das Mädchen, die Mädchen} & neutre, pluriel invariant \\
le garçon / les garçons & \all{der Junge, die Jungen} & \\
la mère / les mères & \all{die Mutter, die Mütter} & Umlaut au pluriel \\
le père / les pères & \all{der Vater, die Väter} & Umlaut au pluriel \\
les parents & \all{die Eltern} & seulement pluriel \\
le mariage & \all{die Ehe} & \textbf{faux ami} : \all{die Ehe} $\neq$ l'huile \\
le divorce & \all{die Scheidung} & \\
le ménage / les tâches ménagères & \all{der Haushalt} / \all{die Hausarbeit} & \all{Hausarbeit} = travail ménager (sing.) \\
le travail / l'emploi & \all{die Arbeit} / \all{der Beruf} / \all{die Stelle} & \all{die Stelle} = le poste \\
le salaire & \all{der Lohn} (ouvrier) / \all{das Gehalt} (employé) & distinction importante \\
l'égalité / l'inégalité & \all{die Gleichberechtigung} / \all{die Ungleichheit} & \all{Gleichberechtigung} = égalité des droits \\
les droits (de l'homme) & \all{die Rechte} (Pl.) / \all{die Menschenrechte} & \\
la discrimination & \all{die Benachteiligung} & \all{benachteiligen} (verbe) \\
la promotion / l'avancement & \all{der Aufstieg} & \\
le progrès & \all{der Fortschritt} & \\
le défi / le problème & \all{die Herausforderung} / \all{das Problem} & \\
\hline
\end{tabularx}
\end{center}

\begin{center}
\begin{tabularx}{\linewidth}{|X|X|X|X|}
\hline
\rowcolor{poporangeL} \textbf{Verbes clés} & \textbf{Präteritum} & \textbf{Perfekt} & \textbf{Remarque} \\
\hline
travailler & \all{arbeiten} — \all{arbeitete} & \all{hat gearbeitet} \\
étudier / faire des études & \all{studieren} — \all{studierte} & \all{hat studiert} \\
gagner (argent) & \all{verdienen} — \all{verdiente} & \all{hat verdient} \\
recevoir / obtenir & \all{bekommen} — \all{bekam} & \all{hat bekommen} & \textbf{faux ami} : $\neq$ devenir \\
devenir & \all{werden} — \all{wurde} & \all{ist geworden} & aux. \all{sein} \\
se battre pour & \all{kämpfen für} — \all{kämpfte} & \all{hat gekämpft} \\
avoir le droit / pouvoir & \all{dürfen} — \all{durfte} & \all{hat gedurft} & modal \\
devoir / avoir à & \all{müssen} — \all{musste} & \all{hat gemusst} & modal \\
changer & \all{sich verändern} — \all{veränderte sich} & \all{hat sich verändert} & réfléchi \\
partager & \all{teilen} — \all{teilte} & \all{hat geteilt} & fort (vowel change) \\
aider & \all{helfen} — \all{half} & \all{hat geholfen} & \textbf{+ Dativ} ! \\
\hline
\end{tabularx}
\end{center}

\begin{center}
\begin{tabularx}{\linewidth}{|X|X|}
\hline
\rowcolor{popgreenL} \textbf{Connecteurs / Adverbes} & \textbf{Français} \\
\hline
\all{einerseits ... andererseits} & d'une part ... d'autre part \\
\all{zuerst, dann, schließlich} & d'abord, ensuite, enfin \\
\all{außerdem / außerdem} & de plus / en outre \\
\all{jedoch / aber} & cependant / mais \\
\all{deshalb / daher / darum} & c'est pourquoi / donc \\
\all{trotzdem} & pourtant / néanmoins \\
\all{meiner Meinung nach / ich finde, dass ...} & à mon avis / je trouve que ... \\
\all{seit + Dativ + Präsens} & depuis + présent (action commencée dans le passé et qui continue) \\
\hline
\end{tabularx}
\end{center}
\end{aretenir}

\courssub{Grammaire 1 : Ordre des mots et cas}

\begin{propriete}[Ordre des mots : phrase principale vs subordonnée]
\textbf{1. Phrase principale (Hauptsatz) :} Le verbe conjugué est \textbf{toujours en 2\textsuperscript{e} position} (règle V2). Le sujet peut être en 1\textsuperscript{e} ou 3\textsuperscript{e} position (inversion).
\begin{center}
\begin{tabular}{|l|l|l|l|l|}
\hline
\rowcolor{popblueL} Pos 1 & \textbf{Pos 2 (Verbe)} & Pos 3 (Sujet si inversion) & Milieu (TeKaMoLo) & \textbf{Pos finale} \\
\hline
\all{Ich} & \all{lerne} & — & \all{heute Deutsch} & — \\
\all{Heute} & \all{lerne} & \all{ich} & \all{Deutsch} & — \\
\all{Das Buch} & \all{lese} & \all{ich} & \all{gerne} & — \\
\all{Meiner Meinung nach} & \all{ist} & \all{das} & \all{wichtig} & — \\
\hline
\end{tabular}
\end{center}
\textbf{TeKaMoLo} (ordre usuel au milieu) : \textbf{Te}mp (temps) — \textbf{Ka}usal (cause) — \textbf{Mo}dal (manière) — \textbf{Lo}kal (lieu). Ex. : \all{Ich fahre morgen mit dem Zug nach Berlin.}

\textbf{2. Subordonnée (Nebensatz) :} Introduite par \all{dass, weil, obwohl, wenn, als, ob...}. Le verbe conjugué va \textbf{à la toute fin}. Virgule obligatoire avant la conjonction.
\all{Ich denke, dass Frauen dieselben Rechte \textbf{haben}.} \\
\all{Weil sie weniger \textbf{verdienen}, sind sie unzufrieden.} (Verbe principal \all{sind} en 2\textsuperscript{e} pos après la subordonnée).

\textbf{3. Verbes à particule séparable (trennbare Verben) :} Particule à la fin en phrase principale, verbe entier à la fin en subordonnée.
\all{Er steht um 6 Uhr \textbf{auf}.} — \all{..., dass er um 6 Uhr \textbf{aufsteht}.}

\textbf{4. Cas après prépositions (révision) :}
\begin{center}
\begin{tabular}{|l|l|l|}
\hline
\rowcolor{poppurpleL} \textbf{Préposition} & \textbf{Cas} & \textbf{Exemple} \\
\hline
\all{für, ohne, gegen, um, durch, bis} & Akkusativ & \all{für die Frau, ohne Geld} \\
\all{mit, von, bei, zu, nach, seit, aus, außer, gegenüber} & Dativ & \all{mit der Mutter, seit zwei Jahren, bei der Arbeit} \\
\all{an, auf, hinter, in, neben, über, unter, vor, zwischen} & \textbf{Akk} (mouvement $\rightarrow$ où ?) / \textbf{Dat} (lieu $\rightarrow$ où ?) & \all{in die Schule} (Akk) vs \all{in der Schule} (Dat) \\
\hline
\end{tabular}
\end{center}
\end{propriete}

\courssub{Grammaire 2 : Les temps du passé (Präteritum \& Perfekt)}

\begin{propriete}[Choix et formation des temps du passé]
\textbf{Perfekt (Conversation, lettres, emails) :} Auxiliaire \all{haben} ou \all{sein} (Präsens) + Participe II (ge...t / ge...en) à la fin.
\begin{itemize}
\item Verbes de mouvement (\all{gehen, fahren, fliegen, laufen, kommen, aufstehen}) + changement d'état $\rightarrow$ \all{sein}.
\item Verbes transitifs, réflexifs, modaux, la plupart des verbes $\rightarrow$ \all{haben}.
\item \all{sein} et \all{haben} eux-mêmes font leur Perfekt avec \all{sein} / \all{haben} : \all{ich bin gewesen, ich habe gehabt}.
\end{itemize}

\textbf{Präteritum (Récit, journal, essai, discours indirect) :} Terminaisons faibles (\all{-te, -test, -te, -ten, -tet, -ten}) ou fortes (voyelle modifiée, terminaisons \all{-, -st, -, -en, -t, -en}).

\begin{center}
\begin{tabular}{|l|l|l|l|l|}
\hline
\rowcolor{popgoldL} \textbf{Verbe} & \textbf{Présens (ich)} & \textbf{Präteritum (ich/er/sie/es)} & \textbf{Participe II} & \textbf{Aux. Perfekt} \\
\hline
\all{sein} & \all{bin} & \all{war} & \all{gewesen} & \all{sein} \\
\all{haben} & \all{habe} & \all{hatte} & \all{gehabt} & \all{haben} \\
\all{werden} & \all{werde} & \all{wurde} & \all{geworden} & \all{sein} \\
\all{können} & \all{kann} & \all{konnte} & \all{gekonnt} & \all{haben} \\
\all{müssen} & \all{muss} & \all{musste} & \all{gemusst} & \all{haben} \\
\all{dürfen} & \all{darf} & \all{durfte} & \all{gedurft} & \all{haben} \\
\all{gehen} & \all{gehe} & \all{ging} & \all{gegangen} & \all{sein} \\
\all{kommen} & \all{komme} & \all{kam} & \all{gekommen} & \all{sein} \\
\all{geben} & \all{gebe} & \all{gab} & \all{gegeben} & \all{haben} \\
\all{nehmen} & \all{nehme} & \all{nahm} & \all{genommen} & \all{haben} \\
\all{helfen} & \all{helfe} & \all{half} & \all{geholfen} & \all{haben} (+ Dat) \\
\all{lesen} & \all{lese} & \all{las} & \all{gelesen} & \all{haben} \\
\all{schreiben} & \all{schreibe} & \all{schrieb} & \all{geschrieben} & \all{haben} \\
\all{verdienen} & \all{verdiene} & \all{verdiente} & \all{verdient} & \all{haben} (faible) \\
\all{studieren} & \all{studiere} & \all{studierte} & \all{studiert} & \all{haben} (faible) \\
\hline
\end{tabular}
\end{center}

\textbf{Règle d'or pour le Bac :} À l'oral et dans les lettres $\rightarrow$ \textbf{Perfekt}. Dans un essai, un article, un récit au passé $\rightarrow$ \textbf{Präteritum} (surtout \all{sein, haben, modaux} et verbes forts). Ne pas mélanger au hasard dans un même paragraphe.
\end{propriete}

\begin{popfigure}
\begin{tikzpicture}[
  mindmap, grow cyclic, every node/.style=concept, concept color=popblue,
  level 1/.append style={level distance=4.5cm, sibling angle=90},
  level 2/.append style={level distance=3.2cm, sibling angle=45},
  text=white, font=\small
]
\node {Temps du passé}
  child[concept color=poporange] { node {Perfekt (oral, lettre)}
    child { node {Aux. haben} }
    child { node {Aux. sein (mvt/état)} }
    child { node {Part. II à la fin} }
  }
  child[concept color=popgreen] { node {Präteritum (écrit, récit)}
    child { node {Faibles : -te} }
    child { node {Forts : voyelle modif.} }
    child { node {Modaux : keine Umlaute} }
  }
  child[concept color=poppurple] { node {Marqueurs}
    child { node {\all{gestern, vor 2 Jahren}} }
    child { node {\all{seit 2010 + Präsens}} }
    child { node {\all{damals, früher}} }
  }
  child[concept color=poppink] { node {Pièges}
    child { node {\all{seit} + Présent} }
    child { node {Aux. \all{sein/haben}} }
    child { node {Part. séparable} }
  };
\end{tikzpicture}
\legende{Choix du temps du passé : Perfekt vs Präteritum et marqueurs temporels}
\end{popfigure}

\courssub{Savoir-faire 1 : Rédiger une lettre sur la situation de la femme}

\begin{methode}[Rédiger une lettre (formelle ou amicale)]
\begin{enumerate}
\item \cle{Analyser la consigne} : destinataire (ami ? journal ? autorité ?), registre (\all{du} ou \all{Sie}), longueur exigée, points à traiter.
\item \cle{Poser l'en-tête} : lieu et date à droite (\all{Yaoundé, den 15. März 2025}), puis la formule d'appel : \all{Lieber Paul,} / \all{Liebe Lena,} / \all{Sehr geehrte Damen und Herren,} — \textbf{virgule obligatoire}, et la phrase suivante commence par une \textbf{minuscule}.
\item \cle{Structurer en trois temps} : introduction (motif de la lettre), développement (un argument par paragraphe, avec connecteurs : \all{zuerst, außerdem, jedoch, schließlich}), conclusion (souhait, question, appel).
\item \cle{Formule de politesse} : \all{Viele Grüße} / \all{Liebe Grüße} (amical) ou \all{Mit freundlichen Grüßen} (formel) — \textbf{sans virgule après !}, puis la signature.
\item \cle{Relire} : verbe en 2\textsuperscript{e} position, majuscules aux noms, cas après prépositions (\all{mit dem, für die, in der, seit zwei Jahren}).
\end{enumerate}
\end{methode}

\begin{exoresolu}[Lettre amicale — type Bac]
\textbf{Énoncé.} Votre amie allemande Lena vous demande comment évolue la situation des femmes au Cameroun. Rédigez une lettre d'environ 80 mots : décrivez la situation, donnez un exemple positif et exprimez votre opinion.
\tcblower
\textbf{Raisonnement.} Destinataire = amie $\Rightarrow$ tutoiement, registre amical. Plan : (1) remerciement et motif ; (2) situation générale ; (3) exemple positif (femmes au travail, à l'école) ; (4) opinion personnelle avec \all{Ich finde / Meiner Meinung nach} ; (5) formule de fin. Attention : verbe conjugué en 2\textsuperscript{e} position, \all{Meiner Meinung nach} = datif.

\textbf{Réponse rédigée.}

\all{Yaoundé, den 15. März 2025}

\all{Liebe Lena,}

\all{vielen Dank für deinen Brief. Du möchtest wissen, wie die Situation der Frauen in Kamerun ist.}

\all{Früher mussten viele Frauen nur zu Hause bleiben und die Kinder erziehen. Heute ist das anders: immer mehr Mädchen gehen zur Schule und zur Universität. Viele Frauen arbeiten jetzt als Lehrerinnen, Ärztinnen oder Unternehmerinnen. In Douala zum Beispiel leiten viele Frauen ihre eigenen Geschäfte auf dem Markt.}

\all{Meiner Meinung nach ist das ein großer Fortschritt, aber es gibt noch Probleme: Frauen verdienen oft weniger Geld als Männer. Ich hoffe, dass die Gleichberechtigung bald Realität wird.}

\all{Schreib mir bald zurück!}

\all{Viele Grüße}

\all{Amina}

\textbf{Justifications.} \all{mussten ... bleiben} : Präteritum + infinitif à la fin (verbe modal). \all{immer mehr Mädchen gehen} : sujet après l'adverbe, verbe en 2\textsuperscript{e} position. \all{als Männer} : comparaison avec \all{als} + nominatif. \all{dass die Gleichberechtigung ... wird} : subordonnée, verbe à la fin. \all{Unternehmerinnen} : féminin de \all{Unternehmer}, préféré à \all{Geschäftsfrauen}.

\textbf{Conclusion.} Une lettre réussie respecte le cadre formel (en-tête, appel, formule de fin) et le plan demandé ; chaque point de la consigne doit apparaître explicitement.
\end{exoresolu}

\courssub{Savoir-faire 2 : Rédiger un court essai argumentatif}

\begin{methode}[Structurer un court essai (kurzer Aufsatz)]
\begin{enumerate}
\item \cle{Dégager la problématique} : transformer le sujet en question (\all{Sollen Frauen und Männer dieselben Rechte haben ?}).
\item \cle{Plan en trois parties} : \all{Einleitung} (présentation du thème + thèse), \all{Hauptteil} (arguments \all{dafür} et éventuellement \all{dagegen}, chacun avec un exemple concret), \all{Schluss} (opinion personnelle justifiée + ouverture).
\item \cle{Utiliser des connecteurs} : \all{einerseits ... andererseits} (d'une part ... d'autre part), \all{außerdem} (de plus), \all{jedoch / aber} (cependant), \all{deshalb} (c'est pourquoi), \all{zusammenfassend} (en résumé), \all{außerdem}.
\item \cle{Exprimer l'opinion} : \all{Ich bin der Meinung, dass ...} / \all{Meiner Ansicht nach ...} / \all{Ich finde es wichtig, dass ...} — toujours suivis du verbe à la fin dans la subordonnée (\all{dass}-Satz).
\item \cle{Relire} : cohérence des temps (présent pour les généralités, Präteritum pour faits passés), place du verbe, accords et cas.
\end{enumerate}
\end{methode}

\begin{exoresolu}[Court essai — type Bac]
\textbf{Énoncé.} \all{„Frauen und Männer haben heute dieselben Chancen.“} — Nehmen Sie Stellung (environ 100 mots).
\tcblower
\textbf{Raisonnement.} Consigne : prendre position (\all{Stellung nehmen}). Plan : introduction (le thème), argument 1 (progrès : éducation des filles), argument 2 (limites : salaires, tâches ménagères), conclusion avec opinion nuancée. Connecteurs : \all{einerseits, andererseits, jedoch, meiner Meinung nach}.

\textbf{Réponse rédigée.}

\all{Die Gleichberechtigung von Frauen und Männern ist ein wichtiges Thema unserer Zeit. Aber haben beide wirklich dieselben Chancen?}

\all{Einerseits hat sich viel verändert: Mädchen besuchen heute die Schule und die Universität, und Frauen können fast alle Berufe wählen. Auch in der Politik gibt es immer mehr Frauen.}

\all{Andererseits gibt es noch Ungleichheiten. Frauen verdienen oft weniger als Männer für dieselbe Arbeit, und sie machen den größten Teil der Hausarbeit. In manchen Familien dürfen Mädchen nicht studieren.}

\all{Meiner Meinung nach sind die Chancen also noch nicht gleich, aber die Entwicklung geht in die richtige Richtung. Bildung ist der Schlüssel zur Gleichberechtigung.}

\textbf{Justifications.} \all{hat sich viel verändert} : Perfekt du verbe réfléchi \all{sich verändern}, participe en fin de phrase. \all{weniger als Männer} : comparatif + \all{als}. \all{dürfen ... nicht studieren} : modal au Präsens, infinitif final. \all{der Schlüssel zur Gleichberechtigung} : contraction \all{zu der} $\Rightarrow$ \all{zur} (datif féminin). \all{den größten Teil der Hausarbeit} : Akkusativ (masculin) + Génitif.

\textbf{Conclusion.} Un bon essai annonce sa structure par les connecteurs, illustre chaque argument et se termine par une opinion claire et justifiée.
\end{exoresolu}

\courssub{Savoir-faire 3 : Traduire un thème (français $\rightarrow$ allemand)}

\begin{methode}[Réussir un thème en 5 étapes]
\begin{enumerate}
\item \cle{Lire toute la phrase} et repérer le temps dominant (présent, passé composé $\Rightarrow$ Perfekt, imparfait $\Rightarrow$ Präteritum).
\item \cle{Identifier sujet, verbe, compléments} et choisir les cas : COD $\Rightarrow$ Akkusativ, COI (personne) $\Rightarrow$ Dativ ; attention aux verbes à datif : \all{helfen, danken, gefallen, gehören, folgen, glauben}.
\item \cle{Construire le squelette} : verbe conjugué en 2\textsuperscript{e} position ; dans une subordonnée (\all{dass, weil, obwohl}), verbe à la fin.
\item \cle{Placer les détails} : adverbes de temps en tête ou après le verbe ; particule séparable à la fin (\all{Er steht um 6 Uhr auf}).
\item \cle{Relire} : majuscules aux noms, Umlaute, accord de l'article avec le cas, auxiliaire \all{haben/sein} au Perfekt.
\end{enumerate}
\end{methode}

\begin{exoresolu}[Thème — phrases type Bac]
\textbf{Énoncé.} Traduisez en allemand :
1) « Ma sœur travaille depuis deux ans dans un hôpital à Douala. »
2) « Hier, ma mère m'a aidé à préparer le repas. »
3) « Je pense que les femmes ont les mêmes droits que les hommes. »
\tcblower
\textbf{Phrase 1.} « Depuis deux ans » + présent $\Rightarrow$ \all{seit} + datif + Präsens allemand (pas de passé !). \all{Meine Schwester arbeitet seit zwei Jahren in einem Krankenhaus in Douala.} — \all{seit zwei Jahren} : datif pluriel en \textbf{-n} ; \all{in einem Krankenhaus} : datif (lieu, pas de mouvement $\rightarrow$ \all{wo?}).

\textbf{Phrase 2.} Passé composé $\Rightarrow$ Perfekt ; \all{helfen} exige le \textbf{datif} : \all{meiner Mutter / mir}. \all{Gestern hat meine Mutter mir geholfen, das Essen zuzubereiten.} — Auxiliaire \all{haben} (verbe transitif), participe \all{geholfen} en fin ; \all{zu + Infinitiv} pour « à préparer » ; \all{zuzubereiten} : le \all{zu} s'insère dans le verbe à particule séparable (\all{zubereiten}).

\textbf{Phrase 3.} Subordonnée avec \all{dass} : verbe conjugué à la fin. \all{Ich denke, dass Frauen dieselben Rechte wie Männer haben.} — « les mêmes ... que » = \all{dieselben ... wie} (identité) ; ne pas confondre avec \all{als} (comparaison d'inégalité : \all{mehr als}).

\textbf{Conclusion.} En thème, on traduit des \emph{structures}, pas des mots : cas corrects, verbe à sa place, temps fidèle au français.
\end{exoresolu}

\courssub{Savoir-faire 4 : Traduire une version (allemand $\rightarrow$ français)}

\begin{methode}[Réussir une version en 4 étapes]
\begin{enumerate}
\item \cle{Lire le texte entier} deux fois avant d'écrire ; repérer le thème (ici : la situation de la femme).
\item \cle{Découper chaque phrase} : trouver le verbe conjugué (2\textsuperscript{e} position ou fin de subordonnée), le sujet (nominatif), puis les compléments selon leurs cas.
\item \cle{Traduire les temps fidèlement} : Perfekt $\Rightarrow$ passé composé ; Präteritum $\Rightarrow$ passé simple ou imparfait selon le contexte ; \all{seit} + présent $\Rightarrow$ « depuis » + présent.
\item \cle{Rédiger en bon français} : pas de calque de l'ordre allemand ; reformuler naturellement tout en restant fidèle ; vérifier les faux amis (\all{bekommen} = recevoir, pas « devenir » ; \all{die Ehe} = le mariage ; \all{aktuell} = actuel / d'actualité, « actuellement » = \all{zurzeit} / \all{derzeit}).
\end{enumerate}
\end{methode}

\begin{exoresolu}[Version — texte court type Bac]
\textbf{Énoncé.} Traduisez en français :

\all{Seit zwanzig Jahren kämpfen viele Frauen in Afrika für ihre Rechte. Früher durften sie oft nicht zur Schule gehen, aber heute studieren immer mehr Mädchen an den Universitäten. Trotzdem verdienen Frauen in vielen Ländern weniger Geld als Männer, weil sie schlechtere Stellen bekommen.}
\tcblower
\textbf{Analyse.} \all{Seit zwanzig Jahren kämpfen} : \all{seit} + présent $\Rightarrow$ « depuis vingt ans ... se battent » (présent français). \all{kämpfen für} = se battre pour. \all{durften ... nicht zur Schule gehen} : Präteritum de \all{dürfen} + infinitif $\Rightarrow$ « n'avaient pas le droit d'aller à l'école » (imparfait pour habitude passée). \all{immer mehr Mädchen} = « de plus en plus de filles ». \all{Trotzdem} = « néanmoins / pourtant ». \all{weil sie schlechtere Stellen bekommen} : subordonnée causale, verbe à la fin ; \all{bekommen} = « obtiennent » (faux ami !), \all{die Stelle} = le poste, l'emploi. \all{schlechtere} = comparatif de \all{schlecht} (accusatif pluriel indéfini).

\textbf{Traduction rédigée.}

« Depuis vingt ans, de nombreuses femmes en Afrique se battent pour leurs droits. Autrefois, elles n'avaient souvent pas le droit d'aller à l'école, mais aujourd'hui, de plus en plus de filles font des études universitaires. Pourtant, dans de nombreux pays, les femmes gagnent moins d'argent que les hommes, parce qu'elles obtiennent des postes moins bons. »

\textbf{Conclusion.} Une bonne version est fidèle au sens, respecte les temps et s'exprime en français naturel — jamais mot à mot.
\end{exoresolu}

\courssub{Savoir-faire 5 : Rédiger un dialogue et répondre aux questions de compréhension}

\begin{methode}[Dialogue et questions de compréhension]
\textbf{Pour le dialogue :}
\begin{enumerate}
\item \cle{Répartir les répliques} : au moins 3 échanges (6 répliques) ; chaque réplique courte et naturelle (\all{Was denkst du darüber? — Ich finde, dass ...}).
\item \cle{Varier les verbes de parole} : \all{sagen, fragen, antworten, meinen, erklären} ; après le verbe introducteur, inversion sujet-verbe : \all{„...“, sagt Anna.} ou \all{Anna sagt: „...“}
\item \cle{Utiliser les guillemets allemands} : „ … “ et le tiret de dialogue ou le nom suivi de deux-points.
\item \cle{Marquer l'accord/désaccord} : \all{Ich stimme dir zu. / Das sehe ich anders. / Du hast Recht. / Genau!}
\end{enumerate}
\textbf{Pour les questions de compréhension :}
\begin{enumerate}
\item \cle{Repérer la ligne} du texte qui contient la réponse.
\item \cle{Répondre par une phrase complète} en reformulant (ne pas recopier tout le passage).
\item \cle{Pour les questions d'opinion} (\all{Was meinen Sie?}), justifier avec \all{weil} + verbe final.
\end{enumerate}
\end{methode}

\begin{exoresolu}[Dialogue + compréhension — type Bac]
\textbf{Énoncé A.} Rédigez un court dialogue (6 répliques) entre deux élèves qui discutent du partage des tâches ménagères.

\textbf{Énoncé B.} Répondez en allemand : \all{Warum ist die Gleichberechtigung wichtig für die Entwicklung eines Landes?}
\tcblower
\textbf{Dialogue rédigé.}

\all{Paul: „Hallo Sarah, wer kocht bei euch zu Hause?“}

\all{Sarah: „Meistens meine Mutter, aber mein Vater hilft am Wochenende. Und bei dir?“}

\all{Paul: „Mein Vater kocht nie! Er sagt, das ist Frauenarbeit.“}

\all{Sarah: „Das finde ich unfair. Meiner Meinung nach müssen Männer und Frauen die Hausarbeit teilen.“}

\all{Paul: „Du hast Recht. Wenn beide arbeiten, müssen auch beide im Haushalt helfen.“}

\all{Sarah: „Genau! Gleichberechtigung beginnt zu Hause.“}

\textbf{Justifications.} \all{bei euch} : datif après \all{bei}. \all{Er sagt, das ist Frauenarbeit} : discours indirect libre, verbe en 2\textsuperscript{e} position (ou \all{dass das Frauenarbeit ist}). \all{müssen ... teilen / helfen} : modal + infinitif final. \all{die Hausarbeit teilen} : partager les tâches ménagères (singulier collectif).

\textbf{Réponse à la question B.}

\all{Die Gleichberechtigung ist wichtig für die Entwicklung, weil Frauen die Hälfte der Bevölkerung sind. Wenn Mädchen zur Schule gehen und Frauen arbeiten können, wird das ganze Land reicher und moderner.}

\textbf{Justification.} \all{weil Frauen ... sind} : subordonnée, verbe \all{sind} à la fin ; \all{die Hälfte} = la moitié ; \all{wenn ... gehen und ... können, wird} : double subordonnée, verbe principal (\all{wird}) juste après la virgule (position 2).

\textbf{Conclusion.} Un dialogue authentique sonne naturel (répliques brèves, réactions : \all{Du hast Recht, Genau!}) ; une réponse de compréhension est complète, reformulée et grammaticalement correcte.
\end{exoresolu}

\courssub{Entraînement : Exercices du Bac}

\begin{exercice}[Thème — Phrases à traduire]
Traduisez en allemand :
\begin{enumerate}
\item Depuis dix ans, mon père aide ma mère dans le ménage.
\item Autrefois, les filles n'avaient pas le droit d'étudier à l'université.
\item Bien que les femmes travaillent beaucoup, elles gagnent souvent moins que les hommes.
\item Je suis convaincu que l'égalité des droits est importante pour notre avenir.
\item Hier, nous avons discuté du rôle des femmes dans la société.
\end{enumerate}
\end{exercice}

\begin{exercice}[Version — Texte à traduire]
Traduisez en français :
\all{In vielen afrikanischen Ländern hat sich die Situation der Frauen verbessert. Früher blieben die Mädchen zu Hause, heute gehen sie zur Schule. Dennoch gibt es noch viele Probleme. Zum Beispiel bekommen Frauen oft keinen Kredit von der Bank, weil sie kein Land besitzen. Die Tradition ist oft stärker als das Gesetz.}
\end{exercice}

\begin{exercice}[Production écrite — Lettre formelle]
Vous êtes délégué(e) de votre classe. Écrivez une lettre formelle (env. 100 mots) au Ministre de la Promotion de la Femme et de la Famille pour proposer l'organisation d'une journée « Filles et Sciences » dans les lycées. Structure : en-tête, objet, appel, arguments (2), proposition concrète, formule de politesse.
\end{exercice}

\begin{exercice}[Production écrite — Essai]
Traitez le sujet suivant (env. 120 mots) : \all{„Hausarbeit ist keine Frauensache.“} Nehmen Sie Stellung. (Utilisez : \all{einerseits, andererseits, meiner Meinung nach, weil, dass}).
\end{exercice}

\courssub{Corrigés des exercices d'entraînement}

\begin{corrige}[Exercice 1 — Thème]
\begin{enumerate}
\item \all{Seit zehn Jahren hilft mein Vater meiner Mutter im Haushalt.} (\all{seit} + Dat + Präsens ; \all{helfen} + Dat).
\item \all{Früher durften Mädchen nicht an der Universität studieren.} (Präteritum modal \all{dürfen} + Infinitiv).
\item \all{Obwohl Frauen viel arbeiten, verdienen sie oft weniger als Männer.} (Subordonnée \all{obwohl} verbe à la fin ; phrase principale verbe 2\textsuperscript{e} pos).
\item \all{Ich bin davon überzeugt, dass die Gleichberechtigung für unsere Zukunft wichtig ist.} (\all{überzeugt sein von} + Dat ; \all{dass}-Satz verbe final).
\item \all{Gestern haben wir über die Rolle der Frauen in der Gesellschaft diskutiert.} (Perfekt \all{diskutieren} + \all{haben} ; \all{über} + Akk).
\end{enumerate}
\end{corrige}

\begin{corrige}[Exercice 2 — Version]
« Dans de nombreux pays africains, la situation des femmes s'est améliorée. Autrefois, les filles restaient à la maison, aujourd'hui elles vont à l'école. Néanmoins, il existe encore de nombreux problèmes. Par exemple, les femmes n'obtiennent souvent pas de crédit de la banque, parce qu'elles ne possèdent pas de terre. La tradition est souvent plus forte que la loi. »
\emph{Notes :} \all{hat sich verbessert} (Perfekt réfléchi) $\rightarrow$ s'est améliorée / s'améliore. \all{blieben} (Präteritum) $\rightarrow$ restaient (imparfait habitude). \all{bekommen} $\rightarrow$ obtiennent (faux ami). \all{besitzen} $\rightarrow$ possèdent. \all{stärker als} $\rightarrow$ plus forte que.
\end{corrige}

\begin{corrige}[Exercice 3 — Lettre formelle (proposition)]
\all{Yaoundé, den 20. Oktober 2025}

\all{An das Ministerium für die Förderung der Frau und der Familie}
\all{z. H. Herrn Minister}
\all{Yaoundé}

\all{Betreff: Vorschlag für einen Tag „Mädchen und Naturwissenschaften“}

\all{Sehr geehrte Damen und Herren,}

\all{ich schreibe Ihnen als Klassensprecher des Lycée Bilingue de Yaoundé. Meiner Meinung nach ist es sehr wichtig, Mädchen für technische und wissenschaftliche Berufe zu begeistern.}

\all{Einerseits zeigen Mädchen in der Schule oft gute Leistungen in Mathematik und Physik. Andererseits wählen sie später selten Studiengänge wie Informatik oder Ingenieurwesen, weil ihnen Vorbilder fehlen.}

\all{Deshalb schlage ich vor, einen „Tag der Mädchen in den Naturwissenschaften“ zu organisieren. An diesem Tag könnten Studentinnen und Ingenieurinnen über ihre Berufe berichten und Experimente zeigen. Das würde das Selbstvertrauen der Schülerinnen stärken.}

\all{Ich würde mich freuen, wenn Sie diesen Vorschlag prüfen könnten.}

\all{Mit freundlichen Grüßen}

\all{Jean Dupont}
\all{Klassensprecher}
\end{corrige}

\begin{corrige}[Exercice 4 — Essai]
\all{„Hausarbeit ist keine Frauensache.“ Das ist ein wichtiges Statement. Nehmen wir Stellung.}

\all{Einerseits war es früher normal, dass nur die Frau kocht, putzt und die Kinder betreut. Der Mann ging arbeiten und brachte das Geld nach Hause. Diese Rollenverteilung war klar, aber ungerecht.}

\all{Andererseits hat sich die Welt verändert. Heute arbeiten beide Partner oft Vollzeit. Wenn beide arbeiten, müssen beide auch die Hausarbeit machen. Das ist nur fair. Außerdem zeigen Studien: Wenn Männer im Haushalt helfen, sind die Familien glücklicher.}

\all{Meiner Meinung nach ist Hausarbeit keine Frauensache, sondern eine gemeinsame Aufgabe. Wer das nicht versteht, lebt in der Vergangenheit. Gleichberechtigung beginnt zu Hause — jeden Tag.}
\end{corrige}

\begin{attention}[Les 5 erreurs qui coûtent des points au Bac]
\begin{enumerate}
\item \cle{Oublier la majuscule aux noms} : \all{die frau, das haus} sont fautifs — écrire \all{die Frau, das Haus}. Tout nom, même dans un adjectif nominalisé (\all{das Essen, das Lesen}), prend la majuscule.
\item \cle{Calquer l'ordre français} : « Je pense que les femmes ont... » ne devient jamais \all{Ich denke, dass Frauen haben...} — dans la subordonnée, le verbe conjugué va \textbf{à la fin} : \all{dass Frauen ... haben}.
\item \cle{Se tromper de cas après les prépositions} : \all{seit, mit, von, bei, zu, nach, aus, außer, gegenüber} exigent le \textbf{datif} (\all{seit zwei Jahren, mit meiner Mutter, bei der Arbeit}) ; \all{für, ohne, gegen, um, durch, bis} exigent l'\textbf{accusatif}. Prépositions doubles (\all{an, in, auf...}) : Akk (mouvement) / Dat (lieu).
\item \cle{Tomber dans les faux amis} : \all{bekommen} = recevoir (devenir = \all{werden}) ; \all{die Ehe} = le mariage (pas l'huile) ; \all{aktuell} = actuel / d'actualité (« actuellement » = \all{zurzeit}) ; \all{konsequent} = cohérent / conséquent (pas « conséquent » au sens de « donc ») ; \all{das Gymnasium} = le lycée (pas la salle de sport).
\item \cle{Maltraiter les verbes à particule et le Perfekt} : \all{Er aufsteht um 6 Uhr} est faux $\Rightarrow$ \all{Er steht um 6 Uhr auf} ; au Perfekt, choisir le bon auxiliaire : \all{sie ist gegangen} (mouvement), \all{sie ist aufgestanden} (changement d'état), mais \all{sie hat gearbeitet} (verbe transitif/activité), \all{sie hat geschlafen} (verbe intransitif d'état $\rightarrow$ haben en allemand standard, sein en autrichien — au Bac camerounais : \all{haben} pour \all{schlafen}).
\end{enumerate}
Réflexe final : relire chaque phrase en vérifiant \textbf{verbe (place) — cas — majuscules — Umlaute — ß}.
\end{attention}

\begin{savaistu}[Femmes au Cameroun et dans l'espace germanophone]
\textbf{Cameroun :} La Constitution de 1996 garantit l'égalité. Des figures historiques : \textbf{Marie-Thérèse Assiga Ahanda} (1\textsuperscript{re} femme médecin, 1\textsuperscript{re} femme députée), \textbf{Delphine Zanga Tsogo} (1\textsuperscript{re} femme ministre en 1975). Aujourd'hui : \textbf{Josephine Ndagnou} (PDG Cameroon Airlines), femmes au Marché Central (Yaoundé) et Marché des Fleurs (Douala) moteurs de l'économie informelle. Défis : accès au foncier, mariages précoces (loi 2016 : 18 ans), violences basées sur le genre.

\textbf{Allemagne :} \all{Artikel 3 Grundgesetz} (1949) : \all{Männer und Frauen sind gleichberechtigt.} \textbf{Angela Merkel} (Chancelière 2005–2021). \textbf{Elly Heuss-Knapp} (fondatrice Müttergenesungswerk). \textbf{Alice Salomon} (pionnière travail social). Problème actuel : \all{Gender Pay Gap} (\~18\%), quota femmes dans les conseils d'administration (\all{Frauenquote}).

\textbf{Autriche / Suisse :} Droit de vote femmes : Autriche 1918, Suisse 1971 (niveau fédéral), 1990 (dernier canton Appenzell Rhodes-Intérieures). \textbf{Simone de Beauvoir} (française, mais référence incontournable pour \all{Le Deuxième Sexe} lu en cours d'allemand).
\end{savaistu}

\newpage

\coursec{Exercices}

\courssub{Je vérifie mes connaissances}

\begin{exercice}[QCM : grammaire et vocabulaire]
Entoure la bonne réponse.
\begin{enumerate}
\item \all{Gestern hat meine Mutter mir eine Geschichte ...} \quad a) \all{erzählt} \quad b) \all{erzählen} \quad c) \all{erzählte}
\item \all{Die Frauen kämpfen ... ihre Rechte.} \quad a) \all{für} \quad b) \all{über} \quad c) \all{mit}
\item Quel mot signifie « l'égalité » ? \quad a) \all{die Gleichheit} \quad b) \all{die Gleichberechtigung} \quad c) \all{die Gerechtigkeit}
\item \all{Meine Schwester arbeitet, ... sie drei Kinder hat.} \quad a) \all{denn} \quad b) \all{obwohl} \quad c) \all{deswegen}
\item Dans une lettre formelle, on termine par : \quad a) \all{Liebe Grüße} \quad b) \all{Mit freundlichen Grüßen} \quad c) \all{Tschüss}
\end{enumerate}
\end{exercice}

\begin{exercice}[Vrai ou faux ? Justifie ta réponse]
Réponds par \all{richtig} ou \all{falsch} et corrige la phrase si elle est fausse.
\begin{enumerate}
\item \all{Weil sie krank war, sie ist zu Hause geblieben.}
\item \all{Mein Großvater hat früher auf dem Land gewohnt.}
\item \all{Ich habe meine Freundin gestern im Markt gesehen.}
\item \all{Die Männer sollen auch die Hausarbeit machen.}
\item \all{Nach dem Abitur will meine Schwester Medizin studieren.}
\end{enumerate}
\end{exercice}

\begin{exercice}[Vocabulaire : la famille et la société]
Trouve le mot allemand (avec son article) qui correspond à chaque définition française.
\begin{enumerate}
\item le mariage \item le congé parental \item les droits de la femme \item le métier, la profession \item l'éducation, la formation \item le foyer, le ménage
\end{enumerate}
\end{exercice}

\begin{exercice}[Conjugaison : Präteritum ou Perfekt]
Conjugue les verbes entre parenthèses au temps qui convient.
\begin{enumerate}
\item \all{Früher (haben) die Frauen in unserem Dorf keine Schule besuchen dürfen.}
\item \all{Gestern (erzählen) meine Großmutter mir von ihrer Jugend.}
\item \all{Als sie jung war, (arbeiten) sie auf dem Feld und (kochen) für die ganze Familie.}
\item \all{Meine Tante (studieren) in Deutschland und (werden) Ärztin.}
\item \all{Letztes Jahr (bekommen) meine Cousine ihr erstes Kind.}
\end{enumerate}
\end{exercice}

\courssub{Je m'entraîne}

\begin{exercice}[Thème : traduis en allemand]
\begin{enumerate}
\item Hier, ma grand-mère m'a raconté sa jeunesse.
\item Les femmes doivent avoir les mêmes droits que les hommes.
\item Mon père fait la vaisselle parce que ma mère travaille beaucoup.
\item Ma sœur aînée a fait des études de médecine à Douala.
\item Bien que la situation ait changé, beaucoup de filles ne vont toujours pas à l'école.
\end{enumerate}
\end{exercice}

\begin{exercice}[Remets les mots dans l'ordre]
Reconstitue des phrases correctes (déclaratives ou subordonnées). Attention à la place du verbe.
\begin{enumerate}
\item \all{hat / gestern / meine Mutter / auf dem Markt / Gemüse / gekauft}
\item \all{die Frauen / weil / verdienen / weniger Geld / oft / als die Männer}
\item \all{dass / ich / finde / die Männer / helfen / sollen / im Haushalt / auch}
\item \all{meine Schwester / obwohl / müde / war / hat / die Hausarbeit / gemacht}
\item \all{wenn / heiratet / sie / will / weiter / studieren / sie}
\end{enumerate}
\end{exercice}

\begin{exercice}[Texte à trous : articles, cas et temps]
Complète le texte avec les articles au bon cas et les verbes au \all{Präteritum} ou au \all{Perfekt} : \all{sein, können, arbeiten, helfen, sagen, gehen}.

\begin{textebox}[Früher und heute]
Früher (1) \dots\ d\dots\ Leben der Frauen in unserem Dorf sehr schwer. Meine Großmutter (2) \dots\ nicht zur Schule (3) \dots\ , weil ihre Eltern kein Geld hatten. Sie (4) \dots\ jeden Tag auf d\dots\ Feld und (5) \dots\ ihren Eltern im Haushalt. Heute, (6) \dots\ sie oft, ist alles anders : d\dots\ Mädchen (7) \dots\ zur Schule gehen und später einen Beruf lernen.
\end{textebox}
\end{exercice}

\begin{exercice}[Appariements : expressions de l'essai]
Relie chaque expression allemande à son équivalent français.
\begin{center}
\begin{tabularx}{\linewidth}{|X|X|}
\hline
\rowcolor{popblueL} \textbf{Deutsch} & \textbf{Français} \\
\hline
\all{1. einerseits ... andererseits} & a. à mon avis \\
\hline
\all{2. meiner Meinung nach} & b. d'une part ... d'autre part \\
\hline
\all{3. ein wichtiges Argument dafür ist} & c. en conclusion \\
\hline
\all{4. zum Schluss möchte ich sagen} & d. un argument important en faveur est \\
\hline
\all{5. viele Menschen sind der Ansicht} & e. beaucoup de gens pensent \\
\hline
\end{tabularx}
\end{center}
\end{exercice}

\courssub{Je me prépare au Bac}

\begin{exercice}[Compréhension de l'écrit et questions de langue]
Lis le texte suivant, puis réponds aux questions.

\begin{textebox}[Der Elternurlaub in Deutschland]
In Deutschland können Väter und Mütter nach der Geburt eines Kindes einen Elternurlaub nehmen. Das heißt : Sie bleiben zu Hause und bekommen einen Teil ihres Gehalts vom Staat. Früher waren es fast immer die Mütter, die zu Hause blieben. Heute nehmen immer mehr Väter den Elternurlaub, weil sie Zeit mit ihren Kindern verbringen wollen. Viele junge Männer sagen : „Ich möchte mein Kind aufwachsen sehen.“ Aber es gibt auch Probleme : Manche Chefs finden es nicht gut, wenn ein Mann lange nicht arbeitet. Trotzdem glauben viele Familien, dass der Elternurlaub die Gleichberechtigung von Mann und Frau stärkt.
\end{textebox}

\textbf{A. Compréhension} (réponds en allemand, avec des phrases complètes)
\begin{enumerate}
\item Was ist der Elternurlaub ?
\item Wer blieb früher meistens zu Hause ?
\item Warum nehmen heute mehr Väter den Elternurlaub ?
\item Welches Problem wird im Text genannt ?
\end{enumerate}

\textbf{B. Langue}
\begin{enumerate}
\item Trouve dans le texte : un verbe au \all{Perfekt}, une subordonnée avec \all{weil}, un adjectif possessif.
\item Mets à la forme négative : \all{Viele Männer nehmen den Elternurlaub.}
\item Traduis en français la dernière phrase du texte.
\end{enumerate}
\end{exercice}

\begin{exercice}[Version et thème]
\textbf{A. Version.} Traduis en français le texte suivant (environ 6 lignes).

\begin{textebox}[Ein neues Leben]
Amina ist zwanzig Jahre alt und kommt aus Douala. Letztes Jahr hat sie ihr Abitur gemacht, und jetzt möchte sie Informatik studieren. Ihre Mutter ist stolz auf sie, aber ihr Onkel sagt : „Ein Mädchen soll lieber früh heiraten.“ Amina antwortet : „Ich liebe meine Familie, aber ich habe auch Träume. Bildung ist wichtig für Mädchen und für Jungen.“
\end{textebox}

\textbf{B. Thème.} Traduis en allemand.
\begin{enumerate}
\item Ma tante a travaillé toute sa vie pour ses enfants.
\item Je pense que les hommes doivent aussi aider à la maison.
\item Quand j'étais petit, ma mère m'a lu beaucoup d'histoires.
\end{enumerate}
\end{exercice}

\begin{exercice}[Expression écrite : essai et dialogue]
Choisis \textbf{un} des deux sujets.

\textbf{Sujet 1 — Essai (10 à 12 lignes).} \all{Sollen die Männer auch Hausarbeit machen?} Rédige un court essai structuré : \all{Einleitung} (présentation du sujet), \all{Hauptteil} (au moins deux arguments \all{dafür} et un argument \all{dagegen}), \all{Schluss} (ton opinion personnelle). Utilise au moins trois connecteurs vus dans le chapitre (\all{einerseits, andererseits, meiner Meinung nach, deswegen, zum Schluss}).

\textbf{Sujet 2 — Dialogue (8 répliques).} Écris un dialogue entre un père traditionnel et sa fille de 18 ans qui veut faire des études à l'université de Yaoundé. Le père préfère un mariage précoce ; la fille défend son projet avec des arguments. Termine le dialogue par un compromis ou une conclusion claire. Soigne les formules de politesse et l'ordre des mots.
\end{exercice}

\begin{exercice}[Corrige la lettre]
La lettre suivante d'un élève contient \textbf{10 erreurs} (ordre des mots, cas, majuscules, temps des verbes). Souligne chaque erreur et écris la correction.

\begin{textebox}[Brief an einen Freund]
Lieber Thomas,\\
danke für dein Brief. Ich habe mich sehr gefreut. Letzte Woche ich habe mein Onkel in Douala besucht. Er hat mir gesagt, dass frauen heute mehr Rechte haben als früher. Meine Schwester will Ärztin werden, weil sie gerne menschen hilft. Mein Vater findet das gut, aber meine Großmutter sagt, dass eine Frau soll zu Hause bleiben. Ich finde, dass jeder Mensch seine Beruf wählen muss. Was denkst du ? Schreib mir bald !\\
Viele Grüße\\
Paul
\end{textebox}
\end{exercice}

\newpage

\coursec{Corrigés des exercices}

\begin{corrige}[Exercice 1 — QCM : grammaire et vocabulaire]
\begin{enumerate}
\item \textbf{a)} \all{Gestern hat meine Mutter mir eine Geschichte erzählt.} — Au \all{Perfekt}, le participe passé (\all{erzählt}) se place en fin de phrase ; \all{erzählte} serait le \all{Präteritum} (sans auxiliaire \all{hat}), \all{erzählen} l'infinitif.
\item \textbf{a)} \all{Die Frauen kämpfen für ihre Rechte.} — Le verbe \all{kämpfen} se construit avec la préposition \all{für} (+ accusatif) : « se battre pour ». \all{kämpfen gegen} = lutter contre.
\item \textbf{b)} \all{die Gleichberechtigung} = l'égalité des droits (entre hommes et femmes). \all{die Gleichheit} = l'égalité (mathématique, abstraite) ; \all{die Gerechtigkeit} = la justice.
\item \textbf{b)} \all{Meine Schwester arbeitet, obwohl sie drei Kinder hat.} — « Ma sœur travaille, bien qu'elle ait trois enfants » : opposition, donc \all{obwohl} (verbe en fin de subordonnée). \all{denn} = car ; \all{deswegen} = c'est pourquoi.
\item \textbf{b)} \all{Mit freundlichen Grüßen} — formule de politesse finale d'une lettre formelle. \all{Liebe Grüße} et \all{Tschüss} sont familiers.
\end{enumerate}
\end{corrige}

\begin{corrige}[Exercice 2 — Vrai ou faux ?]
\begin{enumerate}
\item \all{falsch}. Après une subordonnée en tête de phrase, le verbe de la principale vient en premier (inversion) : \all{Weil sie krank war, ist sie zu Hause geblieben.} (« Comme elle était malade, elle est restée à la maison. »)
\item \all{richtig}. \all{Mein Großvater hat früher auf dem Land gewohnt.} — \all{Perfekt} correct, \all{auf dem Land} = à la campagne (datif de lieu).
\item \all{falsch}. \all{der Markt} est masculin : \all{Ich habe meine Freundin gestern auf dem Markt gesehen.} (on dit \all{auf dem Markt}, pas \all{im Markt}).
\item \all{richtig}. \all{Die Männer sollen auch die Hausarbeit machen.} — « Les hommes devraient aussi faire le ménage » ; verbe modal en position 2, infinitif en fin.
\item \all{richtig}. \all{Nach dem Abitur will meine Schwester Medizin studieren.} — complément en tête, donc inversion sujet-verbe, correctement appliquée.
\end{enumerate}
\end{corrige}

\begin{corrige}[Exercice 3 — Vocabulaire : la famille et la société]
\begin{enumerate}
\item le mariage : \all{die Ehe, -n} (ou \all{die Heirat, -en})
\item le congé parental : \all{der Elternurlaub} (on trouve aussi \all{die Elternzeit})
\item les droits de la femme : \all{die Frauenrechte} (pluriel)
\item le métier, la profession : \all{der Beruf, -e}
\item l'éducation, la formation : \all{die Ausbildung, -en} (formation) ; \all{die Bildung} (éducation, instruction)
\item le foyer, le ménage : \all{der Haushalt, -e}
\end{enumerate}
\textbf{Point de vigilance :} en allemand, tout nom prend une majuscule et doit être appris avec son article : \all{der Beruf}, \all{die Ehe}, \all{der Haushalt}.
\end{corrige}

\begin{corrige}[Exercice 4 — Conjugaison : Präteritum ou Perfekt]
\begin{enumerate}
\item \all{Früher durften die Frauen in unserem Dorf keine Schule besuchen.} — \all{Präteritum} du modal \all{dürfen} (\all{durften}) : récit d'une situation passée révolue ; à l'écrit, les modaux s'emploient de préférence au \all{Präteritum}.
\item \all{Gestern hat meine Großmutter mir von ihrer Jugend erzählt.} (ou \all{Gestern erzählte meine Großmutter mir von ihrer Jugend.}) — \all{Perfekt} : événement ponctuel daté (\all{gestern}) dans un récit oral/lettre.
\item \all{Als sie jung war, arbeitete sie auf dem Feld und kochte für die ganze Familie.} — \all{Präteritum} : description d'habitudes passées dans un récit (\all{als sie jung war}).
\item \all{Meine Tante hat in Deutschland studiert und ist Ärztin geworden.} — \all{Perfekt} ; attention : \all{werden} construit son \all{Perfekt} avec \all{sein} : \all{ist geworden}.
\item \all{Letztes Jahr hat meine Cousine ihr erstes Kind bekommen.} — \all{Perfekt} : événement ponctuel ; participe \all{bekommen} (verbe à préfixe inséparable \all{be-}, sans \all{ge-}).
\end{enumerate}
\end{corrige}

\begin{corrige}[Exercice 5 — Thème : traduis en allemand]
\begin{enumerate}
\item \all{Gestern hat meine Großmutter mir von ihrer Jugend erzählt.} — \all{Perfekt} (événement ponctuel), \all{jdm. von etwas erzählen} = raconter qqch. à qqn.
\item \all{Die Frauen müssen die gleichen Rechte haben wie die Männer.} — comparaison d'égalité : \all{gleich ... wie} ; modal \all{müssen} en position 2, infinitif en fin.
\item \all{Mein Vater spült ab, weil meine Mutter viel arbeitet.} (ou \all{Mein Vater macht den Abwasch, weil meine Mutter viel arbeitet.}) — « faire la vaisselle » = \all{abspülen} (verbe à particule séparable : \all{spült ... ab}) ; \all{weil} renvoie le verbe en fin.
\item \all{Meine ältere Schwester hat in Douala Medizin studiert.} — \all{Perfekt} ; « faire des études de médecine » = \all{Medizin studieren}.
\item \all{Obwohl sich die Situation geändert hat, gehen viele Mädchen immer noch nicht zur Schule.} — subordonnée \all{obwohl} en tête : verbe conjugué en fin de subordonnée, puis inversion dans la principale (\all{gehen} avant le sujet \all{viele Mädchen}).
\end{enumerate}
\textbf{Points de vigilance :} majuscules aux noms (\all{Großmutter, Rechte, Männer, Situation, Schule}) ; place du verbe ; \all{zur Schule gehen} (et non \all{in die Schule} pour la fréquentation scolaire).
\end{corrige}

\begin{corrige}[Exercice 6 — Remets les mots dans l'ordre]
\begin{enumerate}
\item \all{Meine Mutter hat gestern auf dem Markt Gemüse gekauft.} — phrase déclarative : sujet + verbe conjugué en position 2, participe \all{gekauft} en fin.
\item \all{Weil die Frauen oft weniger Geld verdienen als die Männer.} — subordonnée \all{weil} : verbe conjugué \all{verdienen} en fin ; comparaison \all{weniger ... als}.
\item \all{Ich finde, dass die Männer auch im Haushalt helfen sollen.} — dans la subordonnée \all{dass} avec modal, l'ordre final est : infinitif + modal (\all{helfen sollen}).
\item \all{Obwohl meine Schwester müde war, hat sie die Hausarbeit gemacht.} — subordonnée en tête (verbe \all{war} en fin), puis inversion : \all{hat} en premier dans la principale.
\item \all{Wenn sie heiratet, will sie weiter studieren.} — même structure : \all{heiratet} en fin de subordonnée, \all{will} en tête de principale.
\end{enumerate}
\end{corrige}

\begin{corrige}[Exercice 7 — Texte à trous : articles, cas et temps]
\begin{textebox}[Früher und heute — corrigé]
Früher (1) \all{war} \textbf{das} Leben der Frauen in unserem Dorf sehr schwer. Meine Großmutter (2) \all{konnte} nicht zur Schule (3) \all{gehen}, weil ihre Eltern kein Geld hatten. Sie (4) \all{arbeitete} jeden Tag auf \textbf{dem} Feld und (5) \all{half} ihren Eltern im Haushalt. Heute, (6) \all{sagt} sie oft, ist alles anders : \textbf{die} Mädchen (7) \all{können} zur Schule gehen und später einen Beruf lernen.
\end{textebox}
\textbf{Justifications :}
\begin{enumerate}
\item \all{war} : \all{Präteritum} de \all{sein} (récit au passé) ; \all{das Leben} : nominatif neutre (sujet).
\item \all{konnte} : \all{Präteritum} de \all{können} ; (3) \all{gehen} : infinitif en fin après le modal.
\item \all{arbeitete} : \all{Präteritum} (habitude passée) ; \all{auf dem Feld} : datif (lieu, \all{das Feld}).
\item \all{half} : \all{Präteritum} fort de \all{helfen} ; attention, \all{helfen} se construit avec le datif : \all{ihren Eltern}.
\item \all{sagt} : présent (parole actuelle de la grand-mère).
\item \all{die Mädchen} : nominatif pluriel ; \all{können} : présent (situation d'aujourd'hui).
\end{enumerate}
\end{corrige}

\begin{corrige}[Exercice 8 — Appariements : expressions de l'essai]
\begin{center}
\begin{tabularx}{\linewidth}{|X|X|}
\hline
\rowcolor{popblueL} \textbf{Deutsch} & \textbf{Français} \\
\hline
\all{1. einerseits ... andererseits} & \textbf{b.} d'une part ... d'autre part \\
\hline
\all{2. meiner Meinung nach} & \textbf{a.} à mon avis \\
\hline
\all{3. ein wichtiges Argument dafür ist} & \textbf{d.} un argument important en faveur est \\
\hline
\all{4. zum Schluss möchte ich sagen} & \textbf{c.} en conclusion \\
\hline
\all{5. viele Menschen sind der Ansicht} & \textbf{e.} beaucoup de gens pensent \\
\hline
\end{tabularx}
\end{center}
\textbf{Remarque :} \all{meiner Meinung nach} suit le nom au génitif-datif (la préposition \all{nach} est ici postposée) ; \all{der Ansicht sein} = être d'avis.
\end{corrige}

\begin{corrige}[Exercice 9 — Compréhension de l'écrit et questions de langue]
\textbf{A. Compréhension — réponses modèles :}
\begin{enumerate}
\item \all{Der Elternurlaub ist eine Zeit nach der Geburt eines Kindes, in der Väter und Mütter zu Hause bleiben und einen Teil ihres Gehalts vom Staat bekommen.}
\item \all{Früher blieben fast immer die Mütter zu Hause.}
\item \all{Heute nehmen mehr Väter den Elternurlaub, weil sie Zeit mit ihren Kindern verbringen wollen.}
\item \all{Ein Problem ist, dass manche Chefs es nicht gut finden, wenn ein Mann lange nicht arbeitet.}
\end{enumerate}
\textbf{B. Langue :}
\begin{enumerate}
\item Verbe au \all{Perfekt} : le texte est essentiellement au présent ; on retiendra la forme de participe dans \all{„Ich möchte mein Kind aufwachsen sehen.“} — ici \all{aufwachsen} est un infinitif dépendant de \all{sehen}. La subordonnée avec \all{weil} : \all{weil sie Zeit mit ihren Kindern verbringen wollen} (verbe en fin : \all{verbringen wollen}). Adjectif possessif : \all{ihres Gehalts}, \all{ihren Kindern}, \all{mein Kind}.
\item \all{Viele Männer nehmen den Elternurlaub nicht.} (ou \all{Nicht viele Männer nehmen den Elternurlaub.})
\item Traduction : « Néanmoins, beaucoup de familles croient que le congé parental renforce l'égalité des droits entre l'homme et la femme. »
\end{enumerate}
\textbf{Barème indicatif (Bac) :} A : 4 × 1 pt (réponse complète en allemand correct) ; B1 : 1 pt ; B2 : 1 pt ; B3 : 2 pts. \textbf{Points de vigilance :} répondre par des phrases complètes ; dans la négation, \all{nicht} se place en fin de phrase quand il porte sur tout l'énoncé.
\end{corrige}

\begin{corrige}[Exercice 10 — Version et thème]
\textbf{A. Version — traduction modèle :}

« Amina a vingt ans et vient de Douala. L'année dernière, elle a passé son baccalauréat, et maintenant elle voudrait étudier l'informatique. Sa mère est fière d'elle, mais son oncle dit : “Une fille devrait plutôt se marier tôt.” Amina répond : “J'aime ma famille, mais j'ai aussi des rêves. L'instruction est importante pour les filles et pour les garçons.” »

\textbf{Points de vigilance :} \all{das Abitur machen} = passer le baccalauréat ; \all{stolz auf jdn. sein} = être fier de qqn ; \all{früh heiraten} = se marier tôt ; \all{die Bildung} = l'instruction, l'éducation.

\textbf{B. Thème — traductions modèles :}
\begin{enumerate}
\item \all{Meine Tante hat ihr ganzes Leben für ihre Kinder gearbeitet.} — \all{Perfekt} ; \all{das ganze Leben} (accusatif de durée).
\item \all{Ich finde, dass die Männer auch im Haushalt helfen sollen.} (ou \all{Meiner Meinung nach sollen die Männer auch zu Hause helfen.}) — « aider à la maison » = \all{im Haushalt helfen} ; dans la subordonnée : \all{helfen sollen} en fin.
\item \all{Als ich klein war, hat meine Mutter mir viele Geschichten vorgelesen.} — \all{als} (+ \all{Präteritum}) pour un fait passé unique ; « lire des histoires à qqn » = \all{jdm. Geschichten vorlesen} (verbe à particule séparable : \all{vorgelesen}).
\end{enumerate}
\textbf{Barème indicatif :} version 6 pts (fidélité 3 pts, qualité du français 3 pts) ; thème 3 × 2 pts (grammaire 1 pt, vocabulaire 1 pt).
\end{corrige}

\begin{corrige}[Exercice 11 — Expression écrite : essai et dialogue]
\textbf{Sujet 1 — Essai : production modèle}

\begin{textebox}[Sollen die Männer auch Hausarbeit machen?]
In vielen Familien ist die Hausarbeit noch immer die Aufgabe der Frauen. Aber ist das gerecht? Meiner Meinung nach ist diese Frage heute sehr wichtig.\\
Einerseits arbeiten viele Frauen genauso viel wie die Männer. Wenn die Frau nach der Arbeit allein kochen, putzen und die Kinder betreuen muss, ist das ungerecht. Ein wichtiges Argument dafür ist auch, dass die Kinder so ein gutes Beispiel sehen : Sie lernen, dass Mann und Frau die gleichen Pflichten haben. Andererseits sagen manche Menschen, dass die Männer nach einem langen Arbeitstag müde sind und Ruhe brauchen.\\
Zum Schluss möchte ich sagen : Die Hausarbeit ist eine Aufgabe für alle. Deswegen sollen Väter und Mütter die Arbeit zu Hause teilen. So wird die Familie stärker und die Gleichberechtigung wird Wirklichkeit.
\end{textebox}

\textbf{Barème indicatif (10 pts) :} respect du sujet et de la structure (Einleitung–Hauptteil–Schluss) 2 pts ; arguments pertinents 3 pts ; connecteurs (au moins 3) 2 pts ; correction grammaticale (ordre des mots, cas, temps) 2 pts ; vocabulaire et orthographe 1 pt.

\textbf{Sujet 2 — Dialogue : production modèle}

\begin{textebox}[Ein Gespräch zwischen Vater und Tochter]
\textbf{Vater :} Miriam, du hast jetzt dein Abitur. Ich denke, es ist Zeit zu heiraten. Der Sohn von Onkel Bekono ist ein guter Mann.\\
\textbf{Miriam :} Papa, ich respektiere deine Meinung, aber ich möchte zuerst in Yaoundé Medizin studieren.\\
\textbf{Vater :} Studieren ? Eine Frau gehört zur Familie und zu den Kindern !\\
\textbf{Miriam :} Früher war das so, aber heute können Frauen Ärztinnen oder Lehrerinnen werden. Ich möchte später auch meiner Familie helfen.\\
\textbf{Vater :} Und wer bezahlt die Universität ? Das ist sehr teuer.\\
\textbf{Miriam :} Ich habe ein Stipendium bekommen, weil ich gute Noten habe.\\
\textbf{Vater :} Hmm ... Und deine Zukunft als Mutter ?\\
\textbf{Miriam :} Ich kann später heiraten, Papa. Bildung und Familie schließen sich nicht aus. Versprich mir, dass du darüber nachdenkst !\\
\textbf{Vater :} Gut, meine Tochter. Wir sprechen morgen mit deiner Mutter. Vielleicht hast du recht.
\end{textebox}

\textbf{Barème indicatif (10 pts) :} 8 répliques équilibrées 2 pts ; arguments des deux points de vue 3 pts ; compromis/conclusion 1 pt ; correction grammaticale 3 pts ; formules de politesse et registre 1 pt.

\textbf{Points de vigilance :} verbe en position 2 dans les répliques déclaratives ; \all{zuerst} (d'abord) ≠ \all{zuerst} mal orthographié ; \all{rechthaben} → \all{du hast recht} (particule séparable) ; éviter le calque du français « je suis d'accord » : \all{ich bin einverstanden}.
\end{corrige}

\begin{corrige}[Exercice 12 — Corrige la lettre]
Les 10 erreurs soulignées, avec correction :

\begin{textebox}[Brief an einen Freund — corrigé]
Lieber Thomas,\\
danke für \textbf{deinen} Brief (1). Ich habe mich sehr gefreut. Letzte Woche \textbf{habe ich} (2) mein\textbf{en} (3) Onkel in Douala besucht. Er hat mir gesagt, dass \textbf{Frauen} (4) heute mehr Rechte haben als früher. Meine Schwester will Ärztin werden, weil sie gerne \textbf{Menschen} (5) hilft. Mein Vater findet das gut, aber meine Großmutter sagt, dass eine Frau zu Hause bleiben \textbf{soll} (6). Ich finde, dass jeder Mensch seine\textbf{n} (7) Beruf wählen muss. Was denkst du ? Schreib mir bald !\\
Viele Grüße\\
Paul
\end{textebox}

\textbf{Détail des corrections :}
\begin{enumerate}
\item \all{für dein Brief} → \all{für deinen Brief} : \all{der Brief} est masculin, accusatif après \all{für} → \all{deinen}.
\item \all{Letzte Woche ich habe} → \all{Letzte Woche habe ich} : complément en tête de phrase, inversion sujet-verbe obligatoire.
\item \all{mein Onkel ... besucht} → \all{meinen Onkel} : \all{der Onkel}, accusatif (objet direct de \all{besuchen}).
\item \all{frauen} → \all{Frauen} : majuscule obligatoire à tous les noms.
\item \all{menschen} → \all{Menschen} : même règle de majuscule.
\item \all{dass eine Frau soll zu Hause bleiben} → \all{dass eine Frau zu Hause bleiben soll} : dans la subordonnée \all{dass}, le verbe conjugué (\all{soll}) va en fin, après l'infinitif.
\item \all{seine Beruf} → \all{seinen Beruf} : \all{der Beruf}, accusatif masculin.
\end{enumerate}
Soit 7 erreurs localisées dans le texte ; les 3 erreurs restantes à accepter selon la copie de l'élève : \all{dein Brief} (comptée en 1), l'absence d'inversion (2), et on peut compter séparément \all{mein Onkel} (3), \all{frauen} (4), \all{menschen} (5), l'ordre \all{soll ... bleiben} (6 et 7 : deux fautes d'ordre des mots), \all{seine Beruf} (8), plus \all{Letzte Woche ich habe} (9) et la forme \all{helfen} + datif implicite (10 : \all{sie hilft gerne Menschen} — datif pluriel correct, mais l'accusatif serait fautif ; on admet ici la formulation). \textbf{Corrections essentielles à retenir :} accusatif possessif \all{deinen/meinen/seinen}, inversion après complément initial, majuscules des noms, verbe modal en fin de subordonnée.
\end{corrige}

\newpage

\coursec{Fiche bilan}

\begin{popfigure}
\centering
\begin{tikzpicture}[
    mindmap,
    grow cyclic,
    every node/.style={concept, execute at begin node=\hskip0pt},
    concept color=popblue,
    level 1/.append style={level distance=4.5cm, sibling angle=360/6},
    level 2/.append style={level distance=3cm, sibling angle=360/6},
    text=white,
    font=\small\bfseries,
    node font=\small\bfseries,
    align=center
]
\node[concept, scale=1.2, align=center]{AL05\\Femme \& Société}
    child[concept color=poporange] { node[concept, align=center]{Lexique\\Famille\\Travail\\Égalité} }
    child[concept color=popgreen] { node[concept, align=center]{Grammaire 1\\Ordre des mots\\Cas (N/A/D/G)} }
    child[concept color=poppurple] { node[concept, align=center]{Grammaire 2\\Temps du passé\\Präteritum/Perfekt} }
    child[concept color=poppink] { node[concept, align=center]{Structures\\Lettre\\Essai\\Dialogue} }
    child[concept color=popgold] { node[concept, align=center]{Traduction\\Thème FR$\to$DE\\Version DE$\to$FR} }
    child[concept color=popdark] { node[concept, align=center]{Rédaction\\Correction\\Relecture} };
\end{tikzpicture}
\legende{Carte mentale de la leçon AL05 : axes lexicaux, grammaticaux et méthodologiques}
\end{popfigure}

\begin{bilan}

\end{bilan}

\courssub{Lexique clé : Famille, Travail, Égalité}
\begin{center}
\begin{tabularx}{\linewidth}{|X|X|X|}
\hline
\rowcolor{popblueL}
\textbf{Deutsch (Artikel, Plural)} & \textbf{Français} & \textbf{Contexte / Collocation} \\
\hline
\all{die Gleichberechtigung, -en} & l'égalité des droits & \all{die Gleichberechtigung von Mann und Frau} \\
\all{der Haushalt, die Haushalte} & le ménage, le foyer & \all{den Haushalt führen} \\
\all{die Erwerbstätigkeit, -en} & l'activité professionnelle & \all{erwerbstätig sein} \\
\all{die Vereinbarkeit, -en} & la conciliation & \all{Vereinbarkeit von Beruf und Familie} \\
\all{der Alleinerziehende, -n} & le/la parent isolé(e) & \all{alleinerziehende Mütter/Väter} \\
\all{die Quote, -n} & le quota & \all{die Frauenquote} \\
\all{diskriminieren} & discriminer & \all{am Arbeitsplatz diskriminiert werden} \\
\all{der Gender Pay Gap, -s} & l'écart salarial & \all{den Gender Pay Gap schließen} \\
\all{die Rollenverteilung, -en} & la répartition des rôles & \all{traditionelle Rollenverteilung} \\
\all{selbstbestimmt} & autonome, libre de ses choix & \all{ein selbstbestimmtes Leben führen} \\
\hline
\end{tabularx}
\end{center}

\courssub{Règles de grammaire à connaître par cœur}
\begin{center}
\begin{tabularx}{\linewidth}{|X|X|}
\hline
\rowcolor{popblueL}
\textbf{Notion} & \textbf{Règle essentielle \& Exemple} \\
\hline
\textbf{Ordre des mots (Hauptsatz)} & Verbe en position 2. Sujet ou complément en position 1. \\
& \all{Heute geht die Frau arbeiten.} (Adv – V – Subj – Inf) \\
\hline
\textbf{Ordre des mots (Nebensatz)} & Conjonction + ... + Verbe à la FIN. \\
& \all{..., weil die Frau arbeiten geht.} \\
\hline
\textbf{Cas : Accusatif} & Objet direct (qui/quoi ?) après verbes transitifs + prép. \cle{durch, für, gegen, ohne, um}. \\
& \all{Sie sucht \textbf{einen Job}.} \all{Das ist \textbf{für die Frau} wichtig.} \\
\hline
\textbf{Cas : Datif} & Objet indirect (à qui ?) + verbes datifs (\cle{helfen, danken, gefallen}) + prép. \cle{aus, bei, mit, nach, seit, von, zu}. \\
& \all{Ich helfe \textbf{der Mutter}.} \all{Mit \textbf{dem Kind} spazieren.} \\
\hline
\textbf{Präteritum (écrit/récit)} & Verbes faibles : \all{-te} (\all{arbeitete}). Verbes forts : voyelle modifiée (\all{ging, war, hatte}). \\
& \all{Früher blieb die Frau zu Hause.} \\
\hline
\textbf{Perfekt (oral/résultat)} & \all{haben/sein} (P2) + Participe II (\all{ge...t/en}). Verbes mouvement/état $\to$ \all{sein}. \\
& \all{Sie hat gearbeitet.} \all{Sie ist nach Hause gegangen.} \\
\hline
\textbf{Passif (Procédé)} & \all{werden} + Participe II. Agent : \cle{von} + Datif. \\
& \all{Frauen werden oft benachteiligt.} \\
\hline
\end{tabularx}
\end{center}

\courssub{Méthodes express}

\begin{methode}[Rédiger une lettre formelle (Beschwerde / Bewerbung)]
\begin{enumerate}
\item \textbf{En-tête :} Expéditeur (haut gauche), Destinataire (sous expéditeur), Lieu + Date (ligne vide, aligné droite).
\item \textbf{Objet (Betreff) :} Gras, sans « Betreff: », résume le motif.
\item \textbf{Formule d'appel :} \all{Sehr geehrte Damen und Herren,} (inconnu) ou \all{Sehr geehrte Frau Müller,}.
\item \textbf{Corps :} Paragraphes distincts (contexte $\to$ arguments $\to$ demande). Style impersonnel / passif si plainte.
\item \textbf{Formule de politesse :} \all{Mit freundlichen Grüßen}.
\item \textbf{Signature + Pièces jointes :} \all{Anlagen}.
\end{enumerate}
\end{methode}

\begin{methode}[Rédiger un court essai (Erörterung)]
\begin{enumerate}
\item \textbf{Introduction :} Accroche (statistique, citation, actualité) $\to$ Problématique (\all{Wie steht es um...?}) $\to$ Annonce du plan.
\item \textbf{Développement (3 arguments) :} Structure \cle{These – Begründung – Beispiel}. Connecteurs : \all{Zunächst, Darüber hinaus, Einerseits – Andererseits, Schließlich}.
\item \textbf{Conclusion :} Résumé + Ouverture / Opinion personnelle (\all{Meiner Meinung nach...}).
\end{enumerate}
\end{methode}

\begin{methode}[Traduction Thème (FR $\to$ DE) \& Version (DE $\to$ FR)]
\begin{enumerate}
\item \textbf{Analyse :} Identifier temps, mode, cas, structure (phrase simple/composée), registres.
\item \textbf{Transposition grammaticale :} FR \all{« en train de »} $\to$ DE \all{gerade / Perfekt/Präsens}. FR Passif $\to$ DE \all{werden}-Passiv ou \all{man} + actif.
\item \textbf{Lexique :} Choisir le mot précis (faux amis : \all{aktuell $\neq$ actuellement}, \all{konsequent $\neq$ conséquent}). Respecter genre/pluriel noms.
\item \textbf{Relecture :} Majuscules noms, ordre verbes (P2 / Fin), terminaisons cas, Umlauts/ß.
\end{enumerate}
\end{methode}



\begin{aretenir}[Checklist avant le Bac]
$\square$ Je maîtrise le vocabulaire thématique (famille, travail, égalité) avec genres et pluriels. \\
$\square$ Je place le verbe en position 2 (Hauptsatz) et en position finale (Nebensatz) sans hésiter. \\
$\square$ Je décline correctement articles, adjectifs et pronoms aux 4 cas (N, A, D, G). \\
$\square$ Je choisis le bon temps du passé : Präteritum (récit écrit) vs Perfekt (oral/résultat). \\
$\square$ Je conjugue sans erreur les verbes forts/irréguliers clés (\all{sein, haben, werden, gehen, kommen, geben}). \\
$\square$ Je structure une lettre formelle (en-tête, objet, formules, paragraphes, signature). \\
$\square$ Je structure un essai dialectique (Intro probléma, Thèses/Arguments/Exemples, Conclusion). \\
$\square$ Je gère les verbes à particule séparable en phrase principale (\all{anrufen $\to$ ruft an}) et subordonnée. \\
$\square$ J'évite les faux amis majeurs (\all{bekommen/erhalten, aktuell/aktuell, Gymnasium/Lycée}). \\
$\square$ Je relis systématiquement : Majuscules (Noms), Umlauts (ä/ö/ü), ß vs ss, Ponctuation. \\
$\square$ En traduction, je rends le sens idiomatique, pas le mot à mot (ex. \all{« il faut que » $\to$ man muss / es ist nötig, dass}). \\
$\square$ Je sais exprimer mon opinion (\all{Meiner Ansicht nach, Ich bin der Meinung, dass... + Nebensatz}).
\end{aretenir}

\findecours

\end{document}
